%==============================================================================
%  Competitive Programming Notebook — Guidebook
%  Autor: Juan  ·  Universidad del Bosque
%  Formato: doble columna (izquierda / derecha), estilo team notebook ICPC
%  Compilar con: pdflatex guidebook.tex   (dos veces, para el índice)
%==============================================================================
\documentclass[10pt,a4paper,twocolumn]{article}

% ---- Codificación e idioma -------------------------------------------------
\usepackage[T1]{fontenc}
\usepackage[utf8]{inputenc}
\usepackage{lmodern}
\usepackage{microtype}
\usepackage{float}
\usepackage{enumitem}
\usepackage{amsmath}   % \text en modo matematico
\usepackage{amssymb}
% Nombres en español (sin depender de babel-spanish, que no está instalado aquí)
\renewcommand{\contentsname}{Índice}
\renewcommand{\partname}{Parte}

% ---- Diseño de página ------------------------------------------------------
\usepackage[a4paper,margin=1.6cm,columnsep=0.8cm]{geometry}
\setlength{\columnseprule}{0.4pt}             % línea fina que divide izquierda/derecha
\usepackage{xcolor}
\definecolor{acento}{RGB}{20,45,85}           % azul marino formal
\definecolor{acentoClaro}{RGB}{70,110,170}
\definecolor{grisCodigo}{RGB}{245,245,247}
\definecolor{comentario}{RGB}{95,120,95}
\definecolor{palabraClave}{RGB}{30,70,150}
\definecolor{cadena}{RGB}{150,80,40}
\def\columnseprulecolor{\color{gray!40}}

% ---- Encabezados y pie -----------------------------------------------------
\usepackage{fancyhdr}
\usepackage{titlesec}
\pagestyle{fancy}
\fancyhf{}
\fancyhead[L]{\small\itshape Competitive Programming Notebook}
\fancyhead[R]{\small\nouppercase{\itshape\rightmark}}
\fancyfoot[C]{\small\thepage}
\renewcommand{\headrulewidth}{0.4pt}

\titleformat{\section}
  {\normalfont\large\bfseries\color{acento}}{\thesection}{0.6em}{}
\titleformat{\subsection}
  {\normalfont\normalsize\bfseries\color{acentoClaro}}{\thesubsection}{0.5em}{}
\titleformat{\part}
  {\normalfont\Large\bfseries\color{acento}\filcenter}{\partname~\thepart}{0.5em}{}

% ---- Resaltado de código ---------------------------------------------------
\usepackage{listings}
\lstset{
  backgroundcolor=\color{grisCodigo},
  basicstyle=\ttfamily\scriptsize,
  keywordstyle=\color{palabraClave}\bfseries,
  commentstyle=\color{comentario}\itshape,
  stringstyle=\color{cadena},
  numbers=left, numberstyle=\tiny\color{gray}, numbersep=5pt,
  showstringspaces=false, breaklines=true, tabsize=2,
  frame=single, rulecolor=\color{gray!40},
  columns=fullflexible, keepspaces=true,
  extendedchars=true,
  literate=
    {á}{{\'a}}1 {é}{{\'e}}1 {í}{{\'i}}1 {ó}{{\'o}}1 {ú}{{\'u}}1
    {Á}{{\'A}}1 {É}{{\'E}}1 {Í}{{\'I}}1 {Ó}{{\'O}}1 {Ú}{{\'U}}1
    {ñ}{{\~n}}1 {Ñ}{{\~N}}1 {¿}{{?`}}1 {¡}{{!`}}1 {°}{{\textdegree}}1
}
\lstdefinestyle{cpp}{language=C++}
\lstdefinestyle{java}{language=Java}
\lstdefinestyle{py}{language=Python}

% ---- Enlaces (índice navegable) --------------------------------------------
\usepackage[colorlinks=true,linkcolor=acento,urlcolor=acentoClaro]{hyperref}

\setcounter{tocdepth}{2}
\setcounter{secnumdepth}{2}

% Marcador de secciones aún sin contenido
\newcommand{\pendiente}{\textit{\color{gray}Contenido en desarrollo.}\par\medskip}

%==============================================================================
\begin{document}

% ===== PORTADA =============================================================
\onecolumn
\thispagestyle{empty}
\begin{titlepage}
  \centering
  \vspace*{\fill}

  {\large\scshape Universidad del Bosque}\\[4pt]
  {\normalsize Facultad de Ingeniería}\\[2pt]
  {\normalsize Ingeniería de Software \& Ingeniería en Robótica}\\[24pt]

  \color{acento}\rule{\linewidth}{1.2pt}\\[8pt]
  {\Huge\bfseries Competitive Programming}\\[10pt]
  {\Huge\bfseries Notebook}\\[10pt]
  {\large\itshape Algoritmos, estructuras de datos y plantillas}\\[8pt]
  {\large C\texttt{++} \quad}\\[6pt]
  \color{acento}\rule{\linewidth}{1.2pt}\\[28pt]

  \color{black}
  {\large\bfseries Juan Carlos Morales}\\[4pt]
  {\normalsize Doble titulación: Sistemas \& Robótica}\\[36pt]

  {\normalsize Bogotá, Colombia \quad$\cdot$\quad 2026}

  \vspace*{\fill}
\end{titlepage}

% ===== ÍNDICE ==============================================================
\thispagestyle{empty}
\tableofcontents
\clearpage

% ===== CUERPO (doble columna) ==============================================
\twocolumn

\part{Introduccion}
%---------------------------------------------------------------------------
\section{Fundamentos Básicos}
%---------------------------------------------------------------------------

\subsection{Entrada / Salida Rápida (Fast I/O)}
Usar lectura rápida para evitar TLE, hay 3 niveles de lectura rápida, estandarizar Tier 2 y usar Tier 3 (lector por
bytes) solo cuando haya del orden de $10^6$--$10^7$ lecturas.

\begin{table}[h]
\centering
\caption{Casos donde se recomienda Fast I/O}
\begin{tabular}{|p{0.36\columnwidth}|p{0.50\columnwidth}|}
\hline
\textbf{Tipo de problema} & \textbf{Ejemplo} \\
\hline
Arrays gigantes & Leer/ordenar/sumar $10^6$--$10^7$ enteros \\
\hline
Muchos test cases (T grande) & $T \leq 10^5$, cada uno trivial pero se acumula \\
\hline
Grafos grandes & $M \sim 10^6$--$2*10^6$ aristas para BFS/DFS/Dijkstra/MST \\
\hline
Geometria con muchas consultas & $L,S$ grandes + muchos test cases seguidos hasta EOF \\
\hline
Muchos tokens de texto & Strings/lineas que suman millones de caracteres \\
\hline
Judges clasicos "trampa de I/O" & SPOJ INTEST, ARRAYSUB, etc. \\
\hline
\end{tabular}
\end{table}

Regla practica: si N/M es del orden de $10^4-10^5$, cin con
\texttt{sync\_with\_stdio(false);} cin.tie(nullptr); alcanza de sobra. El FastReader
es la "artilleria pesada" para cuando eso deja de ser suficiente.

\textbf{Tier 1 — Básico:} \texttt{cin} / \texttt{cout}.
\begin{lstlisting}[style=cpp]
int entero; double decimal; string palabra, linea;
cin >> entero >> decimal >> palabra;  // >> lee UN token
// ! cin >> deja el '\n' pendiente;
// ! el 1er getline sale VACÍO -> hay que ignorarlo
cin.ignore();
getline(cin, linea);                  // línea completa
\end{lstlisting}
\textbf{Tier 2 — Rápido:} desactiva la sync con C (una vez en
\texttt{main}). Punto dulce para casi todo.
\begin{lstlisting}[style=cpp]
ios_base::sync_with_stdio(false);
cin.tie(nullptr);            // ! no mezclar con scanf/printf
cout << entero << "\n";      // "\n", nunca endl (endl frena)
\end{lstlisting}
\textbf{Tier 3 — Más rápido:} lector por bytes con
\texttt{fread}. Solo con $10^6$–$10^7$ lecturas.
\begin{lstlisting}[style=cpp]
struct FastReader {
    static const int TAM = 1 << 16;   // 64 KB
    char bufer[TAM];
    int puntero = 0, leidos = 0;

    int leer() {
        if (puntero == leidos) {
            leidos = (int) fread(bufer, 1, TAM, stdin);
            puntero = 0;
        }
        return leidos == 0 ? -1 : bufer[puntero++]; // -1=EOF
    }
    int nextInt() {
        int resultado = 0, b = leer();
        while (b <= ' ') b = leer();      // saltar blancos
        bool negativo = (b == '-');
        if (negativo) b = leer();
        while (b >= '0' && b <= '9') {
            resultado = resultado * 10 + (b - '0'); b = leer();
        }
        return negativo ? -resultado : resultado;
    }
    long long nextLong() {
        long long resultado = 0; int b = leer();
        while (b <= ' ') b = leer();
        bool negativo = (b == '-');
        if (negativo) b = leer();
        while (b >= '0' && b <= '9') {
            resultado = resultado * 10 + (b - '0'); b = leer();
        }
        return negativo ? -resultado : resultado;
    }
    double nextDouble() {
        double resultado = 0, divisor = 1; int b = leer();
        while (b <= ' ') b = leer();
        bool negativo = (b == '-');
        if (negativo) b = leer();
        while (b >= '0' && b <= '9') {
            resultado = resultado * 10 + (b - '0'); b = leer();
        }
        if (b == '.') {
            b = leer();
            while (b >= '0' && b <= '9') {
                resultado += (b - '0') / (divisor *= 10); b = leer();
            }
        }
        return negativo ? -resultado : resultado;
    }
    string next() {                       // un token
        string s; int b = leer();
        while (b <= ' ') b = leer();
        while (b > ' ') { s += (char) b; b = leer(); }
        return s;
    }
    string nextLine() {                   // línea completa
        string s; int b = leer();
        while (b != '\n' && b != -1) {
            if (b != '\r') s += (char) b; b = leer();
        }
        return s;
    }
};
\end{lstlisting}

\subsection{Formateo de Salida}


La función \texttt{printf} utiliza una cadena de formato para indicar cómo
debe interpretarse y mostrarse cada argumento. Cada especificador comienza
con \texttt{\%} y corresponde al argumento que aparece después de la cadena.

\begin{lstlisting}[style=cpp]
// %d  -> entero (int)
printf("%d\n", 42);
// Entrada: 42
// Salida: 42

// %lld -> entero largo (long long)
long long x = 1234567890123LL;
printf("%lld\n", x);
// Entrada: 1234567890123
// Salida: 1234567890123

// %5d -> ancho minimo de 5 caracteres, alineado a la derecha
printf("%5d\n", 42);
// Entrada: 42
// Salida: "   42"

// %-5d -> ancho minimo de 5 caracteres, alineado a la izquierda
printf("%-5d|\n", 42);
// Entrada: 42
// Salida: "42   |"

// %05d -> ancho 5, completando con ceros
printf("%05d\n", 42);
// Entrada: 42
// Salida: "00042"

// %+d -> muestra siempre el signo
printf("%+d\n", 42);
// Entrada: 42
// Salida: "+42"

// %x -> hexadecimal usando letras minusculas
printf("%x\n", 255);
// Entrada: 255
// Salida: ff

// %X -> hexadecimal usando letras mayusculas
printf("%X\n", 255);
// Entrada: 255
// Salida: FF

// %.2f -> numero real con exactamente 2 decimales
printf("%.2f\n", 3.14159);
// Entrada: 3.14159
// Salida: 3.14

// %.0f -> numero real sin decimales (redondea)
printf("%.0f\n", 3.7);
// Entrada: 3.7
// Salida: 4

// %e -> notacion cientifica
printf("%e\n", 31415.9);
// Entrada: 31415.9
// Salida: 3.141590e+04

// %c -> un solo caracter
printf("%c\n", 'A');
// Entrada: 'A'
// Salida: A

// %s -> cadena de caracteres
printf("%s\n", "hola");
// Entrada: "hola"
// Salida: hola

// Se pueden combinar varios especificadores
printf("%d %lld %.2f\n", 42, 10000000000LL, 3.14159);
// Entrada: 42, 10000000000, 3.14159
// Salida: 42 10000000000 3.14

// %% -> imprime un '%' literal
printf("100%%\n");
// Entrada: no recibe ningun argumento para %
// Salida: 100%

// ! NO uses %n en C/C++ (puede ser peligroso)
// Para scanf, un double se lee con %lf
\end{lstlisting}

\subsubsection*{Estructura de un especificador}

Un formato puede combinar varias partes:

\begin{lstlisting}[style=cpp]
//   % [flags] [width] [.precision] [specifier]
//   %    0       5         .2          f

printf("%05d\n", 42);
//  %  -> comienza el formato
//  0  -> rellena con ceros
//  5  -> ancho minimo de 5 caracteres
//  d  -> interpreta el argumento como int
//
//  Salida: 00042
\end{lstlisting}

Los modificadores más utilizados en programación competitiva son:

\begin{center}
\begin{tabular}{|c|c|c|}
\hline
\textbf{Formato} & \textbf{Tipo} & \textbf{Ejemplo de salida} \\
\hline
\texttt{\%d} & \texttt{int} & \texttt{42} \\
\hline
\texttt{\%lld} & \texttt{long long} & \texttt{1234567890123} \\
\hline
\texttt{\%c} & \texttt{char} & \texttt{A} \\
\hline
\texttt{\%s} & \texttt{string/char*} & \texttt{hola} \\
\hline
\texttt{\%f} & \texttt{double} & \texttt{3.140000} \\
\hline
\texttt{\%.2f} & \texttt{double} & \texttt{3.14} \\
\hline
\texttt{\%x} & \texttt{int} & \texttt{ff} \\
\hline
\texttt{\%X} & \texttt{int} & \texttt{FF} \\
\hline
\end{tabular}
\end{center}

\textbf{Regla práctica:} en programación competitiva, los formatos más
importantes son \texttt{\%d} para \texttt{int}, \texttt{\%lld} para
\texttt{long long}, \texttt{\%c} para caracteres, \texttt{\%s} para cadenas
y \texttt{\%.2f} (o una precisión diferente) para números reales.


\subsection{Condicionales y Ciclos}

\begin{lstlisting}[style=cpp]
// if / else: ejecuta un bloque segun una condicion
if (x > 0) cout << "pos";       // x > 0 -> "pos"
else if (x < 0) cout << "neg";  // x < 0 -> "neg"
else cout << "cero";             // x == 0 -> "cero"

// Operador ternario: condicion ? valor_si_true : valor_si_false
string s = (x >= 0) ? "no neg" : "neg";
// Si x >= 0, s = "no neg"; si no, s = "neg"

// switch: compara una expresion con diferentes valores
// Solo funciona directamente con tipos integrales como int o char,
// NO con string.
switch (x) {
    case 1: /*...*/ break;       // x == 1
    case 2: case 3: /*...*/ break; // x == 2 o x == 3
    default: /*...*/ ;           // ningun case coincide
}
// ! Sin break, se continua ejecutando el siguiente case (fall-through)

// C++17: se puede declarar una variable dentro del if
if (int r = x % 2; r == 0) {}
// r solo existe dentro del if y sus bloques asociados

// for clasico: usa un indice y se repite mientras i < n
for (int i = 0; i < n; i++) {}

// range-based for: recorre directamente los elementos
// No proporciona el indice.
for (int v : datos) {}

// & obtiene una referencia al elemento original.
// Por tanto, las modificaciones afectan al vector.
for (int &v : datos) v *= 2;

// while: ejecuta mientras la condicion sea verdadera
while (c < n) c++;

// do-while: primero ejecuta el bloque y DESPUES verifica la condicion.
// Por eso se ejecuta AL MENOS una vez.
do {
    c--;
} while (c > 0);

// break -> termina inmediatamente el ciclo o switch
// continue -> salta a la siguiente iteracion del ciclo

// ! C++ NO tiene break con etiqueta.
// Para salir de ciclos anidados se puede usar una bandera,
// encapsular la logica en una funcion y usar return, etc.
\end{lstlisting}




\subsection{Conversiones, Parseos y Casteos}

\begin{lstlisting}[style=cpp]
// string -> numero (parseo)
// Convierte una cadena de texto en el tipo numerico indicado.
int e = stoi("42");
// Entrada: "42"  ->  Salida: 42

long long l = stoll("10000000000");
// Entrada: "10000000000"  ->  Salida: 10000000000

double d = stod("3.14");
// Entrada: "3.14"  ->  Salida: 3.14

// Tambien se puede indicar la base numerica.
int bin = stoi("1010", nullptr, 2);
// "1010" en base 2 -> 10

int hex = stoi("ff", nullptr, 16);
// "ff" en base 16 -> 255


// numero -> string
// Convierte cualquier numero compatible en texto.
string s = to_string(42);
// Entrada: 42  ->  Salida: "42"


// Casteo: fuerza la conversion de un tipo a otro.
// Al pasar de decimal a entero, se TRUNCA hacia 0.
int t = (int) 3.99;
// 3.99 -> 3

int r = static_cast<int>(3.99);
// 3.99 -> 3
// static_cast<int> es la forma recomendada en C++


// ! int / int = division ENTERA
double mal = 7 / 2;
// 7 y 2 son int -> 3
// Luego 3 se convierte a double -> 3.0

double bien = 7 / 2.0;
// 2.0 es double -> division decimal
// 7 / 2.0 -> 3.5


// caracteres
// Los char tienen un valor numerico asociado (ASCII).
int cod = (int) 'A';
// 'A' -> 65

char ch = (char) 66;
// 66 -> 'B'

// Truco muy usado para convertir un digito char a int.
int dig = '7' - '0';
// '7' tiene valor 55 y '0' tiene valor 48
// 55 - 48 -> 7
\end{lstlisting}




%---------------------------------------------------------------------------
\part{Estructuras de Datos}
%---------------------------------------------------------------------------

\section{Basic Structures}

\begin{lstlisting}[style=cpp]
int a[5] = {1,2,3,4,5}; // crear con valores: O(n)

int b[5];               // crear de tamano 5
                        // posiciones: b[0] ... b[4]
                        // ! Valores no inicializados

b[0] = 10;              // INSERT: O(1)
b[1] = 20;              // INSERT: O(1)

int x = b[0];           // READ: O(1)

b[0] = 30;              // UPDATE: O(1)

for (int i = 0; i < 5; i++)
    cout << b[i] << " ";// READ: recorrer -> O(n)

// DELETE: no existe metodo para eliminar
// Se puede sobrescribir o manejar logicamente
b[0] = 0;               // DELETE logico: O(1)

// ! El tamano de un array estatico no cambia.
// Para mas posiciones se necesita otro arreglo.
\end{lstlisting}

\subsection{Matrices (Arreglos 2D)}

Arreglo bidimensional (filas $\times$ columnas). \textbf{Complejidad:}
acceso \texttt{[i][j]} $O(1)$; crear $n\times m$ en $O(nm)$.

\begin{lstlisting}[style=cpp]
vector<vector<int>> m(3, vector<int>(4,0)); // CREATE: O(n*m)
// matriz 3x4, todas las posiciones inicializadas en 0

m[1][2] = 7;              // INSERT: O(1)

int x = m[1][2];           // READ: O(1)

m[1][2] = 10;              // UPDATE: O(1)

for (int i = 0; i < 3; i++)
    for (int j = 0; j < 4; j++)
        cout << m[i][j] << " ";
// READ: recorrer toda la matriz -> O(n*m)

// DELETE: no existe metodo para eliminar una posicion.
// Se puede sobrescribir el valor o manejarlo logicamente.
m[1][2] = 0;               // DELETE logico: O(1)

// ! El tamano de la matriz puede cambiar con vector,
// pero un arreglo 2D estatico no puede cambiar de tamano.
\end{lstlisting}



\subsection{Vectores / Listas dinámicas}

Arreglo que \textbf{crece} solo (ArrayList / \texttt{vector} /
\texttt{list}). \textbf{Complejidad:} acceso por índice $O(1)$; agregar
al final $O(1)$ amortizado; insertar/borrar en medio $O(n)$; buscar
$O(n)$.

\begin{lstlisting}[style=cpp]
vector<int> v;              // CREATE: lista vacia

v.push_back(3);             // INSERT: final -> O(1) amortizado
v.push_back(7);             // INSERT: final -> O(1) amortizado

int x = v[0];               // READ: acceso -> O(1)
int n = v.size();            // READ: tamano -> O(1)

v[0] = 10;                  // UPDATE: modificar -> O(1)

v.insert(v.begin() + 1, 5);  // INSERT: medio -> O(n)

v.erase(v.begin() + 1);      // DELETE: medio -> O(n)

// Buscar un elemento recorriendo el vector
for (int x : v)
    if (x == 10) { /* ... */ }
// SEARCH: recorrer -> O(n)

// DELETE: eliminar el ultimo elemento
v.pop_back();               // DELETE: final -> O(1)

// ! vector mantiene acceso por indice O(1).
// ! Al crecer puede realocar su memoria internamente.
\end{lstlisting}



\subsection{Listas enlazadas (LinkedList)}

Nodos enlazados. \textbf{Complejidad:} insertar/eliminar en los
\textbf{extremos} $O(1)$; acceso por índice $O(n)$; buscar $O(n)$. Útil
como cola o deque.

\begin{lstlisting}[style=cpp]
list<int> l;                // CREATE: lista vacia

l.push_front(1);            // INSERT: inicio -> O(1)
l.push_back(9);             // INSERT: final -> O(1)

int x = l.front();          // READ: primer elemento -> O(1)
int y = l.back();           // READ: ultimo elemento -> O(1)

l.front() = 5;              // UPDATE: primer elemento -> O(1)
l.back() = 10;              // UPDATE: ultimo elemento -> O(1)

l.pop_front();              // DELETE: inicio -> O(1)
l.pop_back();               // DELETE: final -> O(1)

// ! No existe acceso directo por indice.
// Avanzar hasta una posicion requiere recorrer nodos.
auto it = l.begin();
advance(it, 2);             // acceso al tercer elemento -> O(n)

// Buscar un elemento requiere recorrer la lista.
for (int x : l)
    if (x == 10) { /* ... */ }
// SEARCH: recorrer -> O(n)

// ! list no permite l[i].
// Para acceso por indice, vector suele ser mejor.
\end{lstlisting}

\subsection{Pilas (Stack)}

Estructura \textbf{LIFO} (último en entrar, primero en salir).
\textbf{Complejidad:} insertar, eliminar y consultar el tope $O(1)$. Para DFS
iterativo, paréntesis balanceados y deshacer.

\begin{lstlisting}[style=cpp]
stack<int> pila;            // CREATE: pila vacia

pila.push(3);               // INSERT: O(1)
pila.push(7);               // INSERT: O(1)

int tope = pila.top();      // READ: consultar tope -> O(1)

pila.top() = 10;            // UPDATE: modificar tope -> O(1)

pila.pop();                 // DELETE: eliminar tope -> O(1)

// ! Solo se puede acceder directamente al elemento superior.
// ! No existe acceso por indice.
\end{lstlisting}


\subsection{Colas (Queue)}

Estructura \textbf{FIFO} (primero en entrar, primero en salir).
\textbf{Complejidad:} encolar y desencolar $O(1)$. Base del BFS.

\begin{lstlisting}[style=cpp]
queue<int> cola;             // CREATE: cola vacia

cola.push(3);                // INSERT: final -> O(1)
cola.push(7);                // INSERT: final -> O(1)

int frente = cola.front();   // READ: primero -> O(1)
int atras = cola.back();     // READ: ultimo -> O(1)

cola.front() = 10;           // UPDATE: modificar primero -> O(1)

cola.pop();                  // DELETE: primero -> O(1)

// ! push agrega al final y pop elimina del frente.
// ! No existe acceso directo por indice.
\end{lstlisting}


\subsection{Colas de prioridad (Priority Queue / Heap)}

Devuelve siempre el elemento de mayor (o menor) prioridad.
\textbf{Complejidad:} insertar y extraer $O(\log n)$; consultar el tope
$O(1)$. Para Dijkstra, Prim, k-ésimo mayor.

\begin{lstlisting}[style=cpp]
priority_queue<int> pq;      // CREATE: MAX-heap por defecto

pq.push(3);                  // INSERT: O(log n)
pq.push(7);                  // INSERT: O(log n)

int mx = pq.top();           // READ: mayor -> O(1)

pq.pop();                    // DELETE: mayor -> O(log n)

// min-heap: el menor queda en el tope
priority_queue<int, vector<int>, greater<int>> pqMin;

pqMin.push(3);               // INSERT: O(log n)
int mn = pqMin.top();        // READ: menor -> O(1)
pqMin.pop();                 // DELETE: O(log n)

// ! No existe acceso directo a elementos intermedios.
// ! priority_queue no tiene una operacion UPDATE directa.
\end{lstlisting}


\subsection{Comparadores en C++ (Comparators)}

\textbf{Idea:} un comparador es una función (lambda, functor, o
\texttt{operator<} sobrecargado) que le dice a \texttt{sort},
\texttt{set}, \texttt{map}, \texttt{priority\_queue}, etc.
\textbf{cómo ordenar} dos elementos. Debe cumplir la propiedad de
\emph{orden estricto débil} (\textit{strict weak ordering}):
devolver \texttt{true} si el primer argumento va
\textbf{estrictamente antes} que el segundo, y \texttt{false} en
cualquier otro caso —incluyendo cuando se consideran ``iguales'' para
efectos de orden. El error más común es usar \texttt{<=} en vez de
\texttt{<}: eso rompe la propiedad y produce comportamiento
indefinido (desde resultados mal ordenados hasta crashes).

\textbf{¿Por qué usarlo?} El orden por defecto casi nunca alcanza:
necesitas ordenar por varios criterios a la vez, en orden
descendente, por una clave derivada (no el valor mismo), o con un
\texttt{struct}/\texttt{pair} propio que no debe usar el orden
lexicográfico por defecto.

\textbf{Casos de uso:}

\begin{itemize}
    \item Ordenar por múltiples campos (ej. por costo, y en empate
    por índice original).
    \item \texttt{priority\_queue} como min-heap en vez de max-heap
    (o con un orden custom sobre un \texttt{struct}).
    \item \texttt{set}/\texttt{map} con orden propio sobre un tipo
    definido por ti.
    \item Ordenar \textbf{índices} según los valores de otro
    arreglo, sin mover el arreglo original.
\end{itemize}

\textbf{Complejidad:} el comparador en sí no cambia la complejidad
del algoritmo que lo usa (\texttt{sort} sigue siendo
$O(n\log n)$); pero un comparador que no sea un orden estricto débil
puede producir UB o bucles infinitos.


\subsubsection{Comparador con lambda para \texttt{sort} (multi-criterio, índices)}

\begin{lstlisting}[style=cpp]
struct Persona {
    string nombre;
    int edad;
};

vector<Persona> personas = {
    {"Ana", 30}, {"Luis", 25}, {"Eva", 25}
};

// ordena por edad ascendente; en empate, por nombre alfabetico
sort(personas.begin(), personas.end(),
     [](const Persona& a, const Persona& b) {
         if (a.edad != b.edad)
             return a.edad < b.edad;
         return a.nombre < b.nombre;
     });
// resultado: Eva(25), Luis(25), Ana(30)

// ordenar en orden DESCENDENTE: solo invierte la comparacion
sort(personas.begin(), personas.end(),
     [](const Persona& a, const Persona& b) {
         return a.edad > b.edad;
     });
// resultado: Ana(30), luego Luis/Eva(25) en su orden relativo previo
\end{lstlisting}

\textbf{Ordenar índices en vez del arreglo (muy común):}

\begin{lstlisting}[style=cpp]
vector<int> valores = {40, 10, 30, 20};

vector<int> idx(valores.size());
iota(idx.begin(), idx.end(), 0);   // idx = [0, 1, 2, 3]

sort(idx.begin(), idx.end(),
     [&](int a, int b) { return valores[a] < valores[b]; });
// idx queda [1, 3, 2, 0]
// (valores[1]=10, valores[3]=20, valores[2]=30, valores[0]=40)

for (int i : idx)
    cout << valores[i] << " ";
// imprime: 10 20 30 40
// sin haber movido "valores" del lugar
\end{lstlisting}


\subsubsection{\texttt{operator<} dentro de un \texttt{struct} (orden por defecto)}

\textbf{Idea:} si sobrecargas \texttt{operator<} dentro de tu propio
tipo, ese es el orden que usan automáticamente \texttt{sort},
\texttt{set}, \texttt{map} (como clave) y \texttt{priority\_queue}
\textbf{sin pasar ningún comparador extra}.

\begin{lstlisting}[style=cpp]
struct Punto {
    int x, y;

    bool operator<(const Punto& otro) const {
        if (x != otro.x) return x < otro.x;
        return y < otro.y;
    }
};

set<Punto> s;
s.insert({3, 4});
s.insert({1, 2});
s.insert({1, 5});

for (auto& p : s)
    cout << "(" << p.x << "," << p.y << ") ";
// imprime en orden: (1,2) (1,5) (3,4)
// -> set los ordeno automaticamente usando operator<
\end{lstlisting}


\subsubsection{\texttt{priority\_queue} con comparador custom}

\textbf{Idea:} \texttt{priority\_queue} es un \textbf{max-heap} por
defecto (el elemento ``mayor'' según el comparador queda arriba). Para
un min-heap con tipos simples, se usa \texttt{greater<T>}. Para un
\texttt{struct} propio con una regla de prioridad distinta a
\texttt{operator<}, se pasa un \textbf{functor} (struct con
\texttt{operator()}).

\begin{lstlisting}[style=cpp]
// min-heap de enteros: el menor queda arriba
priority_queue<int, vector<int>, greater<int>> minHeap;
minHeap.push(5);
minHeap.push(1);
minHeap.push(3);

cout << minHeap.top();
// 1  (el menor, no el mayor)
\end{lstlisting}

\begin{lstlisting}[style=cpp]
struct Persona {
    string nombre;
    int edad;
};

// functor: queda arriba la persona con MENOR edad
// (priority_queue es max-heap, asi que se invierte el sentido:
// operator() debe devolver true si "a" tiene MENOR prioridad que "b")
struct CompararPorEdad {
    bool operator()(const Persona& a, const Persona& b) const {
        return a.edad > b.edad;
    }
};

priority_queue<Persona, vector<Persona>, CompararPorEdad> pq;
pq.push({"Ana", 30});
pq.push({"Luis", 20});
pq.push({"Eva", 25});

cout << pq.top().nombre;
// "Luis"  (la persona mas joven queda en la cima)
\end{lstlisting}


\subsubsection{Reglas de orden estricto débil (\textit{strict weak ordering})}

Todo comparador \texttt{cmp(a, b)} debe cumplir:

\begin{itemize}
    \item \textbf{Irreflexividad:} \texttt{cmp(a, a)} siempre
    \texttt{false}. Un elemento nunca va estrictamente antes que sí
    mismo.
    \item \textbf{Asimetría:} si \texttt{cmp(a, b)} es
    \texttt{true}, entonces \texttt{cmp(b, a)} debe ser
    \texttt{false}.
    \item \textbf{Transitividad:} si \texttt{cmp(a, b)} y
    \texttt{cmp(b, c)} son \texttt{true}, entonces \texttt{cmp(a, c)}
    también debe serlo.
    \item \textbf{Consistencia en los empates:} si ni
    \texttt{cmp(a,b)} ni \texttt{cmp(b,a)} son \texttt{true}, se
    consideran ``equivalentes'' para efectos de orden (aunque no
    sean idénticos como objetos).
\end{itemize}

\textbf{El error más común:}

\begin{lstlisting}[style=cpp]
// INCORRECTO: usa <= en vez de <
sort(v.begin(), v.end(), [](int a, int b) {
    return a <= b;     // viola irreflexividad: cmp(a,a) da true
});
// esto es Undefined Behavior: en el mejor caso ordena mal,
// en el peor caso crashea o hace un bucle infinito.

// CORRECTO:
sort(v.begin(), v.end(), [](int a, int b) {
    return a < b;
});
\end{lstlisting}

\textbf{Otro error común:} comparadores que no son transitivos al
combinar varios criterios sin desempate consistente (ej. comparar
solo por \texttt{a.x < b.x} cuando a veces \texttt{a.x == b.x} pero
querías tratar esos casos con otro criterio). Si dos elementos
``empatan'' en el primer criterio, siempre hay que \textbf{decidir}
con un segundo criterio, o dejarlos expresamente como equivalentes.

\textbf{Regla práctica:} si tu comparador tiene la forma
\texttt{return a.campo1 < b.campo1;} y sospechas que puede haber
empates que te importan, agrega el siguiente criterio explícitamente
(\texttt{if (a.campo1 != b.campo1) return ...; return
a.campo2 < b.campo2;}) en vez de dejar el desempate al azar del
algoritmo de ordenamiento (que no es estable a menos que uses
\texttt{stable\_sort}).


\subsection{Vector circular (Circular Buffer)}

\textbf{Idea:} un arreglo de tamaño fijo (\texttt{cap}) que se comporta
como una \textbf{cola doble}. En lugar de mover elementos cuando se
inserta o elimina, se reutilizan las posiciones del arreglo de forma
circular.

Se mantienen tres variables principales:

\begin{itemize}
    \item \texttt{cap}: capacidad total del arreglo.
    \item \texttt{ini}: posición física donde comienza el primer elemento.
    \item \texttt{tam}: cantidad actual de elementos almacenados.
\end{itemize}

La posición lógica $i$ no coincide necesariamente con la posición física
del arreglo. Se obtiene mediante:

\[
\texttt{posicionFisica(i) = (ini + i) \% cap}
\]

Por ejemplo, con $\texttt{cap}=5$ y $\texttt{ini}=3$:

\[
\begin{array}{c|ccccc}
\text{Índice lógico} & 0 & 1 & 2 & 3 & 4 \\ \hline
\text{Índice físico} & 3 & 4 & 0 & 1 & 2
\end{array}
\]

Cuando se llega al final del arreglo, el módulo hace que la siguiente
posición vuelva al comienzo. Por eso \textbf{no es necesario mover los
elementos}.

\textbf{¿Por qué usarlo?} En un arreglo normal, insertar al frente
requiere desplazar todos los elementos y cuesta $O(n)$. En un vector
circular solo se mueve \texttt{ini}, por lo que cuesta $O(1)$.

\textbf{Casos de uso:}
\begin{itemize}
    \item Colas de tamaño fijo.
    \item Buffers de audio, video o red.
    \item Ventanas deslizantes.
    \item Sistemas donde llegan y salen datos continuamente.
    \item Simulaciones de procesos en una cola.
    \item Mantener únicamente los últimos $k$ elementos.
\end{itemize}

Si el buffer está lleno y se ejecuta \texttt{pushBack}, el nuevo elemento
\textbf{reemplaza al más antiguo}. Esto permite mantener automáticamente
una ventana de tamaño fijo.

\textbf{Complejidad:} insertar, eliminar, consultar extremos, acceder
por posición y comprobar el tamaño son $O(1)$.


\begin{lstlisting}[style=cpp]
struct VectorCircular {
    vector<int> buf;
    int cap, ini = 0, tam = 0;

    VectorCircular(int c) : buf(c), cap(c) {}

    // CREATE: crea un buffer con capacidad fija.
    // Los elementos inicialmente no estan ocupados.
    
    bool empty() {
        return tam == 0;
    }                           // O(1)

    bool full() {
        return tam == cap;
    }                           // O(1)

    int size() {
        return tam;
    }                           // O(1)

    int capacity() {
        return cap;
    }                           // O(1)

    // Convierte una posicion logica en fisica.
    int pos(int i) {
        return (ini + i) % cap;
    }                           // O(1)

    // INSERT: agrega al final.
    void pushBack(int x) {
        buf[(ini + tam) % cap] = x;

        if (tam < cap)
            tam++;
        else
            ini = (ini + 1) % cap;
    }                           // O(1)

    // INSERT: agrega al inicio.
    void pushFront(int x) {
        ini = (ini - 1 + cap) % cap;
        buf[ini] = x;

        if (tam < cap)
            tam++;
    }                           // O(1)

    // READ: consulta el primer elemento.
    int front() {
        return buf[ini];
    }                           // O(1)

    // READ: consulta el ultimo elemento.
    int back() {
        return buf[(ini + tam - 1) % cap];
    }                           // O(1)

    // READ: acceso por posicion logica.
    int at(int i) {
        return buf[(ini + i) % cap];
    }                           // O(1)

    // UPDATE: modifica un elemento.
    void set(int i, int x) {
        buf[(ini + i) % cap] = x;
    }                           // O(1)

    // DELETE: elimina el primer elemento.
    int popFront() {
        int v = buf[ini];
        ini = (ini + 1) % cap;
        tam--;
        return v;
    }                           // O(1)

    // DELETE: elimina el ultimo elemento.
    int popBack() {
        tam--;
        return buf[(ini + tam) % cap];
    }                           // O(1)

    // DELETE: elimina todos los elementos.
    void clear() {
        ini = 0;
        tam = 0;
    }                           // O(1)

    // SEARCH: busca un elemento.
    // A diferencia de las operaciones anteriores,
    // hay que recorrer los elementos.
    bool contains(int x) {
        for (int i = 0; i < tam; i++)
            if (at(i) == x)
                return true;

        return false;
    }                           // O(n)

    // READ: recorrer los elementos en orden logico.
    void print() {
        for (int i = 0; i < tam; i++)
            cout << at(i) << " ";
    }                           // O(n)
};
\end{lstlisting}

\textbf{Ejemplo de funcionamiento:}

\begin{lstlisting}[style=cpp]
VectorCircular v(5);

v.pushBack(10);
v.pushBack(20);
v.pushBack(30);
// [10, 20, 30, _, _]
// ini = 0, tam = 3

v.pushFront(5);
// [5, 10, 20, 30, _]
// ini cambia de 0 a 4

cout << v.front();
// 5

cout << v.back();
// 30

v.popFront();
// elimina 5
// [_, 10, 20, 30, _]

v.pushBack(40);
v.pushBack(50);
// [_, 10, 20, 30, 40, 50]
// El almacenamiento puede haber dado la vuelta.

v.pushBack(60);
// El buffer estaba lleno.
// Se agrega 60 y se elimina el mas antiguo.
// Resultado logico: [20, 30, 40, 50, 60]
\end{lstlisting}

\textbf{Ventaja principal:} el arreglo físico puede verse ``desordenado'',
pero los elementos siempre tienen un orden lógico determinado por
\texttt{ini}. Por ejemplo:

\begin{lstlisting}[style=cpp]
// Memoria fisica:
// [40, 50, 20, 30,  _]
//       ^
//      ini = 2

// Orden logico:
// 20 -> 30 -> 40 -> 50
\end{lstlisting}

No importa dónde estén físicamente los elementos: \texttt{at(0)} siempre
representa el primero, \texttt{at(1)} el segundo, y así sucesivamente.

\textbf{Diferencia con un vector normal:}

\begin{center}
\small
\begin{tabular}{|l|c|c|}
\hline
\textbf{Operación} & \texttt{vector} & \textbf{Circular} \\
\hline
Agregar al final & $O(1)$* & $O(1)$ \\
\hline
Eliminar al final & $O(1)$ & $O(1)$ \\
\hline
Agregar al frente & $O(n)$ & $O(1)$ \\
\hline
Eliminar al frente & $O(n)$ & $O(1)$ \\
\hline
Acceso por índice & $O(1)$ & $O(1)$ \\
\hline
Buscar & $O(n)$ & $O(n)$ \\
\hline
\end{tabular}
\end{center}

\noindent
*Amortizado.

\textbf{Regla práctica:} si necesitas una estructura de tamaño fijo
donde los elementos entran y salen continuamente por los extremos,
un vector circular evita los desplazamientos de un arreglo normal y
mantiene las operaciones principales en $O(1)$.

\section{Estructuras Hash/Tree}


\subsection{Conjuntos hash (HashSet)}

Colección de elementos \textbf{sin duplicados} y sin orden garantizado.
En C\texttt{++} se implementa mediante \texttt{unordered\_set}. Utiliza
una tabla hash, por lo que las operaciones son $O(1)$ en promedio.
Es útil cuando solo interesa saber si un elemento existe, sin necesidad
de mantener los elementos ordenados.

\textbf{Complejidad:} insertar, buscar y borrar $O(1)$ promedio
($O(n)$ peor caso).

\begin{lstlisting}[style=cpp]
unordered_set<int> s;          // CREATE: conjunto vacio

s.insert(3);                   // INSERT: O(1) prom.
s.insert(7);                   // INSERT: O(1) prom.
s.insert(3);                   // no se agrega: ya existe

bool hay = s.count(3);         // SEARCH: O(1) prom.
// hay = 1 si existe, 0 si no existe

auto it = s.find(7);            // SEARCH: O(1) prom.
// devuelve iterador al elemento o s.end()

s.erase(3);                    // DELETE: O(1) prom.

int n = s.size();              // READ: cantidad -> O(1)
bool vacio = s.empty();        // READ: ¿esta vacio? -> O(1)

// Recorrer todos los elementos
for (int x : s)
    cout << x << " ";
// O(n), pero el orden NO esta garantizado.

// ! No existe acceso por indice: s[0] no es valido.
// ! Los elementos repetidos se ignoran.
// ! El orden puede cambiar entre ejecuciones.
\end{lstlisting}

\textbf{Cuándo usarlo:}
\begin{itemize}
    \item Saber rápidamente si un elemento ya apareció.
    \item Eliminar duplicados.
    \item Mantener un conjunto de elementos visitados.
    \item Detectar elementos repetidos en $O(n)$ promedio.
    \item BFS/DFS: guardar vértices visitados cuando sea conveniente.
\end{itemize}


\subsection{Mapas / Diccionarios hash}

Estructura que almacena pares \textbf{clave $\rightarrow$ valor}.
En C\texttt{++} se implementa mediante \texttt{unordered\_map}.
Cada clave identifica un valor y \textbf{no puede repetirse}. Si se
asigna un nuevo valor a una clave existente, se reemplaza el anterior.

\textbf{Complejidad:} insertar, acceder, buscar y borrar por clave
$O(1)$ promedio ($O(n)$ peor caso).

\begin{lstlisting}[style=cpp]
unordered_map<string,int> m;   // CREATE: mapa vacio

m["a"] = 1;                    // INSERT: O(1) prom.
m["b"] = 2;                    // INSERT: O(1) prom.

int v = m["a"];                // READ: acceder -> O(1) prom.
// v = 1

m["a"] = 10;                   // UPDATE: O(1) prom.
// la clave "a" ahora tiene valor 10

m.erase("a");                  // DELETE: O(1) prom.

bool existe = m.count("b");    // SEARCH: O(1) prom.
// existe = 1

auto it = m.find("b");         // SEARCH: O(1) prom.
if (it != m.end())
    cout << it->second;        // valor asociado a "b"

// Recorrer todas las parejas
for (auto [clave, valor] : m)
    cout << clave << " " << valor;
// O(n), sin orden garantizado.
\end{lstlisting}

\textbf{Importante:} usar \texttt{m[clave]} para consultar una clave
que no existe puede \textbf{crear esa clave} con su valor por defecto.
Para comprobar existencia sin insertar, es preferible usar
\texttt{count()} o \texttt{find()}.

\begin{lstlisting}[style=cpp]
unordered_map<string,int> m;

int x = m["a"];
// Si "a" no existia, se crea:
// "a" -> 0
// ! operator[] puede INSERTAR una clave al consultar.

// count(): comprueba si existe SIN insertar
bool existe = m.count("b");
// Devuelve 0 si no existe, 1 si existe.

// find(): busca la clave SIN insertar
auto it = m.find("b");

if (it != m.end())
    cout << it->second;       // valor asociado a "b"
else
    cout << "No existe";

// ! count() solo indica si existe.
// ! find() permite acceder a la pareja clave -> valor.
\end{lstlisting}


\textbf{Uso clásico: contar frecuencias.}

\begin{lstlisting}[style=cpp]
unordered_map<int,int> freq;

for (int x : datos)
    freq[x]++;                 // frecuencia de cada elemento

cout << freq[5];
// cantidad de veces que aparece 5
\end{lstlisting}

\textbf{Cuándo usarlo:}
\begin{itemize}
    \item Contar frecuencias.
    \item Asociar identificadores con información.
    \item Buscar rápidamente un valor a partir de una clave.
    \item Problemas de conteo y frecuencia.
    \item Memorizar resultados de estados o funciones.
\end{itemize}


\subsection{Conjuntos ordenados (TreeSet)}

Conjunto de elementos \textbf{sin duplicados y ordenados}.
En C\texttt{++} se implementa mediante \texttt{set}, normalmente como
un árbol balanceado.

A diferencia de \texttt{unordered\_set}, los elementos siempre se
mantienen ordenados según el comparador utilizado.

\textbf{Complejidad:} insertar, buscar y borrar $O(\log n)$.
Consultar el menor o mayor elemento mediante \texttt{begin()} o
\texttt{rbegin()} es $O(1)$.

\begin{lstlisting}[style=cpp]
set<int> s;                    // CREATE: conjunto vacio

s.insert(5);                   // INSERT: O(log n)
s.insert(1);                   // INSERT: O(log n)
s.insert(3);                   // INSERT: O(log n)
s.insert(3);                   // no se agrega: ya existe

int menor = *s.begin();        // READ: minimo -> O(1)
int mayor = *s.rbegin();       // READ: maximo -> O(1)

auto it = s.find(3);           // SEARCH: O(log n)
if (it != s.end())
    cout << *it;

s.erase(3);                    // DELETE: O(log n)

int n = s.size();              // READ: cantidad -> O(1)

// Recorrer en orden ascendente
for (int x : s)
    cout << x << " ";
// O(n)
\end{lstlisting}

\textbf{lower\_bound y upper\_bound:}

\begin{lstlisting}[style=cpp]
set<int> s = {1, 3, 5, 7, 9};

auto it1 = s.lower_bound(5);
// primer elemento >= 5
// resultado: 5

auto it2 = s.upper_bound(5);
// primer elemento > 5
// resultado: 7
\end{lstlisting}

Esto permite resolver consultas de \textbf{piso y techo} rápidamente:

\begin{lstlisting}[style=cpp]
// FLOOR: mayor elemento <= x
auto it = s.upper_bound(x);

if (it != s.begin()) {
    --it;
    cout << *it;
}

// CEILING: menor elemento >= x
auto it2 = s.lower_bound(x);

if (it2 != s.end())
    cout << *it2;
\end{lstlisting}

\textbf{Cuándo usarlo:}
\begin{itemize}
    \item Cuando se necesitan elementos únicos y ordenados.
    \item Obtener mínimo o máximo dinámicamente.
    \item Encontrar el siguiente elemento mayor o igual.
    \item Encontrar el anterior elemento menor o igual.
    \item Mantener un conjunto mientras se insertan y eliminan elementos.
\end{itemize}

\textbf{Diferencia clave:} si solo necesitas saber si existe un elemento,
\texttt{unordered\_set} suele ser más rápido. Si necesitas consultar
orden, \texttt{lower\_bound}, \texttt{upper\_bound}, mínimo o máximo,
\texttt{set} es la opción adecuada.


\subsection{Mapas ordenados (TreeMap)}

Mapa \textbf{clave $\rightarrow$ valor} cuyas claves se mantienen
ordenadas. En C\texttt{++} se implementa mediante \texttt{map}, usando
un árbol balanceado.

Cada clave aparece una sola vez. Los valores pueden repetirse.

\textbf{Complejidad:} insertar, buscar, acceder y borrar por clave
$O(\log n)$. \texttt{begin()} permite obtener la menor clave en
$O(1)$.

\begin{lstlisting}[style=cpp]
map<string,int> m;             // CREATE: mapa vacio

m["b"] = 2;                    // INSERT: O(log n)
m["a"] = 1;                    // INSERT: O(log n)
m["c"] = 3;                    // INSERT: O(log n)

int v = m["a"];                // READ: O(log n)
// v = 1

m["a"] = 10;                   // UPDATE: O(log n)

m.erase("b");                  // DELETE: O(log n)

auto it = m.find("c");         // SEARCH: O(log n)
if (it != m.end())
    cout << it->second;

// Primera pareja: menor clave
auto primera = m.begin();      // O(1)
cout << primera->first;

// Ultima pareja: mayor clave
auto ultima = prev(m.end());   // O(1) amortizado
cout << ultima->first;
\end{lstlisting}

\textbf{Las claves se recorren ordenadas:}

\begin{lstlisting}[style=cpp]
map<int,string> m;

m[30] = "C";
m[10] = "A";
m[20] = "B";

for (auto [clave, valor] : m)
    cout << clave << " " << valor;

// Salida:
// 10 A
// 20 B
// 30 C
\end{lstlisting}

\textbf{lower\_bound y upper\_bound} también funcionan sobre las claves:

\begin{lstlisting}[style=cpp]
map<int,string> m;
m[10] = "A";
m[20] = "B";
m[30] = "C";

auto it = m.lower_bound(15);
// primera clave >= 15
// apunta a 20

auto it2 = m.upper_bound(20);
// primera clave > 20
// apunta a 30
\end{lstlisting}

\textbf{Uso clásico: contar frecuencias ordenadas.}

\begin{lstlisting}[style=cpp]
map<int,int> freq;

for (int x : datos)
    freq[x]++;

// Las claves aparecen en orden creciente.
for (auto [x, cantidad] : freq)
    cout << x << ": " << cantidad << "\n";
\end{lstlisting}

\textbf{Cuándo usarlo:}
\begin{itemize}
    \item Cuando se necesitan pares clave--valor ordenados.
    \item Contar frecuencias y procesarlas en orden.
    \item Encontrar claves inmediatamente mayores o menores.
    \item Procesar eventos o valores manteniendo un orden dinámico.
    \item Cuando se necesitan operaciones de árbol como
    \texttt{lower\_bound} y \texttt{upper\_bound}.
\end{itemize}


\subsection*{Comparación rápida}

\begin{center}
\small
\begin{tabular}{|l|c|c|c|}
\hline
\textbf{Estructura} & \textbf{Duplicados} & \textbf{Orden} & \textbf{Buscar} \\
\hline
\texttt{unordered\_set} & No & No & $O(1)$ prom. \\
\hline
\texttt{unordered\_map} & Claves: No & No & $O(1)$ prom. \\
\hline
\texttt{set} & No & Sí & $O(\log n)$ \\
\hline
\texttt{map} & Claves: No & Sí & $O(\log n)$ \\
\hline
\end{tabular}
\end{center}

\textbf{Regla práctica:}

\begin{itemize}[nosep,leftmargin=*]
    \item \texttt{unordered\_set} $\rightarrow$ \textit{¿Ya existe este elemento?}
    \item \texttt{unordered\_map} $\rightarrow$ \textit{¿Qué valor corresponde a esta clave?}
    \item \texttt{set} $\rightarrow$ \textit{Elementos únicos + ordenados.}
    \item \texttt{map} $\rightarrow$ \textit{Claves $\rightarrow$ valores + claves ordenadas.}
\end{itemize}


\section{Trees}

\section{Range Queries}

\subsection{Casos de uso}

\begin{table}[H]
\raggedright
\scriptsize
\setlength{\tabcolsep}{2pt}
\renewcommand{\arraystretch}{1.1}
\begin{tabular}{|p{0.43\columnwidth}|p{0.29\columnwidth}|p{0.20\columnwidth}|}
\hline
\textbf{Cuando el problema te pida:} &
\textbf{Entonces usa:} &
\textbf{Complejidad} \\
\hline

Suma de rango sobre un arreglo que \textbf{nunca cambia}, con muchas consultas
(ej. ``suma de ventas entre el día $l$ y $r$'') &
Prefix Sums &
Build $O(n)$ \newline Query $O(1)$ \\
\hline

Suma de rango con actualizaciones \textbf{puntuales} frecuentes
(ej. ``suma el stock de la posición $i$, luego consulta un rango'') &
Fenwick Tree (BIT) &
Build $O(n)$ \newline Query/Update $O(\log n)$ \\
\hline

Sumar $v$ a \textbf{todo un rango} $[l,r]$ y también consultar la
\textbf{suma} de cualquier rango &
FenwickRangeUpdate (doble BIT) &
Query/Update $O(\log n)$ \\
\hline

Range update + range query con operaciones más complejas:
min/max en rango, asignar un valor a todo un rango, contar ceros, etc. &
Segment Tree + Lazy Propagation &
Query/Update $O(\log n)$ \\
\hline

Mínimo, máximo o gcd de cualquier rango sobre un arreglo
\textbf{estático} (ej. RMQ clásico, LCA) &
Sparse Table &
Build $O(n\log n)$ \\ Query $O(1)$ \\
\hline

Mínimo o máximo de rango cuando el arreglo \textbf{sí cambia}
con el tiempo &
Segment Tree &
Query/Update $O(\log n)$ \\
\hline

Consultas 2D: cuántos puntos o cuánta suma hay dentro de un
rectángulo $[x_1,x_2]\times[y_1,y_2]$ &
Fenwick 2D &
Query/Update \newline $O(\log n\log m)$ \\
\hline

Encontrar el $k$-ésimo elemento, o ``menor índice cuya suma acumulada
alcanza $X$'' (frecuencias) &
Fenwick con \texttt{lowerBound} &
$O(\log n)$ \\
\hline

Muchas queries de rango \textbf{offline} (las conoces todas de antemano)
sobre un arreglo estático o con pocos updates &
Mo's Algorithm &
$O((n+q)\sqrt{n})$ \\
\hline

LCA entre dos nodos, o mínimo en rango sobre un recorrido Euler de un
árbol (valores consecutivos difieren en $\pm1$) &
RMQ $\pm1$ \newline (Euler Tour + Sparse Table) &
Build $O(n)$ \newline Query $O(1)$ \\
\hline

Cuántos elementos en el rango $[l,r]$ están dentro de un rango de
\textbf{valores} $[a,b]$, o cuál es el $k$-ésimo menor en $[l,r]$ &
Wavelet Tree &
Query $O(\log n)$ \\
\hline

Queries o updates sobre \textbf{caminos o subárboles} de un árbol
(ej. ``suma el camino de $u$ a $v$'') &
Heavy-Light Decomposition \newline (+ Segment/Fenwick) &
Query/Update \newline $O(\log^2 n)$ \\
\hline

Necesitas algo simple de implementar bajo presión, con balance
razonable entre update y query, sin memorizar mucha teoría &
Sqrt Decomposition &
Query/Update $O(\sqrt{n})$ \\
\hline

\end{tabular}
\caption{Guía rápida: qué estructura de range queries usar según el problema, con su complejidad.}
\end{table}

\subsection{Prefix Sums}

\textbf{Idea:} si un arreglo \textbf{no cambia}, puedes precalcular
\texttt{pref[i]} $=$ suma de los primeros $i$ elementos en $O(n)$, una
sola vez. Cualquier suma de rango $[l,r]$ se responde restando dos
prefijos: \texttt{pref[r+1] - pref[l]}. Es la base conceptual de todo
lo demás en esta sección: Fenwick y Segment Tree existen justamente
para soportar lo que Prefix Sums no puede (actualizaciones), a costa
de más complejidad de código.

\textbf{¿Por qué usarlo?} Es lo más simple y rápido posible para
consultas de rango cuando el arreglo es estático. Cero estructuras de
datos elaboradas, cero recursión, solo un arreglo auxiliar y una
resta.

\textbf{Casos de uso:}

\begin{itemize}
    \item Suma de rango sobre un arreglo que nunca se modifica, con
    muchas queries.
    \item Suma de submatriz (versión 2D) sobre una matriz estática.
    \item Aplicar $B$ actualizaciones de rango \textbf{offline} y
    consultar el arreglo final una sola vez (Difference Array).
    \item Como paso final de un precómputo de teoría de números
    (criba de $\phi$, $\mu$, número de divisores, etc.) cuando el
    problema pide sumas sobre esos valores en un rango.
\end{itemize}

\textbf{Complejidad:} construir $O(n)$ (u $O(nm)$ en 2D), consulta
$O(1)$. No soporta actualizaciones puntuales eficientes: modificar un
elemento y volver a consultar obliga a reconstruir, $O(n)$.


\subsubsection{Prefix Sum 1D}

\begin{lstlisting}[style=cpp]
struct PrefixSum1D {
    vector<long long> pref;
    int n;

    // CREATE: construye los prefijos a partir de a[0..n-1].
    PrefixSum1D(vector<long long>& a) {
        n = a.size();
        pref.assign(n + 1, 0);

        for (int i = 0; i < n; i++)
            pref[i + 1] = pref[i] + a[i];
    }                                          // O(n)

    // READ: suma del rango [l, r] (0-indexado, inclusivo).
    long long rango(int l, int r) {
        if (l > r) return 0;
        return pref[r + 1] - pref[l];
    }                                          // O(1)
};
\end{lstlisting}

\textbf{Ejemplo de funcionamiento:}

\begin{lstlisting}[style=cpp]
vector<long long> a = {5, 2, 9, 1, 7};

PrefixSum1D ps(a);
// arreglo logico: [5, 2, 9, 1, 7]  (posiciones 0..4)

cout << ps.rango(0, 2);
// 5 + 2 + 9 = 16

cout << ps.rango(1, 4);
// 2 + 9 + 1 + 7 = 19

cout << ps.rango(3, 3);
// 1
\end{lstlisting}


\subsubsection{Prefix Sum 2D}

\begin{lstlisting}[style=cpp]
struct PrefixSum2D {
    vector<vector<long long>> pref;
    int n, m;

    // CREATE: construye la matriz de prefijos a partir de
    // a[0..n-1][0..m-1] usando inclusion-exclusion.
    PrefixSum2D(vector<vector<long long>>& a) {
        n = a.size();
        m = a[0].size();
        pref.assign(n + 1, vector<long long>(m + 1, 0));

        for (int i = 0; i < n; i++)
            for (int j = 0; j < m; j++)
                pref[i + 1][j + 1] = a[i][j] + pref[i][j + 1]
                                    + pref[i + 1][j] - pref[i][j];
    }                                          // O(n * m)

    // READ: suma del rectangulo [x1..x2] x [y1..y2] (0-indexado,
    // inclusivo), usando inclusion-exclusion sobre los prefijos.
    long long rango(int x1, int y1, int x2, int y2) {
        return pref[x2 + 1][y2 + 1] - pref[x1][y2 + 1]
             - pref[x2 + 1][y1] + pref[x1][y1];
    }                                          // O(1)
};
\end{lstlisting}

\textbf{Ejemplo de funcionamiento:}

\begin{lstlisting}[style=cpp]
vector<vector<long long>> a = {
    {1, 2, 3},
    {4, 5, 6},
    {7, 8, 9}
};

PrefixSum2D ps(a);

cout << ps.rango(0, 0, 1, 1);
// 1+2+4+5 = 12

cout << ps.rango(1, 1, 2, 2);
// 5+6+8+9 = 28
\end{lstlisting}


\subsubsection{Difference Array (updates de rango offline)}

\textbf{Idea:} el truco inverso a Prefix Sums: para sumar $v$ a todo
un rango $[l,r]$, marcas \texttt{diff[l] += v} y \texttt{diff[r+1] -=
v} en $O(1)$. Después de aplicar \textbf{todas} las actualizaciones,
un solo recorrido con prefijos reconstruye el arreglo final. Sirve
solo cuando todas las updates son \textbf{offline} (las conoces antes
de necesitar consultar); si necesitas consultar \emph{entre}
actualizaciones, usa FenwickRangeUpdate en su lugar.

\begin{lstlisting}[style=cpp]
struct DifferenceArray {
    vector<long long> diff;
    int n;

    // CREATE: reserva un arreglo de diferencias de tamano n.
    DifferenceArray(int n) : n(n) {
        diff.assign(n + 1, 0);
    }

    // UPDATE: marca +v en todo el rango [l, r] (0-indexado).
    void rangoUpdate(int l, int r, long long v) {
        diff[l] += v;
        if (r + 1 <= n)
            diff[r + 1] -= v;
    }                                          // O(1)

    // READ: reconstruye el arreglo final tras aplicar todas
    // las actualizaciones marcadas hasta el momento.
    vector<long long> aplicar() {
        vector<long long> a(n);
        long long acc = 0;

        for (int i = 0; i < n; i++) {
            acc += diff[i];
            a[i] = acc;
        }

        return a;
    }                                          // O(n)
};
\end{lstlisting}

\textbf{Ejemplo de funcionamiento:}

\begin{lstlisting}[style=cpp]
DifferenceArray d(5);

d.rangoUpdate(1, 3, 4);
d.rangoUpdate(0, 4, 1);
d.rangoUpdate(2, 2, 10);

vector<long long> resultado = d.aplicar();
// resultado: [1, 5, 15, 5, 1]
\end{lstlisting}


\subsubsection{Aplicación: suma de Euler Totient ($\phi$) en un rango}

\textbf{Idea:} $\phi(n)$ cuenta cuántos enteros en $[1,n]$ son
coprimos con $n$. Calcular $\phi$ de \emph{un} número es teoría de
números ($O(\sqrt n)$ factorizando); pero cuando el problema pide la
\textbf{suma de $\phi(i)$ para $i$ en un rango $[l,r]$}, con muchas
queries, el patrón se resuelve en dos fases: (1) precalculas el
arreglo de $\phi$ una sola vez con una criba (teoría de números), y
(2) respondes las queries de rango con Prefix Sums sobre ese arreglo
ya calculado (exactamente esta subsección). Factorizar cada número
del rango por separado en cada query sería $O(q\sqrt n)$ y no escala;
precalcular y sumar con prefijos reduce cada query a $O(1)$.

\paragraph{Criba lineal + Prefix Sums (rango $[1,N]$).}

\begin{lstlisting}[style=cpp]
struct PhiRangeSum {
    vector<long long> phi, pref;
    int n;

    // CREATE: calcula phi[1..n] con criba lineal y sus prefijos.
    // phi[i] = cantidad de enteros en [1, i] coprimos con i.
    PhiRangeSum(int n) : n(n) {
        phi.assign(n + 1, 0);
        pref.assign(n + 1, 0);

        vector<int> primos;
        vector<bool> compuesto(n + 1, false);
        phi[1] = 1;

        for (int i = 2; i <= n; i++) {
            if (!compuesto[i]) {
                primos.push_back(i);
                phi[i] = i - 1;
            }

            for (int p : primos) {
                if ((long long)i * p > n) break;

                compuesto[i * p] = true;

                if (i % p == 0) {
                    phi[i * p] = phi[i] * p;
                    break;
                } else {
                    phi[i * p] = phi[i] * (p - 1);
                }
            }
        }

        for (int i = 1; i <= n; i++)
            pref[i] = pref[i - 1] + phi[i];
    }                                          // O(n)

    // READ: suma de phi(i) para i en [l, r].
    long long rango(int l, int r) {
        if (l > r) return 0;
        return pref[r] - pref[l - 1];
    }                                          // O(1)

    // READ: valor individual de phi(i).
    long long get(int i) {
        return phi[i];
    }                                          // O(1)
};
\end{lstlisting}

\textbf{Ejemplo de funcionamiento:}

\begin{lstlisting}[style=cpp]
PhiRangeSum pr(10);
// phi(1..10) = [1, 1, 2, 2, 4, 2, 6, 4, 6, 4]

cout << pr.get(7);
// phi(7) = 6  (7 es primo)

cout << pr.rango(1, 10);
// 1+1+2+2+4+2+6+4+6+4 = 32

cout << pr.rango(4, 8);
// 2+4+2+6+4 = 18
\end{lstlisting}

\paragraph{Criba segmentada (rango $[L,R]$ grande).}

Si $R$ es muy grande ($10^{12}$) pero $R-L$ es manejable, no puedes
cribar desde 1. En vez de eso, precalculas los primos hasta
$\sqrt R$, y para cada número del segmento $[L,R]$ le quitas sus
factores primos pequeños, aplicando $\phi(n) = n\prod_{p\mid
n}(1-\tfrac1p)$. Si al final queda un residuo $>1$, es un único
primo grande (mayor a $\sqrt R$) que también se aplica a la fórmula.

\begin{lstlisting}[style=cpp]
struct PhiSegmentedRangeSum {
    vector<long long> phi, pref;
    long long L, R;

    // CREATE: calcula phi(i) para cada i en [L, R] factorizando con
    // los primos hasta sqrt(R), y sus prefijos sobre el segmento.
    PhiSegmentedRangeSum(long long L, long long R) : L(L), R(R) {
        int lim = (int)sqrt((double)R) + 1;

        vector<int> primos;
        vector<bool> compuesto(lim + 1, false);
        for (int i = 2; i <= lim; i++) {
            if (!compuesto[i]) primos.push_back(i);
            for (int p : primos) {
                if ((long long)i * p > lim) break;
                compuesto[i * p] = true;
                if (i % p == 0) break;
            }
        }                                      // O(sqrt(R) log log sqrt(R))

        int sz = R - L + 1;
        vector<long long> num(sz);
        phi.assign(sz, 0);

        for (int i = 0; i < sz; i++) {
            num[i] = L + i;
            phi[i] = L + i;
        }

        for (int p : primos) {
            long long start = ((L + p - 1) / p) * p;

            for (long long j = start; j <= R; j += p) {
                int idx = j - L;
                phi[idx] -= phi[idx] / p;

                while (num[idx] % p == 0)
                    num[idx] /= p;
            }
        }

        // lo que quede > 1 es un primo grande (> sqrt(R)).
        for (int i = 0; i < sz; i++)
            if (num[i] > 1)
                phi[i] -= phi[i] / num[i];

        pref.assign(sz + 1, 0);
        for (int i = 0; i < sz; i++)
            pref[i + 1] = pref[i] + phi[i];
    }                                          // O((R-L+1) log log R + sqrt(R))

    // READ: suma de phi(i) para i en [l, r], con L <= l <= r <= R.
    long long rango(long long l, long long r) {
        return pref[r - L + 1] - pref[l - L];
    }                                          // O(1)

    // READ: valor individual de phi(i), con L <= i <= R.
    long long get(long long i) {
        return phi[i - L];
    }                                          // O(1)
};
\end{lstlisting}

\textbf{Ejemplo de funcionamiento:}

\begin{lstlisting}[style=cpp]
PhiSegmentedRangeSum ps(95, 105);
// calcula phi(95), phi(96), ..., phi(105)
// sin cribar desde 1

cout << ps.get(100);
// phi(100) = 40

cout << ps.rango(95, 105);
// suma de phi(95) a phi(105)

cout << ps.rango(100, 102);
// phi(100) + phi(101) + phi(102)
\end{lstlisting}

\textbf{Cuál variante usar:}

\begin{table}[H]
\centering
\tiny
\setlength{\tabcolsep}{1.5pt}
\renewcommand{\arraystretch}{1.1}
\begin{tabular}{|p{0.32\columnwidth}|p{0.30\columnwidth}|p{0.27\columnwidth}|}
\hline
\textbf{Situación} & \textbf{Estructura} & \textbf{Complejidad} \\
\hline
$[1,N]$, $N$ chico/mediano
& \texttt{PhiRangeSum}
& $O(N)$ \\
\hline
$[L,R]$ grande, $R-L$ chico
& \texttt{PhiSegmentedRangeSum}
& $O((R{-}L)\log\log R+\sqrt{R})$ \\
\hline
$\phi$ cambia entre queries
& Fenwick sobre $\phi$
& $O(\log N)$ \\
\hline
\end{tabular}
\caption{Variantes para la suma de $\phi$.}
\label{tab:phi_variants}
\end{table}

\textbf{Regla práctica:} el cálculo de $\phi$ en sí (la criba, lineal
o segmentada) es teoría de números y no cambia según la query. Una vez
que tienes el arreglo de $\phi$ listo, trátalo como \emph{cualquier
otro arreglo estático}: Prefix Sums resuelve el resto. No hace falta
una estructura especial para $\phi$ más allá de construir bien el
arreglo inicial.



\subsection{Árbol de Fenwick (BIT)}

\textbf{Idea:} un arreglo normal da suma de prefijo en $O(1)$ pero
actualizar un elemento cuesta $O(n)$ (o al revés si guardas prefijos).
El Fenwick equilibra ambas a $O(\log n)$ guardando \emph{sumas
parciales inteligentes}: la posición $i$ almacena la suma de un bloque
de tamaño $i\,\&\,(-i)$ (su bit menos significativo) que termina en
$i$. Para consultar un prefijo saltas hacia atrás quitando el bit bajo
($i\mathrel{-}=i\,\&\,(-i)$); para actualizar subes sumándolo
($i\mathrel{+}=i\,\&\,(-i)$). Cada camino toca a lo más $\log n$
posiciones (un bit por paso). Es 1-indexado.

\textbf{¿Por qué usarlo?} Es mucho más simple y liviano que un Segment
Tree cuando solo necesitas sumas (o cualquier operación invertible)
de prefijo/rango con actualizaciones puntuales. Menos memoria, menos
código y una constante más pequeña.

\textbf{Casos de uso:}

\begin{itemize}
    \item Suma de prefijo / suma de rango con updates frecuentes.
    \item Conteo de inversiones en un arreglo.
    \item Frecuencias acumuladas y búsqueda del $k$-ésimo elemento.
    \item Range update + range query cuando no hace falta un Segment
    Tree con Lazy Propagation.
    \item Consultas 2D: conteo/suma de puntos dentro de un rectángulo.
\end{itemize}

\textbf{Complejidad:} construir $O(n)$, update puntual $O(\log n)$,
query de prefijo/rango $O(\log n)$, range update + range query
$O(\log n)$ y versión 2D $O(\log n \cdot \log m)$.


\subsubsection{BIT clásico (Point Update, Range Query)}

\begin{lstlisting}[style=cpp]
struct FenwickTree {
    vector<long long> t;
    int n;

    // CREATE: reserva un arbol de tamano n (1-indexado).
    FenwickTree(int n) : t(n + 1, 0), n(n) {}

    // CREATE: construye el arbol a partir de un arreglo a[0..n-1]
    // propagando cada valor a su "padre" una sola vez, en vez de
    // hacer n updates (que serian O(n log n)).
    void build(vector<long long>& a) {
        for (int i = 1; i <= n; i++) {
            t[i] += a[i - 1];

            int j = i + (i & (-i));

            if (j <= n)
                t[j] += t[i];
        }
    }                                      // O(n)

    // UPDATE: suma v al elemento en la posicion i.
    void update(int i, long long v) {
        for (; i <= n; i += i & (-i))
            t[i] += v;
    }                                      // O(log n)

    // UPDATE: fija el valor absoluto de la posicion i en x.
    void set(int i, long long x) {
        update(i, x - get(i));
    }                                      // O(log n)

    // READ: suma de prefijo [1..i].
    long long query(int i) {
        long long s = 0;

        for (; i > 0; i -= i & (-i))
            s += t[i];

        return s;
    }                                      // O(log n)

    // READ: suma del rango [l, r].
    long long rango(int l, int r) {
        if (l > r)
            return 0;

        return query(r) - query(l - 1);
    }                                      // O(log n)

    // READ: valor puntual en la posicion i.
    long long get(int i) {
        return rango(i, i);
    }                                      // O(log n)

    // SEARCH: menor indice i tal que query(i) >= objetivo.
    // Usa binary lifting sobre el arbol.
    // Requiere que todos los valores sean >= 0.
    int lowerBound(long long objetivo) {
        int pos = 0;
        long long ac = 0;

        int LOG = 0;

        while ((1 << LOG) <= n)
            LOG++;

        for (int k = LOG - 1; k >= 0; k--) {
            int nx = pos + (1 << k);

            if (nx <= n && ac + t[nx] < objetivo) {
                pos = nx;
                ac += t[nx];
            }
        }

        return pos + 1;
    }                                      // O(log n)

    // DELETE: reinicia el arbol a ceros.
    void clear() {
        fill(t.begin(), t.end(), 0);
    }                                      // O(n)
};
\end{lstlisting}


\textbf{Ejemplo de funcionamiento:}

\begin{lstlisting}[style=cpp]
vector<long long> a = {5, 2, 9, 1, 7};

FenwickTree bit(5);
bit.build(a);

// arreglo logico:
// [5, 2, 9, 1, 7]
// posiciones 1..5

cout << bit.query(3);
// 5 + 2 + 9 = 16

cout << bit.rango(2, 4);
// 2 + 9 + 1 = 12

bit.update(2, 3);

// arreglo logico:
// [5, 5, 9, 1, 7]

cout << bit.get(2);
// 5

bit.set(4, 10);

// arreglo logico:
// [5, 5, 9, 10, 7]

cout << bit.lowerBound(15);
// menor indice cuyo prefijo alcanza 15
// -> indice 3 (5 + 5 + 9 = 19 >= 15)
\end{lstlisting}


\subsubsection{BIT con Range Update + Range Query (Doble BIT)}

\textbf{Idea:} con \emph{dos} Fenwick y el truco de diferencias se
puede sumar $v$ a todo un rango $[l,r]$ en $O(\log n)$ y también
consultar la suma de cualquier rango en $O(\log n)$. Esto permite
resolver simultáneamente actualizaciones y consultas de rango.

\begin{lstlisting}[style=cpp]
struct FenwickRangeUpdate {
    FenwickTree b1, b2;
    int n;

    // CREATE: reserva dos arboles auxiliares de tamano n.
    FenwickRangeUpdate(int n) : b1(n), b2(n), n(n) {}

    // UPDATE: suma v a todos los elementos del rango [l, r].
    void rangoUpdate(int l, int r, long long v) {
        b1.update(l, v);
        b1.update(r + 1, -v);

        b2.update(l, v * (l - 1));
        b2.update(r + 1, -v * r);
    }                                      // O(log n)

    // READ: suma de prefijo [1..i].
    long long prefijo(int i) {
        return b1.query(i) * i - b2.query(i);
    }                                      // O(log n)

    // READ: suma del rango [l, r].
    long long rango(int l, int r) {
        if (l > r)
            return 0;

        return prefijo(r) - prefijo(l - 1);
    }                                      // O(log n)

    // READ: valor puntual en la posicion i.
    long long get(int i) {
        return rango(i, i);
    }                                      // O(log n)
};
\end{lstlisting}

\textbf{Ejemplo de funcionamiento:}

\begin{lstlisting}[style=cpp]
FenwickRangeUpdate f(5);
// arreglo logico inicial: [0, 0, 0, 0, 0]

f.rangoUpdate(2, 4, 3);
// arreglo logico: [0, 3, 3, 3, 0]

cout << f.rango(1, 5);
// 0+3+3+3+0 = 9

f.rangoUpdate(1, 3, 2);
// arreglo logico: [2, 5, 5, 3, 0]

cout << f.get(2);
// 5
\end{lstlisting}

\subsubsection{BIT 2D (point update, sum de submatriz)}

\textbf{Idea:} un BIT dentro de otro BIT. Cada update/query en $x$
dispara un update/query completo en $y$, dando $O(\log n \cdot
\log m)$ en vez de $O(nm)$ por consulta.

\begin{lstlisting}[style=cpp]
struct Fenwick2D {
    vector<vector<long long>> t;
    int n, m;

    // CREATE: reserva una matriz de arboles de tamano n x m.
    Fenwick2D(int n, int m) : t(n + 1, vector<long long>(m + 1, 0)),
                               n(n), m(m) {}

    // UPDATE: suma v a la celda (x, y).
    void update(int x, int y, long long v) {
        for (int i = x; i <= n; i += i & (-i))
            for (int j = y; j <= m; j += j & (-j))
                t[i][j] += v;
    }                                          // O(log n * log m)

    // READ: suma de la submatriz [1..x] x [1..y].
    long long query(int x, int y) {
        long long s = 0;
        for (int i = x; i > 0; i -= i & (-i))
            for (int j = y; j > 0; j -= j & (-j))
                s += t[i][j];
        return s;
    }                                          // O(log n * log m)

    // READ: suma del rectangulo [x1..x2] x [y1..y2].
    long long rango(int x1, int y1, int x2, int y2) {
        return query(x2, y2) - query(x1 - 1, y2)
             - query(x2, y1 - 1) + query(x1 - 1, y1 - 1);
    }                                          // O(log n * log m)
};
\end{lstlisting}

\textbf{Ejemplo de funcionamiento:}

\begin{lstlisting}[style=cpp]
Fenwick2D f(4, 4);

f.update(1, 1, 5);
f.update(2, 3, 7);
f.update(4, 4, 2);

cout << f.rango(1, 1, 2, 3);
// 5 + 7 = 12

cout << f.rango(3, 3, 4, 4);
// 2
\end{lstlisting}

\textbf{Diferencia con un arreglo de prefijos simple:}

\begin{table}[H]
\raggedright
\scriptsize
\setlength{\tabcolsep}{3pt}
\begin{tabular}{|l|c|c|c|}
\hline
\textbf{Estructura} & \textbf{Construir} & \textbf{Query} & \textbf{Update} \\
\hline
Prefijos simples & $O(n)$ & $O(1)$ & $O(n)$ \\
\hline
\texttt{FenwickTree} & $O(n)$ & $O(\log n)$ & $O(\log n)$ \\
\hline
\texttt{FenwickRangeUpdate} & $O(n)$ & $O(\log n)$ & $O(\log n)$* \\
\hline
\texttt{Fenwick2D} & $O(nm)$ & $O(\log n\log m)$ & $O(\log n\log m)$ \\
\hline
\end{tabular}
\caption{Comparación de complejidades: Fenwick vs. prefijos simples.}
\end{table}

\noindent
*Range update + range query, con la variante de doble BIT.

\textbf{Regla práctica:} si el arreglo cambia con frecuencia y solo
necesitas sumas (o cualquier operación invertible) de prefijo o rango,
el Fenwick da el mejor balance entre simplicidad de código y
velocidad. Si necesitas min/max, operaciones no invertibles, o updates
de rango con lazy propagation general, usa Segment Tree en su lugar.




\subsection{Sqrt Decomposition}

\textbf{Idea:} divide el arreglo en bloques de tamaño
$\approx\sqrt n$ y guarda un resumen (suma, mínimo, etc.) por bloque.
Un update toca solo un bloque; una query de rango toca a lo más dos
bloques \emph{parciales} (recorridos elemento por elemento) y varios
bloques \emph{completos} en el medio (usados de golpe con su resumen
precalculado). Como hay $\approx\sqrt n$ bloques y cada bloque tiene
$\approx\sqrt n$ elementos, todo cuesta $O(\sqrt n)$: ni tan lento
como un arreglo plano, ni tan rápido como Fenwick/Segment Tree, pero
mucho más simple de escribir y de adaptar a operaciones raras.

\textbf{¿Por qué usarlo?} Es la estructura más flexible cuando no
tienes tiempo de pensar un Segment Tree con lazy propagation
correcto: soporta casi cualquier operación (suma, min, max, contar
algo con una condición) con el mismo esqueleto de código, aunque sea
$O(\sqrt n)$ en vez de $O(\log n)$. También es la base conceptual de
Mo's Algorithm.

\textbf{Casos de uso:}

\begin{itemize}
    \item Necesitas implementar algo rápido bajo presión, sin
    memorizar casos de Segment Tree con lazy propagation.
    \item Operaciones donde min/max cambian de forma que Fenwick no
    puede manejar (Fenwick solo sirve para operaciones invertibles
    como la suma).
    \item $n$ no es tan grande ($\le 10^5$–$10^6$) y $O(\sqrt n)$ por
    operación es suficiente para el límite de tiempo.
    \item Base para Mo's Algorithm (ver esa sección).
\end{itemize}

\textbf{Complejidad:} construir $O(n)$, update puntual $O(1)$ (suma)
u $O(\sqrt n)$ (min/max), query de rango $O(\sqrt n)$, range update +
range query $O(\sqrt n)$ con lazy por bloque.


\subsubsection{Sqrt Decomposition clásico (point update, range query — suma)}

\begin{lstlisting}[style=cpp]
struct SqrtDecomposition {
    vector<long long> a, bloque;
    int n, tam;

    // CREATE: construye el arreglo en bloques de tamano ~sqrt(n).
    SqrtDecomposition(vector<long long>& arr) {
        a = arr;
        n = a.size();
        tam = max(1, (int)sqrt((double)n));

        bloque.assign((n + tam - 1) / tam, 0);

        for (int i = 0; i < n; i++)
            bloque[i / tam] += a[i];
    }                                          // O(n)

    // UPDATE: suma v al elemento en la posicion i.
    void update(int i, long long v) {
        a[i] += v;
        bloque[i / tam] += v;
    }                                          // O(1)

    // READ: suma del rango [l, r] (0-indexado, inclusivo).
    long long rango(int l, int r) {
        long long s = 0;
        int bl = l / tam, br = r / tam;

        if (bl == br) {
            for (int i = l; i <= r; i++)
                s += a[i];
            return s;
        }

        for (int i = l; i < (bl + 1) * tam; i++)
            s += a[i];

        for (int b = bl + 1; b < br; b++)
            s += bloque[b];

        for (int i = br * tam; i <= r; i++)
            s += a[i];

        return s;
    }                                          // O(sqrt(n))
};
\end{lstlisting}

\textbf{Ejemplo de funcionamiento:}

\begin{lstlisting}[style=cpp]
vector<long long> a = {5, 2, 9, 1, 7, 3, 8, 4};

SqrtDecomposition sd(a);
// arreglo logico: [5, 2, 9, 1, 7, 3, 8, 4]  (posiciones 0..7)
// tam = 2 -> bloques: [5+2, 9+1, 7+3, 8+4] = [7, 10, 10, 12]

cout << sd.rango(1, 5);
// 2 + 9 + 1 + 7 + 3 = 22

sd.update(3, 10);
// arreglo logico: [5, 2, 9, 11, 7, 3, 8, 4]

cout << sd.rango(1, 5);
// 2 + 9 + 11 + 7 + 3 = 32
\end{lstlisting}


\subsubsection{Sqrt Decomposition con range update + range query (lazy por bloque)}

\textbf{Idea:} igual que arriba, pero cada bloque guarda un valor
\texttt{lazy} pendiente que representa ``sumar $v$ a todo el bloque''
sin tocar cada elemento individualmente. Solo se aplica a los
elementos reales cuando un update o query \textbf{parcial} necesita
tocar ese bloque específico.

\begin{lstlisting}[style=cpp]
struct SqrtRangeUpdate {
    vector<long long> a, bloqueSum, lazy;
    int n, tam;

    // CREATE: construye el arreglo en bloques de tamano ~sqrt(n).
    SqrtRangeUpdate(vector<long long>& arr) {
        a = arr;
        n = a.size();
        tam = max(1, (int)sqrt((double)n));

        int nb = (n + tam - 1) / tam;
        bloqueSum.assign(nb, 0);
        lazy.assign(nb, 0);

        for (int i = 0; i < n; i++)
            bloqueSum[i / tam] += a[i];
    }                                          // O(n)

    // READ auxiliar: cantidad de elementos reales del bloque b.
    int tamBloque(int b) {
        return min(n, (b + 1) * tam) - b * tam;
    }                                          // O(1)

    // UPDATE: suma v a todo el rango [l, r].
    void rangoUpdate(int l, int r, long long v) {
        int bl = l / tam, br = r / tam;

        if (bl == br) {
            for (int i = l; i <= r; i++)
                a[i] += v;
            bloqueSum[bl] += v * (r - l + 1);
            return;
        }

        for (int i = l; i < (bl + 1) * tam; i++)
            a[i] += v;
        bloqueSum[bl] += v * ((bl + 1) * tam - l);

        for (int b = bl + 1; b < br; b++)
            lazy[b] += v;

        for (int i = br * tam; i <= r; i++)
            a[i] += v;
        bloqueSum[br] += v * (r - br * tam + 1);
    }                                          // O(sqrt(n))

    // READ: suma del rango [l, r].
    long long rango(int l, int r) {
        long long s = 0;
        int bl = l / tam, br = r / tam;

        if (bl == br) {
            for (int i = l; i <= r; i++)
                s += a[i] + lazy[bl];
            return s;
        }

        for (int i = l; i < (bl + 1) * tam; i++)
            s += a[i] + lazy[bl];

        for (int b = bl + 1; b < br; b++)
            s += bloqueSum[b] + lazy[b] * tamBloque(b);

        for (int i = br * tam; i <= r; i++)
            s += a[i] + lazy[br];

        return s;
    }                                          // O(sqrt(n))
};
\end{lstlisting}

\textbf{Ejemplo de funcionamiento:}

\begin{lstlisting}[style=cpp]
vector<long long> a = {1, 1, 1, 1, 1, 1};

SqrtRangeUpdate sr(a);
// tam = 2 -> bloques: [1+1, 1+1, 1+1] = [2, 2, 2]

sr.rangoUpdate(1, 4, 10);
// arreglo logico "efectivo": [1, 11, 11, 11, 11, 1]

cout << sr.rango(0, 5);
// 1+11+11+11+11+1 = 46

cout << sr.rango(2, 3);
// 11+11 = 22
\end{lstlisting}


\subsubsection{Variante para mínimo / máximo (point update, range query)}

\textbf{Idea:} a diferencia de Fenwick (que solo sirve para
operaciones \emph{invertibles} como la suma), Sqrt Decomposition
soporta min/max con el mismo esqueleto: la única diferencia es que un
update ya no puede ajustar el resumen del bloque en $O(1)$, porque
quitar el valor viejo no dice si sigue habiendo otro elemento igual
de chico. Hay que \textbf{reconstruir el mínimo de todo el bloque},
lo cual sigue siendo $O(\sqrt n)$.

\begin{lstlisting}[style=cpp]
struct SqrtMinQuery {
    vector<long long> a, bloqueMin;
    int n, tam;

    // CREATE: construye el arreglo en bloques de tamano ~sqrt(n).
    SqrtMinQuery(vector<long long>& arr) {
        a = arr;
        n = a.size();
        tam = max(1, (int)sqrt((double)n));

        bloqueMin.assign((n + tam - 1) / tam, LLONG_MAX);

        for (int i = 0; i < n; i++)
            bloqueMin[i / tam] = min(bloqueMin[i / tam], a[i]);
    }                                          // O(n)

    // UPDATE: fija el valor absoluto de la posicion i en x.
    // Reconstruye el minimo de todo el bloque, ya que a diferencia
    // de la suma, el minimo no se puede ajustar incrementalmente.
    void set(int i, long long x) {
        a[i] = x;

        int b = i / tam;
        long long mn = LLONG_MAX;

        for (int j = b * tam; j < min(n, (b + 1) * tam); j++)
            mn = min(mn, a[j]);

        bloqueMin[b] = mn;
    }                                          // O(sqrt(n))

    // READ: minimo del rango [l, r].
    long long rango(int l, int r) {
        long long mn = LLONG_MAX;
        int bl = l / tam, br = r / tam;

        if (bl == br) {
            for (int i = l; i <= r; i++)
                mn = min(mn, a[i]);
            return mn;
        }

        for (int i = l; i < (bl + 1) * tam; i++)
            mn = min(mn, a[i]);

        for (int b = bl + 1; b < br; b++)
            mn = min(mn, bloqueMin[b]);

        for (int i = br * tam; i <= r; i++)
            mn = min(mn, a[i]);

        return mn;
    }                                          // O(sqrt(n))
};
\end{lstlisting}

\textbf{Ejemplo de funcionamiento:}

\begin{lstlisting}[style=cpp]
vector<long long> a = {5, 2, 9, 1, 7, 3, 8, 4};

SqrtMinQuery sm(a);
// tam = 2 -> bloques min: [2, 1, 3, 4]

cout << sm.rango(0, 4);
// min(5,2,9,1,7) = 1

sm.set(3, 100);
// arreglo logico: [5, 2, 9, 100, 7, 3, 8, 4]

cout << sm.rango(0, 4);
// min(5,2,9,100,7) = 2
\end{lstlisting}

\textbf{Comparación con Fenwick y Segment Tree:}

\begin{table}[H]
\raggedright
\scriptsize
\setlength{\tabcolsep}{3pt}
\begin{tabular}{|l|c|c|c|}
\hline
\textbf{Estructura} & \textbf{Update} & \textbf{Query} & \textbf{Min/Max?} \\
\hline
Sqrt Decomposition & $O(1)$–$O(\sqrt n)$ & $O(\sqrt n)$ & Sí \\
\hline
\texttt{FenwickTree} & $O(\log n)$ & $O(\log n)$ & No \\
\hline
Segment Tree & $O(\log n)$ & $O(\log n)$ & Sí \\
\hline
\end{tabular}
\caption{Sqrt Decomposition vs. Fenwick vs. Segment Tree.}
\end{table}

\textbf{Regla práctica:} usa Sqrt Decomposition cuando quieras algo
simple y flexible de escribir rápido, sobre todo si la operación no
es una suma (Fenwick no sirve) y no quieres arriesgarte con un
Segment Tree + Lazy Propagation bajo presión de tiempo. Si $n$ es muy
grande o el límite de tiempo es ajustado, $O(\log n)$ de Segment Tree
gana; si la operación es solo suma, Fenwick es más simple y más
rápido.


\subsection{Sparse Table (RMQ)}

\textbf{Idea:} responde consultas de rango en $O(1)$ sobre un arreglo
que \textbf{no cambia}, precalculando \texttt{sp[k][i]} $=$ el
resultado de combinar el bloque que empieza en $i$ y tiene largo
$2^k$. La clave está en la consulta: para el rango $[l,r]$ tomas
$k=\lfloor\log_2(r-l+1)\rfloor$ y lo cubres con \emph{dos} bloques de
largo $2^k$ —uno pegado a $l$ y otro pegado a $r$— que
\textbf{se solapan}. Si la operación es \emph{idempotente}
($f(x,x)=x$, como min, max o gcd), contar el solape dos veces no
cambia el resultado, así que la consulta son solo dos lecturas y una
combinación. Para operaciones \textbf{no idempotentes} (como suma o
xor) hace falta la variante \emph{Disjoint Sparse Table}, que evita el
solape por construcción.

\textbf{¿Por qué usarlo?} Es la estructura más rápida posible para
consultas de rango ($O(1)$ por query) cuando el arreglo es estático.
A cambio de esa velocidad, no soporta actualizaciones: reconstruir
cuesta $O(n\log n)$, así que si el arreglo cambia, conviene un
Segment Tree.

\textbf{Casos de uso:}

\begin{itemize}
    \item RMQ clásico: mínimo o máximo de cualquier rango, arreglo
    estático.
    \item gcd o AND/OR de rango sobre un arreglo estático.
    \item LCA vía Euler Tour + RMQ $\pm1$.
    \item Suma, xor o cualquier operación no idempotente de rango
    estático, en $O(1)$, usando Disjoint Sparse Table.
    \item Cualquier problema con muchísimas queries offline sobre datos
    que nunca se modifican.
\end{itemize}

\textbf{Complejidad:} construir $O(n\log n)$, consulta $O(1)$. (Si el
arreglo \emph{cambia}, usa Segment Tree: $O(\log n)$ por operación.)


\subsubsection{Sparse Table clásico (operaciones idempotentes: min, max, gcd)}

\begin{lstlisting}[style=cpp]
struct SparseTable {
    vector<vector<long long>> sp;
    vector<int> lg;
    int n;

    // CREATE: precalcula la tabla a partir de un arreglo a[0..n-1].
    // sp[k][i] = minimo del bloque [i, i + 2^k - 1].
    SparseTable(vector<long long>& a) {
        n = a.size();

        lg.assign(n + 1, 0);
        for (int i = 2; i <= n; i++)
            lg[i] = lg[i / 2] + 1;

        int K = lg[n] + 1;
        sp.assign(K, vector<long long>(n));
        sp[0] = a;

        for (int k = 1; k < K; k++)
            for (int i = 0; i + (1 << k) <= n; i++)
                sp[k][i] = min(sp[k - 1][i],
                                sp[k - 1][i + (1 << (k - 1))]);
    }                                          // O(n log n)

    // READ: minimo del rango [l, r] (0-indexado, inclusivo).
    // Cubre [l, r] con dos bloques de largo 2^k que se solapan;
    // como min es idempotente, el solape no afecta el resultado.
    long long query(int l, int r) {
        int k = lg[r - l + 1];
        return min(sp[k][l], sp[k][r - (1 << k) + 1]);
    }                                          // O(1)
};
\end{lstlisting}

\textbf{Ejemplo de funcionamiento:}

\begin{lstlisting}[style=cpp]
vector<long long> a = {5, 2, 9, 1, 7, 3};

SparseTable st(a);
// arreglo logico: [5, 2, 9, 1, 7, 3]  (posiciones 0..5)

cout << st.query(0, 2);
// min(5, 2, 9) = 2

cout << st.query(2, 5);
// min(9, 1, 7, 3) = 1

cout << st.query(4, 4);
// min(7) = 7
\end{lstlisting}

\textbf{Nota:} para \texttt{max} o \texttt{gcd} solo cambia
\texttt{min(...)} por \texttt{max(...)} o \texttt{\_\_gcd(...)} en el
constructor y en \texttt{query}. Si necesitas saber \emph{qué índice}
alcanza el mínimo (no solo el valor), guarda índices en vez de
valores y compara \texttt{a[sp[k][i]]} en lugar de \texttt{sp[k][i]}.


\subsubsection{Disjoint Sparse Table (operaciones no idempotentes: suma, xor)}

\textbf{Idea:} en vez de dos bloques que se solapan, cada nivel $k$
particiona el arreglo en bloques de largo $2^{k+1}$ y precalcula, para
cada bloque, un prefijo hacia la izquierda desde su punto medio y un
prefijo hacia la derecha desde ese mismo punto medio. Toda consulta
$[l,r]$ cae exactamente en un nivel donde $l$ y $r$ quedan en mitades
distintas del mismo bloque (el nivel lo da el bit más alto donde
difieren $l$ y $r$), así que se combina sin solapar y sirve para
operaciones \textbf{no idempotentes} como la suma.

\begin{lstlisting}[style=cpp]
struct DisjointSparseTable {
    vector<vector<long long>> pre;
    vector<long long> a;
    int n, LOG;

    // CREATE: precalcula prefijos izquierda/derecha por nivel,
    // partiendo el arreglo en bloques de tamano 2^(k+1).
    DisjointSparseTable(vector<long long>& arr) {
        a = arr;
        n = a.size();
        LOG = 1;
        while ((1 << LOG) < n) LOG++;
        LOG++;

        pre.assign(LOG, vector<long long>(n));

        for (int k = 0; k < LOG; k++) {
            int block = 1 << (k + 1);

            for (int start = 0; start < n; start += block) {
                int mid = min(start + block / 2, n);
                int end = min(start + block, n);

                if (mid > n) continue;

                // prefijo hacia la izquierda desde mid-1
                if (mid - 1 >= start) {
                    pre[k][mid - 1] = a[mid - 1];
                    for (int i = mid - 2; i >= start; i--)
                        pre[k][i] = a[i] + pre[k][i + 1];
                }

                // prefijo hacia la derecha desde mid
                if (mid < end) {
                    pre[k][mid] = a[mid];
                    for (int i = mid + 1; i < end; i++)
                        pre[k][i] = pre[k][i - 1] + a[i];
                }
            }
        }
    }                                          // O(n log n)

    // READ: suma del rango [l, r] (0-indexado, inclusivo).
    // El nivel k lo da el bit mas alto donde l y r difieren:
    // ahi es donde el bloque se parte en dos mitades disjuntas.
    long long query(int l, int r) {
        if (l == r) return a[l];

        int k = 31 - __builtin_clz(l ^ r);
        return pre[k][l] + pre[k][r];
    }                                          // O(1)
};
\end{lstlisting}

\textbf{Ejemplo de funcionamiento:}

\begin{lstlisting}[style=cpp]
vector<long long> a = {5, 2, 9, 1, 7, 3};

DisjointSparseTable dst(a);
// arreglo logico: [5, 2, 9, 1, 7, 3]  (posiciones 0..5)

cout << dst.query(0, 2);
// 5 + 2 + 9 = 16

cout << dst.query(2, 5);
// 9 + 1 + 7 + 3 = 20

cout << dst.query(1, 1);
// 2
\end{lstlisting}

\textbf{Diferencia entre las dos variantes:}

\begin{table}[H]
\raggedright
\scriptsize
\setlength{\tabcolsep}{3pt}
\begin{tabular}{|l|c|c|c|}
\hline
\textbf{Estructura} & \textbf{Construir} & \textbf{Query} & \textbf{Updates} \\
\hline
Sparse Table (min/max/gcd) & $O(n\log n)$ & $O(1)$ & No \\
\hline
Disjoint Sparse Table (suma/xor) & $O(n\log n)$ & $O(1)$ & No \\
\hline
Segment Tree & $O(n)$ & $O(\log n)$ & Sí, $O(\log n)$ \\
\hline
\end{tabular}
\caption{Sparse Table vs. Disjoint Sparse Table vs. Segment Tree.}
\end{table}

\textbf{Regla práctica:} si el arreglo \textbf{no cambia}, usa Sparse
Table para min/max/gcd, o Disjoint Sparse Table si necesitas suma,
xor, o cualquier operación no idempotente. En cuanto el arreglo
\textbf{cambia} —aunque sea una sola actualización— reconstruir en
$O(n\log n)$ deja de valer la pena y conviene un Segment Tree.




\subsection{Segment Tree}

\textbf{Idea:} un árbol binario implícito (guardado en un arreglo)
donde cada nodo representa el resultado combinado de un rango del
arreglo original. La raíz cubre todo $[0,n-1]$; cada nodo se parte en
dos mitades hasta llegar a las hojas, que son elementos individuales.
Un update baja por el árbol tocando $O(\log n)$ nodos (uno por nivel);
una query de rango se descompone en a lo más $O(\log n)$ nodos que
cubren exactamente el rango pedido, sin necesidad de bloques
parciales como en Sqrt Decomposition. A diferencia de Fenwick, sirve
para \textbf{cualquier} operación asociativa (min, max, gcd, and, or,
no solo sumas), y con Lazy Propagation soporta actualizaciones de
rango completo.

\textbf{¿Por qué usarlo?} Es la estructura más versátil de esta
sección: $O(\log n)$ tanto para update como para query, funciona con
casi cualquier operación asociativa, y con lazy propagation resuelve
range update + range query para operaciones que Fenwick no puede
(min, max, asignar un valor a todo un rango, etc.).

\textbf{Casos de uso:}

\begin{itemize}
    \item Min/max/gcd de rango con actualizaciones (a diferencia de
    Sparse Table, que exige arreglo estático).
    \item Range update + range query con operaciones no invertibles:
    asignar un valor a todo un rango, min/max con suma en rango, etc.
    \item Cualquier operación asociativa que no sea una suma simple
    (donde Fenwick ya alcanzaría y sería más simple).
    \item Base para variantes más avanzadas: Segment Tree 2D, Merge
    Sort Tree, Persistent Segment Tree, Segment Tree Beats.
\end{itemize}

\textbf{Complejidad:} construir $O(n)$, update puntual $O(\log n)$,
query de rango $O(\log n)$, range update + range query con lazy
propagation $O(\log n)$.


\subsubsection{Segment Tree clásico (point update, range query — suma)}

\begin{lstlisting}[style=cpp]
struct SegmentTree {
    vector<long long> tree, a;
    int n;

    // CREATE: reserva el arbol (tamano 4n es cota segura) y
    // construye recursivamente a partir de a[0..n-1].
    SegmentTree(vector<long long>& arr) {
        a = arr;
        n = a.size();
        tree.assign(4 * n, 0);
        if (n > 0)
            build(1, 0, n - 1);
    }                                          // O(n)

    void build(int node, int l, int r) {
        if (l == r) {
            tree[node] = a[l];
            return;
        }

        int mid = (l + r) / 2;
        build(2 * node, l, mid);
        build(2 * node + 1, mid + 1, r);
        tree[node] = tree[2 * node] + tree[2 * node + 1];
    }

    void updateUtil(int node, int l, int r, int i, long long x) {
        if (l == r) {
            tree[node] = x;
            return;
        }

        int mid = (l + r) / 2;
        if (i <= mid)
            updateUtil(2 * node, l, mid, i, x);
        else
            updateUtil(2 * node + 1, mid + 1, r, i, x);

        tree[node] = tree[2 * node] + tree[2 * node + 1];
    }

    // UPDATE: fija el valor absoluto de la posicion i en x.
    void update(int i, long long x) {
        updateUtil(1, 0, n - 1, i, x);
    }                                          // O(log n)

    long long queryUtil(int node, int l, int r, int ql, int qr) {
        if (qr < l || r < ql)
            return 0;                          // neutro para suma

        if (ql <= l && r <= qr)
            return tree[node];

        int mid = (l + r) / 2;
        return queryUtil(2 * node, l, mid, ql, qr)
             + queryUtil(2 * node + 1, mid + 1, r, ql, qr);
    }

    // READ: suma del rango [l, r] (0-indexado, inclusivo).
    long long rango(int l, int r) {
        return queryUtil(1, 0, n - 1, l, r);
    }                                          // O(log n)
};
\end{lstlisting}

\textbf{Ejemplo de funcionamiento:}

\begin{lstlisting}[style=cpp]
vector<long long> a = {5, 2, 9, 1, 7};

SegmentTree st(a);
// arreglo logico: [5, 2, 9, 1, 7]  (posiciones 0..4)

cout << st.rango(1, 3);
// 2 + 9 + 1 = 12

st.update(2, 100);
// arreglo logico: [5, 2, 100, 1, 7]

cout << st.rango(1, 3);
// 2 + 100 + 1 = 103
\end{lstlisting}

\textbf{Nota:} para min/max/gcd solo cambia el neutro del caso base
(\texttt{LLONG\_MAX} para min, \texttt{0} para gcd, etc.) y el
operador de combinación (\texttt{min(...)}, \texttt{\_\_gcd(...)}) en
\texttt{build}, \texttt{updateUtil} y \texttt{queryUtil}.


\subsubsection{Segment Tree con Lazy Propagation (range update + range query — suma)}

\textbf{Idea:} cada nodo guarda un valor \texttt{lazy} pendiente que
representa ``sumar $v$ a todo este rango'' sin aplicarlo todavía a
los hijos. Se aplica (\emph{se propaga}) solo cuando hace falta bajar
más en el árbol para otra operación.

\begin{lstlisting}[style=cpp]
struct SegmentTreeLazyAdd {
    vector<long long> tree, lazy, a;
    int n;

    // CREATE: reserva arbol + lazy, construye a partir de a[0..n-1].
    SegmentTreeLazyAdd(vector<long long>& arr) {
        a = arr;
        n = a.size();
        tree.assign(4 * n, 0);
        lazy.assign(4 * n, 0);
        if (n > 0)
            build(1, 0, n - 1);
    }                                          // O(n)

    void build(int node, int l, int r) {
        if (l == r) {
            tree[node] = a[l];
            return;
        }

        int mid = (l + r) / 2;
        build(2 * node, l, mid);
        build(2 * node + 1, mid + 1, r);
        tree[node] = tree[2 * node] + tree[2 * node + 1];
    }

    // Aplica el lazy pendiente de este nodo y lo empuja a los hijos.
    void propagar(int node, int l, int r) {
        if (lazy[node] == 0)
            return;

        tree[node] += lazy[node] * (r - l + 1);

        if (l != r) {
            lazy[2 * node] += lazy[node];
            lazy[2 * node + 1] += lazy[node];
        }

        lazy[node] = 0;
    }

    void rangoUpdateUtil(int node, int l, int r, int ql, int qr,
                          long long v) {
        propagar(node, l, r);

        if (qr < l || r < ql)
            return;

        if (ql <= l && r <= qr) {
            lazy[node] += v;
            propagar(node, l, r);
            return;
        }

        int mid = (l + r) / 2;
        rangoUpdateUtil(2 * node, l, mid, ql, qr, v);
        rangoUpdateUtil(2 * node + 1, mid + 1, r, ql, qr, v);
        tree[node] = tree[2 * node] + tree[2 * node + 1];
    }

    // UPDATE: suma v a todo el rango [l, r].
    void rangoUpdate(int l, int r, long long v) {
        rangoUpdateUtil(1, 0, n - 1, l, r, v);
    }                                          // O(log n)

    long long queryUtil(int node, int l, int r, int ql, int qr) {
        propagar(node, l, r);

        if (qr < l || r < ql)
            return 0;

        if (ql <= l && r <= qr)
            return tree[node];

        int mid = (l + r) / 2;
        return queryUtil(2 * node, l, mid, ql, qr)
             + queryUtil(2 * node + 1, mid + 1, r, ql, qr);
    }

    // READ: suma del rango [l, r].
    long long rango(int l, int r) {
        return queryUtil(1, 0, n - 1, l, r);
    }                                          // O(log n)
};
\end{lstlisting}

\textbf{Ejemplo de funcionamiento:}

\begin{lstlisting}[style=cpp]
vector<long long> a = {1, 1, 1, 1, 1, 1};

SegmentTreeLazyAdd st(a);

st.rangoUpdate(1, 4, 10);
// arreglo logico efectivo: [1, 11, 11, 11, 11, 1]

cout << st.rango(0, 5);
// 1+11+11+11+11+1 = 46

cout << st.rango(2, 3);
// 11+11 = 22
\end{lstlisting}


\subsubsection{Lazy Propagation con asignación (range assign + range query — mínimo)}

\textbf{Idea:} en vez de \emph{sumar} un valor al rango, lo
\textbf{reemplaza} por completo (\texttt{asigna x a todo
$[l,r]$}). El lazy aquí no se acumula como en la suma: el último
assign pisa a cualquier assign anterior pendiente, así que se necesita
una bandera aparte para distinguir ``no hay lazy pendiente'' de
``hay que asignar 0''.

\begin{lstlisting}[style=cpp]
struct SegmentTreeAssignMin {
    vector<long long> tree, lazy, a;
    vector<bool> tieneLazy;
    int n;

    // CREATE: reserva arbol + lazy + banderas, construye desde a.
    SegmentTreeAssignMin(vector<long long>& arr) {
        a = arr;
        n = a.size();
        tree.assign(4 * n, LLONG_MAX);
        lazy.assign(4 * n, 0);
        tieneLazy.assign(4 * n, false);
        if (n > 0)
            build(1, 0, n - 1);
    }                                          // O(n)

    void build(int node, int l, int r) {
        if (l == r) {
            tree[node] = a[l];
            return;
        }

        int mid = (l + r) / 2;
        build(2 * node, l, mid);
        build(2 * node + 1, mid + 1, r);
        tree[node] = min(tree[2 * node], tree[2 * node + 1]);
    }

    // Empuja el assign pendiente de este nodo hacia sus hijos.
    void propagar(int node) {
        if (!tieneLazy[node])
            return;

        for (int hijo : {2 * node, 2 * node + 1}) {
            tree[hijo] = lazy[node];
            lazy[hijo] = lazy[node];
            tieneLazy[hijo] = true;
        }

        tieneLazy[node] = false;
    }

    void rangoAssignUtil(int node, int l, int r, int ql, int qr,
                          long long x) {
        if (qr < l || r < ql)
            return;

        if (ql <= l && r <= qr) {
            tree[node] = x;
            lazy[node] = x;
            tieneLazy[node] = true;
            return;
        }

        propagar(node);
        int mid = (l + r) / 2;
        rangoAssignUtil(2 * node, l, mid, ql, qr, x);
        rangoAssignUtil(2 * node + 1, mid + 1, r, ql, qr, x);
        tree[node] = min(tree[2 * node], tree[2 * node + 1]);
    }

    // UPDATE: asigna x a todo el rango [l, r].
    void rangoAssign(int l, int r, long long x) {
        rangoAssignUtil(1, 0, n - 1, l, r, x);
    }                                          // O(log n)

    long long queryUtil(int node, int l, int r, int ql, int qr) {
        if (qr < l || r < ql)
            return LLONG_MAX;

        if (ql <= l && r <= qr)
            return tree[node];

        propagar(node);
        int mid = (l + r) / 2;
        return min(queryUtil(2 * node, l, mid, ql, qr),
                    queryUtil(2 * node + 1, mid + 1, r, ql, qr));
    }

    // READ: minimo del rango [l, r].
    long long rango(int l, int r) {
        return queryUtil(1, 0, n - 1, l, r);
    }                                          // O(log n)
};
\end{lstlisting}

\textbf{Ejemplo de funcionamiento:}

\begin{lstlisting}[style=cpp]
vector<long long> a = {5, 2, 9, 1, 7, 3};

SegmentTreeAssignMin st(a);

cout << st.rango(0, 5);
// min(5,2,9,1,7,3) = 1

st.rangoAssign(0, 2, 100);
// arreglo logico efectivo: [100, 100, 100, 1, 7, 3]

cout << st.rango(0, 2);
// min(100,100,100) = 100

cout << st.rango(0, 5);
// min(100,100,100,1,7,3) = 1
\end{lstlisting}

\textbf{Comparación con las otras estructuras de rango:}

\begin{table}[H]
\raggedright
\scriptsize
\setlength{\tabcolsep}{3pt}
\begin{tabular}{|l|c|c|c|}
\hline
\textbf{Estructura} & \textbf{Update} & \textbf{Query} & \textbf{Min/Max?} \\
\hline
Segment Tree & $O(\log n)$ & $O(\log n)$ & Sí \\
\hline
Segment Tree + Lazy & $O(\log n)$ range & $O(\log n)$ range & Sí \\
\hline
\texttt{FenwickTree} & $O(\log n)$ punto & $O(\log n)$ rango & No \\
\hline
Sqrt Decomposition & $O(1)$–$O(\sqrt n)$ & $O(\sqrt n)$ & Sí \\
\hline
\end{tabular}
\caption{Segment Tree vs. Fenwick vs. Sqrt Decomposition.}
\end{table}

\textbf{Regla práctica:} usa Segment Tree cuando necesites min/max/gcd
(operaciones no invertibles que Fenwick no soporta), o cuando
necesites range update + range query con una operación que no sea
suma simple. Si la operación es solo suma y no necesitas range
update, Fenwick es más simple y más rápido en la práctica (menor
constante). Si no confías en escribir el lazy propagation correcto
bajo presión, Sqrt Decomposition es una alternativa más simple a
costa de $O(\sqrt n)$ en vez de $O(\log n)$.


\subsection{RMQ (Range Minimum Query)}

\textbf{Idea:} RMQ es el problema general de "dado un arreglo,
encuentra el mínimo en cualquier rango $[l,r]$". Ya tienes las dos
estructuras que lo resuelven: \textbf{Sparse Table} si el arreglo es
estático ($O(1)$ por query) y \textbf{Segment Tree} si cambia
($O(\log n)$ por query). Esta sección no repite esas estructuras,
sino que muestra las dos aplicaciones clásicas donde un problema que
\emph{no parece} de rango se puede \textbf{reformular} como RMQ sobre
un arreglo auxiliar: LCA en árboles, y LCP entre sufijos arbitrarios
de un string.

\textbf{¿Por qué usarlo?} Reconocer cuándo un problema es "RMQ
disfrazado" es lo que realmente separa resolver estos problemas en
$O(1)$–$O(\log n)$ de resolverlos con fuerza bruta $O(n)$ por
consulta. El patrón se repite: construyes un arreglo auxiliar donde
la respuesta a tu pregunta real equivale al mínimo de un rango, y le
aplicas Sparse Table o Segment Tree normal.

\textbf{Casos de uso:}

\begin{itemize}
    \item LCA (ancestro común más bajo) de dos nodos en un árbol, con
    muchas queries: Euler Tour + RMQ.
    \item LCP (prefijo común más largo) entre dos sufijos
    \emph{cualquiera} de un string (no necesariamente consecutivos en
    el Suffix Array): Suffix Array + Kasai + RMQ sobre el LCP Array.
    \item Cualquier problema donde la respuesta se pueda expresar como
    "el mínimo de algo" sobre un rango de índices auxiliares.
\end{itemize}

\textbf{Complejidad:} en ambas aplicaciones, el precómputo (Euler
Tour o Suffix Array + Kasai) es $O(n)$ a $O(n\log^2 n)$ según el
método de Suffix Array usado; una vez armado el arreglo auxiliar, el
RMQ es $O(n\log n)$ de build y $O(1)$ por query con Sparse Table
(ambos casos son estáticos: el árbol y el string no cambian).


\subsubsection{LCA vía Euler Tour + RMQ (Sparse Table)}

\textbf{Idea:} recorres el árbol con un Euler Tour, anotando el nodo
visitado cada vez que \textbf{entras o vuelves} a él (arreglo de
tamaño $2n-1$) junto con su profundidad. El LCA de $u$ y $v$ es el
nodo de \textbf{menor profundidad} que aparece entre la primera
aparición de $u$ y la primera aparición de $v$ en ese recorrido —
exactamente un RMQ (por profundidad) sobre el arreglo del Euler Tour.

\begin{lstlisting}[style=cpp]
struct LCA {
    vector<int> euler, primeraApar, profundidad;
    vector<vector<int>> sp;   // sp[k][i] = nodo con menor profundidad
    vector<int> lg;

    // CREATE: construye el Euler Tour del arbol (arr. de tamano
    // 2n-1) y la Sparse Table sobre profundidad, para RMQ O(1).
    LCA(vector<vector<int>>& adj, int raiz, int n) {
        primeraApar.assign(n, -1);
        profundidad.assign(n, 0);
        euler.reserve(2 * n);

        dfs(adj, raiz, -1, 0);
        build();
    }                                          // O(n log n)

    void dfs(vector<vector<int>>& adj, int u, int p, int d) {
        primeraApar[u] = euler.size();
        profundidad[u] = d;
        euler.push_back(u);

        for (int v : adj[u]) {
            if (v != p) {
                dfs(adj, v, u, d + 1);
                euler.push_back(u);
            }
        }
    }

    void build() {
        int m = euler.size();
        lg.assign(m + 1, 0);
        for (int i = 2; i <= m; i++)
            lg[i] = lg[i / 2] + 1;

        int K = lg[m] + 1;
        sp.assign(K, vector<int>(m));
        sp[0] = euler;

        for (int k = 1; k < K; k++)
            for (int i = 0; i + (1 << k) <= m; i++) {
                int l = sp[k - 1][i];
                int r = sp[k - 1][i + (1 << (k - 1))];
                sp[k][i] = profundidad[l] < profundidad[r] ? l : r;
            }
    }

    // READ: LCA de u y v.
    int query(int u, int v) {
        int l = primeraApar[u], r = primeraApar[v];
        if (l > r) swap(l, r);

        int k = lg[r - l + 1];
        int a = sp[k][l], b = sp[k][r - (1 << k) + 1];
        return profundidad[a] < profundidad[b] ? a : b;
    }                                          // O(1)
};
\end{lstlisting}

\textbf{Ejemplo de funcionamiento:}

\begin{lstlisting}[style=cpp]
// arbol:      0
//            / \
//           1   2
//          / \
//         3   4

vector<vector<int>> adj(5);
adj[0] = {1, 2};
adj[1] = {0, 3, 4};
adj[2] = {0};
adj[3] = {1};
adj[4] = {1};

LCA lca(adj, 0, 5);

cout << lca.query(3, 4);
// 1  (padre comun de 3 y 4)

cout << lca.query(3, 2);
// 0  (raiz)

cout << lca.query(1, 4);
// 1  (uno es ancestro del otro)
\end{lstlisting}


\subsubsection{LCP entre dos sufijos arbitrarios (Suffix Array + Kasai + RMQ)}

\textbf{Idea:} el LCP Array (via Kasai) solo te da el LCP entre
sufijos \textbf{consecutivos} en el Suffix Array. Para el LCP entre
dos sufijos \textbf{cualquiera} $i,j$: los ubicas por su rango en el
SA ($\texttt{rnk}[i]$, $\texttt{rnk}[j]$) y el LCP entre ambos es el
\textbf{mínimo} del LCP Array en el rango entre esas dos posiciones
—otra vez RMQ, esta vez sobre el LCP Array.

\begin{lstlisting}[style=cpp]
struct SuffixLCPQuery {
    vector<int> sa, rnk, lcp;
    vector<vector<int>> sp;
    vector<int> lg;
    int n;

    // CREATE: construye SA, LCP (Kasai) y Sparse Table sobre lcp,
    // todo en un solo paso.
    SuffixLCPQuery(const string& s) {
        n = s.size();
        sa = buildSuffixArray(s);             // O(n log^2 n)
        lcp = buildLCPArray(s, sa);           // O(n), Kasai

        rnk.assign(n, 0);
        for (int i = 0; i < n; i++)
            rnk[sa[i]] = i;

        buildSparse();
    }                                          // O(n log^2 n)

    void buildSparse() {
        lg.assign(n + 1, 0);
        for (int i = 2; i <= n; i++)
            lg[i] = lg[i / 2] + 1;

        int K = lg[n] + 1;
        sp.assign(K, vector<int>(n));
        sp[0] = lcp;

        for (int k = 1; k < K; k++)
            for (int i = 0; i + (1 << k) <= n; i++)
                sp[k][i] = min(sp[k - 1][i],
                                sp[k - 1][i + (1 << (k - 1))]);
    }

    // READ: LCP entre el sufijo que empieza en i y el que
    // empieza en j (0-indexado sobre el string original).
    int lcpEntreSufijos(int i, int j) {
        if (i == j) return n - i;

        int a = rnk[i], b = rnk[j];
        if (a > b) swap(a, b);

        // el minimo va de [a+1, b] en el LCP array
        int k = lg[b - a];
        return min(sp[k][a + 1], sp[k][b - (1 << k) + 1]);
    }                                          // O(1)
};
\end{lstlisting}

\textbf{Ejemplo de funcionamiento:}

\begin{lstlisting}[style=cpp]
SuffixLCPQuery slq("banana");

cout << slq.lcpEntreSufijos(1, 3);
// sufijo desde 1: "anana"
// sufijo desde 3: "ana"
// LCP = "ana" -> 3

cout << slq.lcpEntreSufijos(0, 3);
// sufijo desde 0: "banana"
// sufijo desde 3: "ana"
// LCP = "" -> 0
\end{lstlisting}

\textbf{Nota:} \texttt{buildSuffixArray} y \texttt{buildLCPArray} son
las funciones estándar de Suffix Array + Kasai (ver la sección de
Algoritmos de Strings); aquí solo se reutilizan como precómputo antes
de aplicar RMQ.

\textbf{Regla práctica:} cuando un problema pida "algo" entre dos
elementos que no son adyacentes en su estructura original (dos nodos
de un árbol, dos sufijos de un string, dos elementos de cualquier
secuencia con una noción de "orden" y "profundidad/afinidad"),
pregúntate si esa relación se puede expresar como el mínimo de un
arreglo auxiliar construido con un recorrido lineal. Si sí, el
problema se reduce a Sparse Table (estático) o Segment Tree
(dinámico), ya cubiertos en las secciones anteriores.


\subsection{Mo's Algorithm}

\textbf{Idea:} responde $q$ queries de rango \textbf{offline} (las
conoces todas de antemano) en $O((n+q)\sqrt n)$ total, sin ninguna
estructura de datos elaborada. La clave es \textbf{ordenar las
queries} —no procesarlas en el orden dado— agrupándolas por bloque de
tamaño $\sqrt n$ según su extremo izquierdo $l$, y dentro de cada
bloque por su extremo derecho $r$ (alternando la dirección en bloques
pares/impares para aprovechar mejor la localidad). Luego mantienes un
puntero $[curL, curR]$ que se \textbf{desliza} de una query a la
siguiente con funciones \texttt{add}/\texttt{remove} en $O(1)$ cada
una. Como los punteros nunca retroceden más de lo necesario gracias al
orden elegido, el total de movimientos amortizado es
$O((n+q)\sqrt n)$ en vez de $O(nq)$ con fuerza bruta.

\textbf{¿Por qué usarlo?} Sirve para respuestas que \textbf{no tienen
una estructura de rango obvia} (por ejemplo: contar elementos
distintos, sumar cuadrados de frecuencias, contar pares con cierta
propiedad) donde escribir un Segment Tree específico sería
complicado o imposible, pero sí es fácil mantener la respuesta
incrementalmente al agregar/quitar un elemento a la vez.

\textbf{Casos de uso:}

\begin{itemize}
    \item Contar elementos distintos en un rango, con muchas queries.
    \item Sumas o conteos que dependen de \emph{frecuencias} dentro
    del rango (no solo de valores individuales), donde \texttt{add}/
    \texttt{remove} son fáciles pero no hay forma clara de combinar
    dos sub-rangos con un Segment Tree.
    \item Todas las queries se conocen de antemano (offline). Si
    necesitas responder \textbf{online} (una query a la vez, sin ver
    las siguientes), Mo's no aplica: usa Segment Tree o Sqrt
    Decomposition.
    \item El arreglo tiene pocos updates puntuales intercalados con
    las queries: usa la variante \emph{Mo's con updates} (3D Mo's).
\end{itemize}

\textbf{Complejidad:} $O((n+q)\sqrt n)$ total para $q$ queries sin
updates; $O(n^{5/3})$ (aprox.) para la variante con updates, usando
bloques de tamaño $n^{2/3}$.


\subsubsection{Mo's clásico (queries offline, sin updates)}

\textbf{Ejemplo guía:} contar cuántos valores \textbf{distintos} hay
en cada rango $[l,r]$ preguntado.

\begin{lstlisting}[style=cpp]
struct MoAlgorithm {
    int n, blockSize;
    vector<int> a;
    vector<long long> cnt;
    int curL = 0, curR = -1;
    long long curDistinct = 0;

    // CREATE: reserva el arreglo de frecuencias segun el maximo
    // valor posible, y fija el tamano de bloque en ~sqrt(n).
    MoAlgorithm(vector<int>& arr, int maxVal) {
        a = arr;
        n = a.size();
        blockSize = max(1, (int)sqrt((double)n));
        cnt.assign(maxVal, 0);
    }                                          // O(n)

    // UPDATE: agrega el elemento en pos a la ventana actual.
    void add(int pos) {
        if (cnt[a[pos]]++ == 0)
            curDistinct++;
    }                                          // O(1)

    // UPDATE: quita el elemento en pos de la ventana actual.
    void remove(int pos) {
        if (--cnt[a[pos]] == 0)
            curDistinct--;
    }                                          // O(1)

    // READ: procesa todas las queries offline (l, r inclusive,
    // 0-indexadas) y devuelve las respuestas en el orden original.
    vector<long long> resolver(vector<array<int, 2>>& queries) {
        int q = queries.size();
        vector<int> idx(q);
        iota(idx.begin(), idx.end(), 0);

        sort(idx.begin(), idx.end(), [&](int i, int j) {
            int bi = queries[i][0] / blockSize;
            int bj = queries[j][0] / blockSize;
            if (bi != bj) return bi < bj;
            return (bi & 1) ? (queries[i][1] > queries[j][1])
                             : (queries[i][1] < queries[j][1]);
        });

        vector<long long> respuestas(q);
        for (int id : idx) {
            int l = queries[id][0], r = queries[id][1];

            while (curR < r) add(++curR);
            while (curL > l) add(--curL);
            while (curR > r) remove(curR--);
            while (curL < l) remove(curL++);

            respuestas[id] = curDistinct;
        }
        return respuestas;
    }                                          // O((n+q) sqrt(n))
};
\end{lstlisting}

\textbf{Ejemplo de funcionamiento:}

\begin{lstlisting}[style=cpp]
vector<int> a = {1, 2, 1, 3, 2, 4, 1};
// valores en [1, 4]

MoAlgorithm mo(a, 5);

vector<array<int, 2>> queries = {
    {0, 2},   // [1, 2, 1] -> 2 distintos
    {1, 4},   // [2, 1, 3, 2] -> 3 distintos
    {0, 6}    // arreglo completo -> 4 distintos
};

vector<long long> resp = mo.resolver(queries);
// resp = [2, 3, 4]
\end{lstlisting}


\subsubsection{Mo's con updates (3D Mo's)}

\textbf{Idea:} si además de queries hay \textbf{updates puntuales}
intercaladas, se agrega una tercera dimensión al ordenamiento: el
\emph{tiempo} $t$ (cuántas updates ya ocurrieron antes de cada
query). Se ordenan las queries por
$(\text{bloque}(l), \text{bloque}(r), t)$, y además de deslizar
$[curL,curR]$, se avanza o retrocede $curT$ aplicando o revirtiendo
updates una por una. El bloque óptimo cambia a
$\approx n^{2/3}$ (en vez de $\sqrt n$) porque ahora hay tres
dimensiones que balancear.

\begin{lstlisting}[style=cpp]
struct MoUpdates {
    int n, blockSize;
    vector<int> a;
    vector<long long> cnt;
    vector<pair<int, int>> updates;   // (pos, valor); ver nota abajo
    int curL = 0, curR = -1, curT = 0;
    long long curDistinct = 0;

    // CREATE: reserva frecuencias y fija el bloque en ~n^(2/3).
    MoUpdates(vector<int>& arr, int maxVal) {
        a = arr;
        n = a.size();
        blockSize = max(1, (int)cbrt((double)n) *
                            (int)cbrt((double)n));
        cnt.assign(maxVal, 0);
    }                                          // O(n)

    void add(int pos) {
        if (cnt[a[pos]]++ == 0) curDistinct++;
    }                                          // O(1)

    void remove(int pos) {
        if (--cnt[a[pos]] == 0) curDistinct--;
    }                                          // O(1)

    // CREATE: registra una update (aplicarla no es inmediato,
    // solo queda en la lista para procesarse en su momento).
    void agregarUpdate(int pos, int val) {
        updates.push_back({pos, val});
    }                                          // O(1)

    // UPDATE: aplica (o revierte, es la misma operacion) la
    // update numero t. El truco: como updates[t].second guarda
    // "el valor a poner", tras aplicarla queda guardado el valor
    // viejo -> llamar de nuevo revierte exactamente al estado
    // anterior.
    void alternarUpdate(int t) {
        int pos = updates[t].first;
        int val = updates[t].second;

        if (pos >= curL && pos <= curR)
            remove(pos);

        swap(a[pos], val);
        updates[t].second = val;

        if (pos >= curL && pos <= curR)
            add(pos);
    }                                          // O(1)

    // READ: procesa todas las queries offline (l, r, t) y
    // devuelve las respuestas en el orden original.
    vector<long long> resolver(vector<array<int, 3>>& queries) {
        int q = queries.size();
        vector<int> idx(q);
        iota(idx.begin(), idx.end(), 0);

        sort(idx.begin(), idx.end(), [&](int i, int j) {
            int bi = queries[i][0] / blockSize;
            int bj = queries[j][0] / blockSize;
            if (bi != bj) return bi < bj;

            int ri = queries[i][1] / blockSize;
            int rj = queries[j][1] / blockSize;
            if (ri != rj) return ri < rj;

            return queries[i][2] < queries[j][2];
        });

        vector<long long> respuestas(q);
        for (int id : idx) {
            int l = queries[id][0], r = queries[id][1],
                t = queries[id][2];

            while (curR < r) add(++curR);
            while (curL > l) add(--curL);
            while (curR > r) remove(curR--);
            while (curL < l) remove(curL++);

            while (curT < t) alternarUpdate(curT++);
            while (curT > t) alternarUpdate(--curT);

            respuestas[id] = curDistinct;
        }
        return respuestas;
    }                                          // O(n^(5/3)) aprox.
};
\end{lstlisting}

\textbf{Ejemplo de funcionamiento:}

\begin{lstlisting}[style=cpp]
vector<int> a = {1, 2, 3};

MoUpdates mo(a, 5);
mo.agregarUpdate(0, 2);   // update #0: a[0] = 2
mo.agregarUpdate(1, 1);   // update #1: a[1] = 1

// queries: (l, r, cuantas updates ya se aplicaron)
vector<array<int, 3>> queries = {
    {0, 2, 0},   // arreglo original [1,2,3] -> 3 distintos
    {0, 2, 1},   // tras update #0:  [2,2,3] -> 2 distintos
    {0, 2, 2}    // tras update #1:  [2,1,3] -> 3 distintos
};

vector<long long> resp = mo.resolver(queries);
// resp = [3, 2, 3]
\end{lstlisting}

\textbf{Nota:} para problemas sobre \textbf{árboles} (queries sobre
un camino o subárbol) existe \emph{Mo's on Trees}, que aplana el
árbol con un Euler Tour y aplica el mismo algoritmo sobre ese
arreglo; el resto del razonamiento (ordenar queries, deslizar
punteros) es idéntico.

\textbf{Comparación con otras estructuras de rango:}

\begin{table}[H]
\raggedright
\scriptsize
\setlength{\tabcolsep}{3pt}
\renewcommand{\arraystretch}{1.1}
\begin{tabular}{|p{0.43\columnwidth}|p{0.30\columnwidth}|p{0.20\columnwidth}|}
\hline
\textbf{Situación} & \textbf{Estructura} & \textbf{Complejidad} \\
\hline
Queries offline, sin updates &
Mo's clásico &
$O((n+q)\sqrt{n})$ \\
\hline
Queries offline, con updates &
Mo's con updates &
$O(n^{5/3})$ \\
\hline
Queries online (una a la vez) &
Segment Tree / Sqrt Decomp. &
$O(\log n)$ / $O(\sqrt{n})$ \\
\hline
\end{tabular}
\caption{Cuándo usar Mo's Algorithm frente a estructuras online.}
\end{table}

\textbf{Regla práctica:} usa Mo's solo cuando (1) todas las queries
son offline, y (2) la respuesta es fácil de mantener
incrementalmente con \texttt{add}/\texttt{remove}, pero difícil de
combinar con un Segment Tree (min/max/suma sí se combinan fácil, así
que ahí Mo's no aporta nada). Si el problema admite Segment Tree o
Fenwick directamente, esas opciones son más simples y no necesitan
que las queries sean offline.

\subsection{Wavelet Tree}

\textbf{Idea:} es un árbol binario sobre el \textbf{rango de valores}
del arreglo (no sobre las posiciones, como Segment Tree). La raíz
cubre $[\text{lo},\text{hi}]$ de valores posibles; en cada nodo, los
elementos con valor $\le$ mediana van al hijo izquierdo y el resto al
derecho, \textbf{preservando su orden relativo} (partición estable).
Cada nodo guarda un arreglo \texttt{b[]} con cuántos elementos, hasta
cada posición, fueron hacia la izquierda. Con eso puedes, en cada
nivel, saber exactamente qué sub-rango de posiciones le corresponde a
cada hijo — y bajar por el árbol respondiendo preguntas sobre
\textbf{valores dentro de un rango de posiciones} en $O(\log C)$,
donde $C$ es la cantidad de valores distintos.

\textbf{¿Por qué usarlo?} Es la única estructura de esta sección que
responde eficientemente preguntas que combinan \textbf{posición y
valor} a la vez: "el $k$-ésimo menor en $[l,r]$", "cuántos elementos
en $[l,r]$ son $\le X$", "cuántas veces aparece exactamente $X$ en
$[l,r]$". Ninguna de las estructuras anteriores (Fenwick, Segment
Tree, Sparse Table) resuelve esto directamente sin re-plantear el
problema.

\textbf{Casos de uso:}

\begin{itemize}
    \item $k$-ésimo elemento más chico (o más grande) dentro de un
    rango $[l,r]$.
    \item Conteo de elementos en $[l,r]$ cuyo valor está en un rango
    $[a,b]$.
    \item Frecuencia de un valor específico dentro de un rango de
    posiciones (rank query).
    \item Arreglo \textbf{estático} — no soporta updates sin
    reconstruir (para eso existe la variante persistente, más
    avanzada).
\end{itemize}

\textbf{Complejidad:} construir $O(n\log C)$, cada query
$O(\log C)$, donde $C$ es la cantidad de valores distintos tras
comprimir coordenadas.


\subsubsection{Wavelet Tree: $k$-ésimo, conteo $\le X$ y frecuencia}

\begin{lstlisting}[style=cpp]
struct WaveletTree {
    int lo, hi;
    WaveletTree *izq = nullptr, *der = nullptr;
    vector<int> b;   // b[i] = cuantos de a[0..i-1] fueron a la izq.

    // CREATE: construye recursivamente sobre el segmento [from, to)
    // de un arreglo mutable, particionando por la mediana de valor
    // [lo, hi]. Los valores deben venir ya comprimidos a [0, C-1].
    WaveletTree(int* from, int* to, int lo, int hi) : lo(lo), hi(hi) {
        if (lo == hi || from >= to) return;

        int mid = (lo + hi) / 2;
        auto vaIzq = [mid](int x) { return x <= mid; };

        b.reserve(to - from + 1);
        b.push_back(0);
        for (auto it = from; it != to; it++)
            b.push_back(b.back() + vaIzq(*it));

        auto pivote = stable_partition(from, to, vaIzq);
        izq = new WaveletTree(from, pivote, lo, mid);
        der = new WaveletTree(pivote, to, mid + 1, hi);
    }                                          // O(n log C)

    // READ: k-esimo valor mas chico en [l, r] (1-indexado, k desde 1).
    int kth(int l, int r, int k) {
        if (l > r) return -1;
        if (lo == hi) return lo;

        int enIzq = b[r] - b[l - 1];
        int lb = b[l - 1], rb = b[r];

        if (k <= enIzq)
            return izq->kth(lb + 1, rb, k);
        return der->kth(l - lb, r - rb, k - enIzq);
    }                                          // O(log C)

    // READ: cuantos elementos de [l, r] tienen valor <= val.
    int contarMenorIgual(int l, int r, int val) {
        if (l > r || val < lo) return 0;
        if (hi <= val) return r - l + 1;

        int lb = b[l - 1], rb = b[r];
        return izq->contarMenorIgual(lb + 1, rb, val)
             + der->contarMenorIgual(l - lb, r - rb, val);
    }                                          // O(log C)

    // READ: cuantas veces aparece exactamente val en [l, r].
    int contarIgual(int l, int r, int val) {
        if (l > r || val < lo || val > hi) return 0;
        if (lo == hi) return r - l + 1;

        int mid = (lo + hi) / 2;
        int lb = b[l - 1], rb = b[r];

        if (val <= mid)
            return izq->contarIgual(lb + 1, rb, val);
        return der->contarIgual(l - lb, r - rb, val);
    }                                          // O(log C)

    ~WaveletTree() {
        delete izq;
        delete der;
    }
};
\end{lstlisting}

\textbf{Ejemplo de funcionamiento} (con compresión de coordenadas):

\begin{lstlisting}[style=cpp]
vector<int> original = {40, 10, 30, 10, 20, 40, 30};

// comprimir a [0, C-1] para que el wavelet tree trabaje con
// valores pequenos y contiguos.
vector<int> ordenado = original;
sort(ordenado.begin(), ordenado.end());
ordenado.erase(unique(ordenado.begin(), ordenado.end()), ordenado.end());

vector<int> comprimido(original.size());
for (int i = 0; i < (int)original.size(); i++)
    comprimido[i] = lower_bound(ordenado.begin(), ordenado.end(),
                                  original[i]) - ordenado.begin();
// comprimido: [3, 0, 2, 0, 1, 3, 2]   (valores 10,20,30,40 -> 0,1,2,3)

WaveletTree wt(comprimido.data(),
                comprimido.data() + comprimido.size(),
                0, (int)ordenado.size() - 1);

// 3er elemento mas chico en el rango [2, 6] (1-indexado)
int idx = wt.kth(2, 6, 3);
cout << ordenado[idx];
// posiciones 2..6 (1-idx): 10,30,10,20,40,30 (valores 0-idx 2..6)
// -> valores reales: 30,10,20,40,30 ordenados: 10,20,30,30,40
// 3er mas chico = 30

cout << wt.contarMenorIgual(1, 7, 1);
// cuantos valores comprimidos <= 1 (o sea, <= 20) hay en todo el arreglo
// original: 40,10,30,10,20,40,30 -> <=20 son: 10,10,20 -> 3

cout << wt.contarIgual(1, 7, 0);
// cuantas veces aparece el valor comprimido 0 (o sea, 10) -> 2
\end{lstlisting}

\textbf{Comparación con las otras estructuras de rango:}

\begin{table}[H]
\raggedright
\scriptsize
\setlength{\tabcolsep}{3pt}
\begin{tabular}{|p{3.3cm}|p{3.3cm}|}
\hline
\textbf{Pregunta} & \textbf{Estructura} \\
\hline
Suma/min/max de rango & Fenwick / Segment Tree / Sparse Table \\
\hline
$k$-ésimo menor, o conteo por rango de \textbf{valores}, en un rango de \textbf{posiciones} & Wavelet Tree \\
\hline
Elementos distintos, moda, funciones no lineales de frecuencia & Mo's Algorithm \\
\hline
\end{tabular}
\caption{Cuándo el problema necesita Wavelet Tree en vez de las otras estructuras.}
\end{table}

\textbf{Regla práctica:} si el enunciado combina una condición sobre
\textbf{valores} (``menor que'', ``entre $a$ y $b$'', ``el $k$-ésimo
más chico'') con un rango de \textbf{posiciones}, y el arreglo es
estático, es casi siempre Wavelet Tree. Si solo te piden agregar
(suma/min/max) sin condición sobre el valor, usa las estructuras más
simples de las secciones anteriores.


\subsection{Heavy-Light Decomposition}

\textbf{Idea:} descompone un árbol en \textbf{cadenas} de forma que
cualquier camino desde la raíz hasta cualquier nodo cruza a lo más
$O(\log n)$ cadenas distintas. En cada nodo, el hijo con el
\textbf{subárbol más grande} (el hijo ``pesado'') continúa la misma
cadena que su padre; los demás hijos inician cadenas nuevas. Al
numerar los nodos en el orden en que se visitan (DFS que siempre
sigue primero al hijo pesado), cada cadena queda como un \textbf{rango
contiguo} de posiciones — así que una query sobre un \textbf{camino}
del árbol se convierte en $O(\log n)$ queries de rango sobre un
Segment Tree normal, una por cada cadena que el camino atraviesa.

\textbf{¿Por qué usarlo?} Es la forma estándar de llevar range
queries de un arreglo a un \textbf{árbol}: sumas/mínimos/updates en
el camino entre dos nodos, o en el subárbol de un nodo. Sin HLD, un
camino arbitrario no tiene ninguna estructura contigua sobre la cual
aplicar Fenwick o Segment Tree directamente.

\textbf{Casos de uso:}

\begin{itemize}
    \item Suma/mínimo/máximo en el camino entre dos nodos $u$ y $v$.
    \item Actualizar todos los nodos (o aristas) de un camino $u$-$v$.
    \item Consultas o updates sobre el \textbf{subárbol} completo de un
    nodo (más simple: el subárbol siempre es un rango contiguo, sin
    necesidad de descomponer en cadenas).
    \item Cualquier problema de árbol que ``se sentiría'' como Range
    Queries si el árbol fuera un arreglo lineal.
\end{itemize}

\textbf{Complejidad:} construir $O(n)$ (dos DFS) más $O(n)$ del
Segment Tree subyacente; query o update de \textbf{camino}
$O(\log^2 n)$ ($O(\log n)$ cadenas $\times$ $O(\log n)$ por consulta
al Segment Tree); query o update de \textbf{subárbol} $O(\log n)$
(un solo rango contiguo).


\subsubsection{HLD: suma en camino, suma en subárbol}

\textbf{Idea de la implementación:} dos DFS \textbf{iterativas} (con
pila explícita, para evitar stack overflow en árboles muy profundos
o muy desbalanceados — un riesgo real con recursión cuando
$n\ge10^5$ y el árbol es una cadena larga). El primero
(\texttt{dfsTamano}) calcula el tamaño de cada subárbol y decide el
hijo pesado de cada nodo. El segundo (\texttt{dfsDescomponer})
recorre siempre primero al hijo pesado (así la cadena sigue
consecutiva en las posiciones) y solo después explora los hijos
livianos, cada uno iniciando su propia cadena nueva. Las posiciones
resultantes (\texttt{pos[]}) alimentan un Segment Tree normal —aquí
se reutiliza \texttt{SegmentTreeLazyAdd} de la sección anterior—
sobre el cual se hacen las queries reales.

\begin{lstlisting}[style=cpp]
struct HLD {
    vector<vector<int>> adj;
    vector<int> padre, profundidad, tamano, pesado, cabeza, pos;
    int n, timer = 0;
    SegmentTreeLazyAdd* seg;   // Segment Tree sobre el orden pos[]

    // CREATE: dos DFS iterativas para calcular cadenas, luego arma
    // el Segment Tree sobre el arreglo reordenado por pos[].
    HLD(vector<vector<int>>& adjacency, vector<long long>& valores,
        int raiz) {
        adj = adjacency;
        n = adj.size();
        padre.assign(n, -1);
        profundidad.assign(n, 0);
        tamano.assign(n, 1);
        pesado.assign(n, -1);
        cabeza.assign(n, raiz);
        pos.assign(n, 0);

        dfsTamano(raiz);
        dfsDescomponer(raiz);

        vector<long long> reordenado(n);
        for (int v = 0; v < n; v++)
            reordenado[pos[v]] = valores[v];

        seg = new SegmentTreeLazyAdd(reordenado);
    }                                          // O(n)

    // Fase 1: obtiene un orden donde cada padre aparece ANTES que
    // sus hijos (equivalente a un preorden), usando una pila.
    // Fase 2: recorre ese orden al reves (hijos antes que padres)
    // para acumular tamano[] y decidir el hijo pesado, igual que
    // haria la version recursiva al "volver" de cada llamada.
    void dfsTamano(int raiz) {
        vector<int> orden;
        stack<int> st;
        st.push(raiz);
        padre[raiz] = -1;
        profundidad[raiz] = 0;

        while (!st.empty()) {
            int u = st.top();
            st.pop();
            orden.push_back(u);

            for (int v : adj[u]) {
                if (v != padre[u]) {
                    padre[v] = u;
                    profundidad[v] = profundidad[u] + 1;
                    st.push(v);
                }
            }
        }

        for (int i = (int)orden.size() - 1; i >= 0; i--) {
            int u = orden[i];
            tamano[u] = 1;
            int maxSub = 0;

            for (int v : adj[u]) {
                if (v == padre[u]) continue;

                tamano[u] += tamano[v];
                if (tamano[v] > maxSub) {
                    maxSub = tamano[v];
                    pesado[u] = v;
                }
            }
        }
    }                                          // O(n)

    // Pila de pares (nodo, cabeza de su cadena). Para que el hijo
    // pesado quede en el tope de la pila (y por lo tanto se
    // procese inmediatamente despues, como en la version
    // recursiva), se empuja de ULTIMO: los hijos livianos
    // (cadena nueva) se empujan primero.
    void dfsDescomponer(int raiz) {
        stack<pair<int, int>> st;
        st.push({raiz, raiz});

        while (!st.empty()) {
            auto [u, h] = st.top();
            st.pop();

            cabeza[u] = h;
            pos[u] = timer++;

            for (int v : adj[u]) {
                if (v == padre[u] || v == pesado[u]) continue;
                st.push({v, v});                // cadena nueva
            }

            if (pesado[u] != -1)
                st.push({pesado[u], h});        // misma cadena
        }
    }                                          // O(n)

    // READ: suma de todos los nodos en el camino u-v.
    long long consultarCamino(int u, int v) {
        long long resultado = 0;

        while (cabeza[u] != cabeza[v]) {
            if (profundidad[cabeza[u]] < profundidad[cabeza[v]])
                swap(u, v);

            resultado += seg->rango(pos[cabeza[u]], pos[u]);
            u = padre[cabeza[u]];
        }

        if (profundidad[u] > profundidad[v]) swap(u, v);
        resultado += seg->rango(pos[u], pos[v]);
        return resultado;
    }                                          // O(log^2 n)

    // UPDATE: suma val a todos los nodos en el camino u-v.
    void actualizarCamino(int u, int v, long long val) {
        while (cabeza[u] != cabeza[v]) {
            if (profundidad[cabeza[u]] < profundidad[cabeza[v]])
                swap(u, v);

            seg->rangoUpdate(pos[cabeza[u]], pos[u], val);
            u = padre[cabeza[u]];
        }

        if (profundidad[u] > profundidad[v]) swap(u, v);
        seg->rangoUpdate(pos[u], pos[v], val);
    }                                          // O(log^2 n)

    // READ: suma de todo el subarbol de u (rango contiguo por
    // construccion, gracias al orden DFS de dfsDescomponer).
    long long consultarSubarbol(int u) {
        return seg->rango(pos[u], pos[u] + tamano[u] - 1);
    }                                          // O(log n)

    // UPDATE: suma val a todo el subarbol de u.
    void actualizarSubarbol(int u, long long val) {
        seg->rangoUpdate(pos[u], pos[u] + tamano[u] - 1, val);
    }                                          // O(log n)
};
\end{lstlisting}

\textbf{Ejemplo de funcionamiento:}

\begin{lstlisting}[style=cpp]
// arbol:        0
//              /|\
//             1 2 3
//            /|
//           4 5

vector<vector<int>> adj(6);
adj[0] = {1, 2, 3};
adj[1] = {0, 4, 5};
adj[2] = {0};
adj[3] = {0};
adj[4] = {1};
adj[5] = {1};

vector<long long> valores = {1, 2, 3, 4, 5, 6};
// nodo:                     0  1  2  3  4  5

HLD hld(adj, valores, 0);

cout << hld.consultarCamino(4, 3);
// camino 4 -> 1 -> 0 -> 3 : valores 5+2+1+4 = 12

cout << hld.consultarSubarbol(1);
// subarbol de 1: nodos {1,4,5} : 2+5+6 = 13

hld.actualizarCamino(4, 5, 10);
// suma 10 a cada nodo del camino 4 -> 1 -> 5

cout << hld.consultarSubarbol(0);
// subarbol de 0 (todo el arbol) tras el update:
// 1+(2+10)+3+4+(5+10)+(6+10) = 51
\end{lstlisting}

\textbf{Nota:} al ser iterativas, \texttt{dfsTamano} y
\texttt{dfsDescomponer} no tienen riesgo de stack overflow incluso en
árboles con $n\ge10^5$ y forma de ``cadena larga'' (el peor caso para
recursión). \texttt{SegmentTreeLazyAdd} es exactamente el struct de
la sección de Segment Tree — cualquier otra variante (min, max,
assign) funciona igual de bien como motor subyacente.

\textbf{Regla práctica:} usa HLD en cuanto el problema pida algo
sobre un \textbf{camino} entre dos nodos de un árbol con muchas
queries. Si solo pide algo sobre \textbf{subárboles} (sin caminos
arbitrarios), no hace falta la descomposición en cadenas completa:
basta un solo DFS para las posiciones (\texttt{pos[]} y
\texttt{tamano[]}) y un Segment Tree normal, ya que el subárbol
siempre es un rango contiguo por sí solo.



\part{Grafos}

\subsection{¿Qué algoritmo usar?}

La siguiente tabla resume situaciones típicas de Competitive Programming
y el algoritmo que debe considerarse. La idea es reconocer las palabras
clave o condiciones del enunciado y asociarlas rápidamente con una técnica.

\begin{table}[H]
\raggedright
\scriptsize
\setlength{\tabcolsep}{2.5pt}
\renewcommand{\arraystretch}{1.15}
\begin{tabular}{|p{0.42\columnwidth}|p{0.23\columnwidth}|p{0.25\columnwidth}|}
\hline
\textbf{Cuando el ejercicio diga / pida} &
\textbf{Usa} &
\textbf{Pista rápida} \\
\hline

``Recorrer el grafo por niveles'' &
BFS &
Usa una cola. \\
\hline

``Encontrar la distancia mínima'' con aristas de igual costo &
BFS &
Cada arista cuesta lo mismo. \\
\hline

``Encontrar caminos mínimos'' con pesos $0$ y $1$ &
0-1 BFS &
Deque; pesos únicamente $0$ o $1$. \\
\hline

``Camino mínimo'' con pesos no negativos &
Dijkstra &
Usa una cola de prioridad. \\
\hline

``Camino mínimo'' con pesos negativos &
Bellman-Ford &
Permite detectar ciclos negativos. \\
\hline

``Camino mínimo entre todos los pares'' y $n$ pequeño &
Floyd-Warshall &
DP sobre todos los pares. \\
\hline

``¿Existe un camino entre $u$ y $v$?'' &
DFS / BFS &
Problema de conectividad. \\
\hline

``Encontrar componentes conexas'' &
DFS / BFS &
Ejecutar traversal desde cada nodo no visitado. \\
\hline

``Detectar si existe un ciclo'' en grafo no dirigido &
DFS / DSU &
DFS para recorrer; DSU si se agregan aristas. \\
\hline

``Detectar ciclos'' en grafo dirigido &
DFS / Kahn &
Usar estados en DFS o grados de entrada. \\
\hline

``Ordenar tareas con dependencias'' &
Topological Sort &
El grafo debe ser dirigido y acíclico. \\
\hline

``Encontrar puentes'' &
Tarjan &
DFS + tiempos de descubrimiento. \\
\hline

``Encontrar puntos de articulación'' &
Tarjan &
Vértices cuya eliminación desconecta el grafo. \\
\hline

``Encontrar componentes fuertemente conexas'' &
Kosaraju / Tarjan &
Grafo dirigido. \\
\hline

``Conectar todos los nodos con costo mínimo'' &
MST &
Kruskal o Prim. \\
\hline

``MST procesando aristas ordenadas por peso'' &
Kruskal &
Ordenar aristas + DSU. \\
\hline

``MST expandiendo desde un nodo'' &
Prim &
Priority Queue + lista de adyacencia. \\
\hline

``¿Dos nodos pertenecen al mismo conjunto?'' &
DSU &
Union-Find. \\
\hline

``Agrupar elementos mediante uniones'' &
DSU &
\texttt{union} + \texttt{find}. \\
\hline

``Encontrar si un grafo es bipartito'' &
2-Coloring &
Colorear con dos colores. \\
\hline

``Máximo número de parejas entre dos grupos'' &
Bipartite Matching &
Emparejamiento entre conjuntos. \\
\hline

``Máximo flujo de una fuente a un destino'' &
Dinic &
Network Flow. \\
\hline

``Separar fuente y sumidero con capacidad mínima'' &
Min-Cut &
Por el teorema Max-Flow Min-Cut. \\
\hline

``Distancia entre dos nodos de un árbol'' &
LCA &
$\operatorname{dist}(u,v)=
d[u]+d[v]-2d[\operatorname{LCA}(u,v)]$. \\
\hline

``Encontrar el ancestro $k$ niveles arriba'' &
Binary Lifting &
Saltar potencias de $2$. \\
\hline

``Consultas sobre subárboles'' &
Euler Tour + Fenwick / Segment Tree &
Convertir subárbol en intervalo. \\
\hline

``Consultas o actualizaciones sobre caminos de un árbol'' &
HLD &
Descomponer el camino en cadenas. \\
\hline

``Encontrar el diámetro de un árbol'' &
BFS / DFS &
Dos traversals desde extremos. \\
\hline

``Encontrar un camino que use todas las aristas'' &
Eulerian Path / Hierholzer &
Trabaja con aristas, no con vértices. \\
\hline

``Encontrar todas las rutas posibles'' &
DFS + Backtracking &
Explorar y deshacer decisiones. \\
\hline

``Rellenar una región conectada'' &
Flood Fill &
DFS / BFS desde una celda inicial. \\
\hline

``Moverse por una cuadrícula buscando la menor cantidad de pasos'' &
BFS &
Cada movimiento tiene costo $1$. \\
\hline

``Variables booleanas con restricciones de la forma OR'' &
2-SAT + SCC &
Construir grafo de implicaciones. \\
\hline

``Resolver muchas consultas sobre un árbol'' &
LCA / HLD &
Depende de si son ancestros o caminos. \\
\hline

\end{tabular}
\caption{Guía rápida para seleccionar algoritmos de grafos.}
\end{table}
\begin{table}[H]
\raggedright
\scriptsize
\setlength{\tabcolsep}{3pt}
\renewcommand{\arraystretch}{1.15}
\begin{tabular}{|p{0.43\columnwidth}|p{0.27\columnwidth}|p{0.22\columnwidth}|}
\hline
\textbf{Condición del problema} &
\textbf{Algoritmo} &
\textbf{Complejidad} \\
\hline
Aristas sin peso / costo $1$ &
BFS &
$O(n+m)$ \\
\hline
Pesos únicamente $0$ y $1$ &
0-1 BFS &
$O(n+m)$ \\
\hline
Pesos no negativos &
Dijkstra &
$O((n+m)\log n)$ \\
\hline
Puede haber pesos negativos &
Bellman-Ford &
$O(nm)$ \\
\hline
Todos los pares y $n$ pequeño &
Floyd-Warshall &
$O(n^3)$ \\
\hline
\end{tabular}
\caption{Selección del algoritmo para caminos mínimos.}
\end{table}



% =========================================================
% TRAVERSAL
% =========================================================

\section{Traversal}


\subsection{Representación de Grafos}

\textbf{Idea:} antes de aplicar cualquier algoritmo sobre un grafo,
primero debemos decidir cómo almacenar sus vértices y aristas. Las
tres representaciones más comunes son la \textbf{lista de adyacencia},
la \textbf{matriz de adyacencia} y la \textbf{lista de aristas}. La
elección depende principalmente de la cantidad de vértices y aristas
y del tipo de consultas que realizará el algoritmo.

La \textbf{lista de adyacencia} almacena, para cada vértice, todos los
vértices conectados directamente con él. Es la representación más
utilizada en Competitive Programming porque ocupa $O(n+m)$ memoria y
permite recorrer eficientemente todos los vecinos de un vértice.

La \textbf{matriz de adyacencia} utiliza una matriz $n\times n$ donde
la posición $(u,v)$ indica si existe una arista entre $u$ y $v$.
Permite comprobar la existencia de una arista en $O(1)$, pero requiere
$O(n^2)$ memoria.

La \textbf{lista de aristas} simplemente almacena cada arista como un
par $(u,v)$, o como una terna $(u,v,w)$ si tiene peso. Es especialmente
útil para algoritmos que necesitan procesar todas las aristas, como
Kruskal o Bellman-Ford.

\textbf{¿Por qué usar cada una?}

\begin{itemize}
    \item \textbf{Lista de adyacencia:} opción general para BFS, DFS,
    Dijkstra, Prim, búsqueda de ciclos y prácticamente la mayoría de
    algoritmos sobre grafos dispersos.

    \item \textbf{Matriz de adyacencia:} útil cuando $n$ es pequeño o
    cuando se necesita comprobar constantemente si existe una arista
    entre dos vértices en $O(1)$.

    \item \textbf{Lista de aristas:} conveniente cuando el algoritmo
    necesita recorrer todas las aristas directamente, especialmente
    en Kruskal y Bellman-Ford.
\end{itemize}

\textbf{Tipos de grafos:} una arista puede ser \textbf{dirigida} o
\textbf{no dirigida}, y puede tener o no un \textbf{peso}. En un grafo
no dirigido, una arista $(u,v)$ se representa agregando $v$ a los
vecinos de $u$ y $u$ a los vecinos de $v$. En un grafo dirigido,
solamente se agrega $v$ a los vecinos de $u$.

\textbf{Complejidad:}

\begin{table}[H]
\raggedright
\scriptsize
\setlength{\tabcolsep}{2.5pt}
\renewcommand{\arraystretch}{1.1}
\begin{tabular}{|p{0.27\columnwidth}|c|c|c|}
\hline
\textbf{Representación} &
\textbf{Memoria} &
\textbf{Consulta} &
\textbf{Vecinos} \\
\hline
Lista de adyacencia &
$O(n+m)$ &
$O(\deg(u))$ &
$O(\deg(u))$ \\
\hline
Matriz de adyacencia &
$O(n^2)$ &
$O(1)$ &
$O(n)$ \\
\hline
Lista de aristas &
$O(m)$ &
$O(m)$ &
$O(m)$ \\
\hline
\end{tabular}
\caption{Comparación de representaciones de grafos.}
\end{table}

\subsubsection{Lista de adyacencia}

\textbf{Idea de la implementación:} se utiliza un
\texttt{vector<vector<int>>} donde la posición $u$ contiene todos los
vértices conectados directamente con $u$. Para grafos ponderados se
puede almacenar un par $(v,w)$, donde $v$ es el destino y $w$ el peso
de la arista.

\begin{lstlisting}[style=cpp]
// Grafo no dirigido y no ponderado.
int n = 5;
vector<vector<int>> adj(n);

// Agrega una arista u-v.
auto agregarArista = [&](int u, int v) {
    adj[u].push_back(v);
    adj[v].push_back(u);
};

agregarArista(0, 1);
agregarArista(0, 2);
agregarArista(1, 3);
agregarArista(1, 4);
\end{lstlisting}

\textbf{Grafo resultante:}

\begin{lstlisting}[style=cpp]
//        0
//       / \
//      1   2
//     / \
//    3   4
//
// adj[0] = {1, 2}
// adj[1] = {0, 3, 4}
// adj[2] = {0}
// adj[3] = {1}
// adj[4] = {1}
\end{lstlisting}

Para un grafo dirigido, solamente se agrega la arista en una dirección:

\begin{lstlisting}[style=cpp]
// Grafo dirigido.
vector<vector<int>> adj(n);

auto agregarArista = [&](int u, int v) {
    adj[u].push_back(v);
};

agregarArista(0, 1);
agregarArista(0, 2);
agregarArista(1, 3);
\end{lstlisting}

Para un grafo ponderado, cada vecino puede almacenar también el peso:

\begin{lstlisting}[style=cpp]
// Grafo ponderado.
vector<vector<pair<int, long long>>> adj(n);

auto agregarArista = [&](int u, int v, long long w) {
    adj[u].push_back({v, w});
    adj[v].push_back({u, w});
};

agregarArista(0, 1, 5);
agregarArista(0, 2, 3);
agregarArista(1, 3, 7);
\end{lstlisting}

\textbf{Ejemplo de recorrido:}

\begin{lstlisting}[style=cpp]
void dfs(int u, int padre, vector<vector<int>>& adj) {

    cout << u << ' ';

    for (int v : adj[u]) {
        if (v == padre) continue;

        dfs(v, u, adj);
    }
}
\end{lstlisting}

La lista de adyacencia permite que DFS o BFS procesen cada vértice y
cada arista una cantidad constante de veces, obteniendo una
complejidad total de $O(n+m)$.

\subsubsection{Matriz de adyacencia}

\textbf{Idea:} una matriz $adj[u][v]$ indica directamente si existe
una conexión desde $u$ hacia $v$. Si el grafo es ponderado, la
posición puede almacenar directamente el peso de la arista.

\begin{lstlisting}[style=cpp]
int n = 5;
vector<vector<int>> adj(n, vector<int>(n, 0));

auto agregarArista = [&](int u, int v) {
    adj[u][v] = 1;
    adj[v][u] = 1;
};

agregarArista(0, 1);
agregarArista(0, 2);
agregarArista(1, 3);
agregarArista(1, 4);
\end{lstlisting}

Ahora comprobar si existe una arista entre $u$ y $v$ es inmediato:

\begin{lstlisting}[style=cpp]
if (adj[u][v]) {
    cout << "Existe una arista";
}
\end{lstlisting}

La consulta anterior cuesta $O(1)$, pero recorrer todos los vecinos de
$u$ requiere revisar las $n$ posiciones de la fila:

\begin{lstlisting}[style=cpp]
for (int v = 0; v < n; v++) {
    if (adj[u][v]) {
        cout << v << ' ';
    }
}
\end{lstlisting}

\subsubsection{Lista de aristas}

\textbf{Idea:} almacena directamente todas las aristas del grafo.
Para un grafo no ponderado basta guardar $(u,v)$; para uno ponderado
se guarda $(u,v,w)$.

\begin{lstlisting}[style=cpp]
// Grafo ponderado.
struct Edge {
    int u, v;
    long long w;
};

vector<Edge> edges;

edges.push_back({0, 1, 5});
edges.push_back({0, 2, 3});
edges.push_back({1, 3, 7});
\end{lstlisting}

Esta representación resulta especialmente conveniente cuando el
algoritmo necesita procesar todas las aristas:

\begin{lstlisting}[style=cpp]
for (Edge e : edges) {
    cout << e.u << " -> " << e.v
         << " peso = " << e.w << '\n';
}
\end{lstlisting}

Por ejemplo, \textbf{Kruskal} ordena directamente esta lista por peso,
mientras que \textbf{Bellman-Ford} puede recorrerla completa en cada
iteración.

\textbf{Nota:} para la mayoría de problemas de Competitive Programming,
la \textbf{lista de adyacencia} debe ser la representación por defecto.
Es especialmente importante cuando el grafo es disperso, es decir,
cuando $m$ es mucho menor que $n^2$. La matriz de adyacencia solamente
conviene cuando $n$ es pequeño o cuando las consultas de existencia de
aristas en $O(1)$ son importantes.

\textbf{Regla práctica:} si no tienes una razón específica para usar
otra representación, usa \textbf{lista de adyacencia}. Si necesitas
preguntar constantemente ``¿existe una arista $u\rightarrow v$?'',
considera una \textbf{matriz de adyacencia}. Si el algoritmo trabaja
directamente con todas las aristas, considera una \textbf{lista de
aristas}.

```latex
\subsection{BFS (Breadth-First Search)}

\textbf{Idea:} BFS (\textit{Breadth-First Search}) es un algoritmo de
recorrido de grafos que explora los vértices por \textbf{niveles}. A
partir de un vértice inicial, primero visita todos sus vecinos
directos, después los vecinos de estos, y continúa de esta manera
hasta procesar todos los vértices alcanzables.

Para mantener este orden se utiliza una \textbf{cola}
(\texttt{queue}), siguiendo el principio FIFO (\textit{First In,
First Out}). Cada vértice se marca como visitado cuando se agrega a
la cola, evitando procesarlo varias veces.

Una de las propiedades más importantes de BFS es que, en un grafo
\textbf{no ponderado}, encuentra la distancia mínima desde un vértice
origen hasta cualquier otro vértice alcanzable, donde la distancia se
mide como el número de aristas recorridas.

\textbf{¿Por qué usarlo?}

\begin{itemize}
    \item \textbf{Camino mínimo:} encuentra el camino con menor número
    de aristas en grafos no ponderados.

    \item \textbf{Distancias:} permite calcular la distancia mínima
    desde un origen hasta todos los vértices alcanzables.

    \item \textbf{Recorrido por niveles:} permite procesar los
    vértices según su distancia al origen.

    \item \textbf{Reconstrucción de caminos:} almacenando el padre de
    cada vértice se puede recuperar el camino mínimo.

    \item \textbf{Bipartición:} permite determinar si un grafo es
    bipartito utilizando dos colores.

    \item \textbf{Componentes conexas:} ejecutando BFS desde cada
    vértice no visitado se pueden encontrar todas las componentes
    conexas.

    \item \textbf{Problemas en grids:} permite encontrar la cantidad
    mínima de movimientos cuando cada movimiento tiene el mismo costo.

    \item \textbf{Multi-Source BFS:} permite calcular la distancia
    mínima desde cualquiera de varios vértices iniciales.
\end{itemize}

\textbf{Funcionamiento:}

El algoritmo sigue los siguientes pasos:

\begin{enumerate}
    \item Elegir un vértice inicial $s$.
    \item Marcar $s$ como visitado y agregarlo a la cola.
    \item Mientras la cola no esté vacía, extraer el primer vértice.
    \item Recorrer todos los vecinos del vértice extraído.
    \item Si un vecino no ha sido visitado, marcarlo y agregarlo a
    la cola.
\end{enumerate}

La cola garantiza que los vértices se procesen por niveles. Primero
se procesan los vértices a distancia $0$, después los de distancia
$1$, luego los de distancia $2$, y así sucesivamente.

\textbf{Implementación básica:}

\begin{lstlisting}[style=cpp]
void bfs(int inicio, vector<vector<int>>& adj) {

    int n = adj.size();

    vector<bool> visitado(n, false);
    queue<int> q;

    visitado[inicio] = true;
    q.push(inicio);

    while (!q.empty()) {

        int u = q.front();
        q.pop();

        cout << u << ' ';

        for (int v : adj[u]) {

            if (visitado[v]) continue;

            visitado[v] = true;
            q.push(v);
        }
    }
}
\end{lstlisting}

\textbf{Ejemplo:}

Consideremos el siguiente grafo:

\begin{lstlisting}[style=cpp]
//        0
//       / \
//      1   2
//     / \
//    3   4
\end{lstlisting}

Si comenzamos BFS desde el vértice $0$, los niveles son:

\begin{lstlisting}[style=cpp]
// Nivel 0: 0
// Nivel 1: 1 2
// Nivel 2: 3 4
\end{lstlisting}

Por lo tanto, un posible orden de recorrido es:

\begin{lstlisting}[style=cpp]
0 1 2 3 4
\end{lstlisting}

El orden entre vértices del mismo nivel depende del orden en que
aparezcan en la lista de adyacencia.

\textbf{BFS para encontrar distancias:}

Para calcular la distancia mínima desde el origen hasta cada vértice,
se utiliza un arreglo \texttt{dist}. Inicialmente todos los valores
son $-1$, indicando que el vértice todavía no ha sido visitado.

Cuando BFS descubre un nuevo vértice $v$ desde $u$, su distancia se
calcula como:

\[
dist[v] = dist[u] + 1
\]

\begin{lstlisting}[style=cpp]
vector<int> bfsDist(int inicio, vector<vector<int>>& adj) {

    int n = adj.size();

    vector<int> dist(n, -1);
    queue<int> q;

    dist[inicio] = 0;
    q.push(inicio);

    while (!q.empty()) {

        int u = q.front();
        q.pop();

        for (int v : adj[u]) {

            if (dist[v] != -1) continue;

            dist[v] = dist[u] + 1;
            q.push(v);
        }
    }

    return dist;
}
\end{lstlisting}

Para el ejemplo anterior, comenzando desde $0$, se obtiene:

\begin{lstlisting}[style=cpp]
// dist[0] = 0
// dist[1] = 1
// dist[2] = 1
// dist[3] = 2
// dist[4] = 2
\end{lstlisting}

Por ejemplo, $dist[3] = 2$ significa que el camino mínimo desde
$0$ hasta $3$ utiliza dos aristas.

\textbf{Reconstrucción del camino:}

Si además de la distancia se necesita conocer el camino utilizado,
se puede almacenar el \textbf{padre} de cada vértice. Cuando BFS
descubre $v$ desde $u$, se establece:

\[
padre[v] = u
\]

\begin{lstlisting}[style=cpp]
vector<int> dist(n, -1);
vector<int> padre(n, -1);

queue<int> q;

dist[inicio] = 0;
q.push(inicio);

while (!q.empty()) {

    int u = q.front();
    q.pop();

    for (int v : adj[u]) {

        if (dist[v] != -1) continue;

        dist[v] = dist[u] + 1;
        padre[v] = u;

        q.push(v);
    }
}
\end{lstlisting}

Una vez terminado BFS, el camino hasta un destino se reconstruye
siguiendo los padres desde el destino hasta el origen:

\begin{lstlisting}[style=cpp]
vector<int> camino;

for (int v = destino; v != -1; v = padre[v]) {
    camino.push_back(v);
}

reverse(camino.begin(), camino.end());
\end{lstlisting}

Por ejemplo, si:

\begin{lstlisting}[style=cpp]
// padre[3] = 1
// padre[1] = 0
// padre[0] = -1
\end{lstlisting}

el camino mínimo desde $0$ hasta $3$ es:

\begin{lstlisting}[style=cpp]
0 -> 1 -> 3
\end{lstlisting}

\textbf{BFS en grafos no conectados:}

Un BFS iniciado desde un único vértice solamente visita los vértices
que pueden alcanzarse desde él. Para recorrer todo el grafo, se puede
ejecutar BFS desde cada vértice que todavía no haya sido visitado.

\begin{lstlisting}[style=cpp]
vector<bool> visitado(n, false);

for (int s = 0; s < n; s++) {

    if (visitado[s]) continue;

    queue<int> q;

    visitado[s] = true;
    q.push(s);

    while (!q.empty()) {

        int u = q.front();
        q.pop();

        for (int v : adj[u]) {

            if (visitado[v]) continue;

            visitado[v] = true;
            q.push(v);
        }
    }
}
\end{lstlisting}

Esta técnica permite encontrar, entre otras cosas, el número de
\textbf{componentes conexas} de un grafo.

\textbf{BFS para comprobar si un grafo es bipartito:}

Un grafo es bipartito si sus vértices pueden dividirse en dos grupos
de manera que ninguna arista conecte dos vértices del mismo grupo.
BFS puede comprobar esta propiedad asignando alternativamente los
colores $0$ y $1$ a los vértices.

Si un vértice $u$ tiene color $c$, sus vecinos deben tener el color
$c \mathbin{\hat{}} 1$. Si una arista conecta dos vértices del mismo
color, el grafo no es bipartito.

\begin{lstlisting}[style=cpp]
bool esBipartito(vector<vector<int>>& adj) {

    int n = adj.size();

    vector<int> color(n, -1);

    for (int inicio = 0; inicio < n; inicio++) {

        if (color[inicio] != -1) continue;

        queue<int> q;

        color[inicio] = 0;
        q.push(inicio);

        while (!q.empty()) {

            int u = q.front();
            q.pop();

            for (int v : adj[u]) {

                if (color[v] == -1) {

                    color[v] = color[u] ^ 1;
                    q.push(v);

                } else if (color[v] == color[u]) {

                    return false;
                }
            }
        }
    }

    return true;
}
\end{lstlisting}

El arreglo \texttt{color} también funciona como arreglo de visitados:
un valor de $-1$ significa que el vértice todavía no ha sido
procesado.

\textbf{BFS en un grid:}

BFS también puede utilizarse sobre matrices o tableros. Cada celda
se considera un vértice y cada movimiento válido entre celdas se
considera una arista.

Por ejemplo, si solamente se permiten movimientos hacia arriba,
abajo, izquierda y derecha:

\begin{lstlisting}[style=cpp]
int dx[] = {-1, 1, 0, 0};
int dy[] = {0, 0, -1, 1};

queue<pair<int, int>> q;

q.push({sx, sy});
dist[sx][sy] = 0;

while (!q.empty()) {

    auto [x, y] = q.front();
    q.pop();

    for (int d = 0; d < 4; d++) {

        int nx = x + dx[d];
        int ny = y + dy[d];

        if (nx < 0 || nx >= n ||
            ny < 0 || ny >= m)
            continue;

        if (dist[nx][ny] != -1)
            continue;

        if (grid[nx][ny] == '#')
            continue;

        dist[nx][ny] = dist[x][y] + 1;
        q.push({nx, ny});
    }
}
\end{lstlisting}

Esta técnica es especialmente útil en problemas donde se pregunta
por la menor cantidad de movimientos necesarios para llegar desde
una celda hasta otra.

\textbf{Multi-Source BFS:}

En algunos problemas existen varios vértices iniciales al mismo
tiempo. En lugar de ejecutar un BFS separado desde cada uno, todos
los orígenes se introducen inicialmente en la misma cola con
distancia $0$.

\begin{lstlisting}[style=cpp]
queue<int> q;
vector<int> dist(n, -1);

for (int s : fuentes) {

    dist[s] = 0;
    q.push(s);
}

while (!q.empty()) {

    int u = q.front();
    q.pop();

    for (int v : adj[u]) {

        if (dist[v] != -1) continue;

        dist[v] = dist[u] + 1;
        q.push(v);
    }
}
\end{lstlisting}

De esta manera, cada vértice obtiene la distancia mínima hasta
\textbf{cualquiera} de las fuentes.

\textbf{Complejidad:}

Utilizando una lista de adyacencia, cada vértice se procesa una sola
vez y cada arista se revisa una cantidad constante de veces.

\begin{table}[H]
\raggedright
\scriptsize
\setlength{\tabcolsep}{3pt}
\renewcommand{\arraystretch}{1.1}
\begin{tabular}{|l|c|c|}
\hline
\textbf{Representación} &
\textbf{Tiempo} &
\textbf{Memoria adicional} \\
\hline
Lista de adyacencia &
$O(n+m)$ &
$O(n)$ \\
\hline
Matriz de adyacencia &
$O(n^2)$ &
$O(n)$ \\
\hline
\end{tabular}
\caption{Complejidad de BFS según la representación del grafo.}
\end{table}

La complejidad $O(n+m)$ hace que BFS sea especialmente eficiente
para grafos dispersos representados mediante listas de adyacencia.

\textbf{Nota:} BFS encuentra caminos mínimos cuando todas las aristas
tienen el mismo costo. No debe utilizarse directamente para encontrar
caminos mínimos en grafos con pesos arbitrarios. Para pesos no
negativos se puede utilizar \textbf{Dijkstra}; si los pesos solamente
son $0$ y $1$, una alternativa eficiente es \textbf{0-1 BFS}.

\textbf{Regla práctica:} si un problema pide minimizar el número de
\textbf{pasos, movimientos o aristas} y cada movimiento tiene el mismo
costo, considera primero \textbf{BFS}. Si necesitas conocer la
distancia, utiliza \texttt{dist}; si necesitas recuperar el camino,
utiliza \texttt{padre}; y si existen múltiples puntos de origen,
considera \textbf{Multi-Source BFS}.
```


```latex
\subsection{BFS en Grafos Desconectados}

\textbf{Idea:} un BFS iniciado desde un único vértice solamente visita
los vértices que son alcanzables desde él. Si el grafo tiene varias
\textbf{componentes conexas}, es necesario iniciar un nuevo BFS desde
cada vértice que todavía no haya sido visitado.

La idea consiste en recorrer todos los vértices y, cada vez que se
encuentra uno que no ha sido visitado, iniciar un nuevo BFS desde él.
Cada ejecución de BFS procesa una componente conexa completa.

\textbf{Implementación:}

\begin{lstlisting}[style=cpp]
void bfs(int inicio, vector<vector<int>>& adj,
         vector<bool>& visitado) {

    queue<int> q;

    visitado[inicio] = true;
    q.push(inicio);

    while (!q.empty()) {

        int u = q.front();
        q.pop();

        for (int v : adj[u]) {

            if (visitado[v]) continue;

            visitado[v] = true;
            q.push(v);
        }
    }
}
\end{lstlisting}

Para recorrer todo el grafo:

\begin{lstlisting}[style=cpp]
vector<bool> visitado(n, false);

for (int u = 0; u < n; u++) {

    if (visitado[u]) continue;

    bfs(u, adj, visitado);
}
\end{lstlisting}

Cada vez que se ejecuta \texttt{bfs}, se procesa una componente conexa
que todavía no había sido visitada.

\textbf{Contar componentes conexas:}

Una aplicación directa consiste en contar cuántas componentes conexas
tiene el grafo. Basta incrementar un contador cada vez que se inicia
un nuevo BFS.

\begin{lstlisting}[style=cpp]
int componentes = 0;

vector<bool> visitado(n, false);

for (int u = 0; u < n; u++) {

    if (visitado[u]) continue;

    componentes++;

    queue<int> q;

    visitado[u] = true;
    q.push(u);

    while (!q.empty()) {

        int v = q.front();
        q.pop();

        for (int w : adj[v]) {

            if (visitado[w]) continue;

            visitado[w] = true;
            q.push(w);
        }
    }
}

cout << componentes << '\n';
\end{lstlisting}

\textbf{Ejemplo:}

Consideremos el siguiente grafo:

\begin{lstlisting}[style=cpp]
//    0 --- 1       3 --- 4
//          |
//          2
//
//    5 --- 6
\end{lstlisting}

El grafo tiene tres componentes conexas:

\begin{lstlisting}[style=cpp]
// Componente 1: {0, 1, 2}
// Componente 2: {3, 4}
// Componente 3: {5, 6}
\end{lstlisting}

Si recorremos los vértices en orden, el primer BFS puede comenzar
desde $0$ y visitar $\{0,1,2\}$. Después, los vértices $1$ y $2$ ya
están visitados, por lo que el siguiente BFS comienza desde $3$ y
visita $\{3,4\}$. Finalmente, se inicia otro BFS desde $5$ para
visitar $\{5,6\}$.

\textbf{Complejidad:}

Aunque se ejecuten varios BFS, cada vértice se visita una sola vez y
cada arista se procesa una cantidad constante de veces. Por lo tanto,
la complejidad total utilizando una lista de adyacencia sigue siendo:

\[
O(n+m)
\]

La memoria adicional utilizada por el arreglo de visitados y la cola
es $O(n)$.

\begin{table}[H]
\raggedright
\scriptsize
\setlength{\tabcolsep}{3pt}
\renewcommand{\arraystretch}{1.1}
\begin{tabular}{|l|c|}
\hline
\textbf{Operación} & \textbf{Complejidad} \\
\hline
Recorrer todas las componentes & $O(n+m)$ \\
\hline
Contar componentes conexas & $O(n+m)$ \\
\hline
Memoria adicional & $O(n)$ \\
\hline
\end{tabular}
\caption{Complejidad de BFS en grafos desconectados.}
\end{table}

\textbf{Regla práctica:} si el grafo puede estar desconectado y se
necesita procesar \textbf{todos los vértices}, no basta con ejecutar
BFS una sola vez. Recorre todos los vértices y comienza un nuevo BFS
cada vez que encuentres uno que todavía no haya sido visitado.
```


\subsection{BFS para Todos los Caminos}

\textbf{Idea:} BFS no solamente permite encontrar la distancia mínima
desde un vértice origen hacia los demás vértices. También puede
utilizarse para reconstruir el \textbf{camino mínimo} desde el origen
hasta cualquier vértice alcanzable.

Para reconstruir un camino, además de almacenar la distancia, se
guarda para cada vértice su \textbf{padre}, es decir, el vértice desde
el cual fue descubierto por primera vez. Como BFS visita los vértices
por niveles, el primer padre asignado a un vértice pertenece a un
camino mínimo.

\textbf{Implementación:}

\begin{lstlisting}[style=cpp]
vector<int> dist(n, -1);
vector<int> padre(n, -1);

queue<int> q;

dist[inicio] = 0;
q.push(inicio);

while (!q.empty()) {

    int u = q.front();
    q.pop();

    for (int v : adj[u]) {

        if (dist[v] != -1) continue;

        dist[v] = dist[u] + 1;
        padre[v] = u;

        q.push(v);
    }
}
\end{lstlisting}

Después de ejecutar BFS, \texttt{dist[v]} contiene la distancia mínima
desde \texttt{inicio} hasta $v$, mientras que \texttt{padre[v]} indica
qué vértice se debe seguir para reconstruir el camino.

\textbf{Reconstrucción del camino:}

Para obtener el camino hasta un vértice destino, se comienza en
\texttt{destino} y se siguen los padres hasta llegar al origen.

\begin{lstlisting}[style=cpp]
vector<int> camino;

for (int v = destino; v != -1; v = padre[v]) {
    camino.push_back(v);
}

reverse(camino.begin(), camino.end());
\end{lstlisting}

El vector \texttt{camino} contiene ahora los vértices en el orden
correcto, desde el origen hasta el destino.

\textbf{Ejemplo:}

Supongamos que BFS produce los siguientes padres:

\begin{lstlisting}[style=cpp]
// padre[0] = -1
// padre[1] = 0
// padre[2] = 0
// padre[3] = 1
// padre[4] = 3
\end{lstlisting}

Si el origen es $0$ y queremos reconstruir el camino hasta $4$,
seguimos los padres:

\[
4 \rightarrow 3 \rightarrow 1 \rightarrow 0
\]

Después de invertirlo obtenemos:

\[
0 \rightarrow 1 \rightarrow 3 \rightarrow 4
\]

Por lo tanto, este es un camino mínimo desde $0$ hasta $4$.

\textbf{Función para reconstruir el camino:}

La reconstrucción puede encapsularse en una función:

\begin{lstlisting}[style=cpp]
vector<int> reconstruirCamino(int destino,
                              vector<int>& padre) {

    vector<int> camino;

    for (int v = destino; v != -1; v = padre[v]) {
        camino.push_back(v);
    }

    reverse(camino.begin(), camino.end());

    return camino;
}
\end{lstlisting}

El BFS completo puede combinarse con esta función:

\begin{lstlisting}[style=cpp]
vector<int> dist(n, -1);
vector<int> padre(n, -1);

queue<int> q;

dist[inicio] = 0;
q.push(inicio);

while (!q.empty()) {

    int u = q.front();
    q.pop();

    for (int v : adj[u]) {

        if (dist[v] != -1) continue;

        dist[v] = dist[u] + 1;
        padre[v] = u;

        q.push(v);
    }
}

vector<int> camino =
    reconstruirCamino(destino, padre);
\end{lstlisting}

\textbf{¿Qué significa ``todos los caminos''?}

En este contexto, BFS obtiene un camino mínimo desde el origen hacia
\textbf{cada vértice alcanzable}. No significa que encuentre todos los
caminos posibles entre dos vértices, ya que la cantidad de caminos
posibles puede ser exponencial.

Por ejemplo, después de ejecutar BFS desde $s$:

\begin{itemize}
    \item \texttt{dist[v]} indica la distancia mínima de $s$ a $v$.
    \item \texttt{padre[v]} permite reconstruir un camino mínimo hacia
    $v$.
    \item Si \texttt{dist[v]} es $-1$, $v$ no es alcanzable desde $s$.
\end{itemize}

\textbf{Camino inexistente:}

Antes de reconstruir el camino conviene comprobar si el destino fue
alcanzado:

\begin{lstlisting}[style=cpp]
if (dist[destino] == -1) {
    cout << "No existe camino\n";
}
else {
    vector<int> camino =
        reconstruirCamino(destino, padre);

    for (int v : camino) {
        cout << v << ' ';
    }
}
\end{lstlisting}

\textbf{Complejidad:}

La ejecución de BFS utilizando una lista de adyacencia cuesta
$O(n+m)$. La reconstrucción de un camino cuesta $O(k)$, donde $k$ es
la cantidad de vértices que contiene el camino.

\begin{table}[H]
\raggedright
\scriptsize
\setlength{\tabcolsep}{3pt}
\renewcommand{\arraystretch}{1.1}
\begin{tabular}{|l|c|}
\hline
\textbf{Operación} & \textbf{Complejidad} \\
\hline
BFS & $O(n+m)$ \\
\hline
Reconstruir un camino & $O(k)$ \\
\hline
Memoria adicional & $O(n)$ \\
\hline
\end{tabular}
\caption{Complejidad de BFS para reconstrucción de caminos.}
\end{table}

\textbf{Nota:} BFS encuentra caminos mínimos en grafos donde todas las
aristas tienen el mismo costo. Si las aristas tienen pesos diferentes,
deben utilizarse algoritmos como \textbf{Dijkstra}, \textbf{0-1 BFS} o
\textbf{Bellman-Ford}, dependiendo de las características del problema.

\textbf{Regla práctica:} si el problema pide no solamente la distancia
mínima sino también \textbf{qué vértices forman el camino}, guarda un
arreglo \texttt{padre} durante el BFS. No es necesario almacenar todos
los caminos: basta con guardar un padre por vértice.

\subsection{DFS}

\textbf{Idea:} Depth-First Search (DFS) explora un grafo siguiendo un
camino lo más profundamente posible antes de regresar y explorar otras
ramas. A diferencia de BFS, que utiliza una \textbf{cola}, DFS utiliza
una \textbf{pila}.

DFS puede implementarse de forma recursiva o iterativa. En Competitive
Programming, la versión iterativa es útil cuando el grafo puede tener
muchos vértices y una implementación recursiva podría provocar un
\textbf{stack overflow}.

\textbf{Implementación iterativa:}

Se utiliza un \texttt{stack<int>} para almacenar los vértices que deben
ser procesados. Un vértice se marca como visitado cuando se introduce
en la pila, evitando que pueda agregarse múltiples veces.

\begin{lstlisting}[style=cpp]
void dfs(int inicio, vector<vector<int>>& adj,
         vector<bool>& visitado) {

    stack<int> st;

    visitado[inicio] = true;
    st.push(inicio);

    while (!st.empty()) {

        int u = st.top();
        st.pop();

        cout << u << ' ';

        for (int v : adj[u]) {

            if (visitado[v]) continue;

            visitado[v] = true;
            st.push(v);
        }
    }
}
\end{lstlisting}

El orden exacto en que se visitan los vértices depende del orden de los
vecinos en la lista de adyacencia. Lo importante es que cada vértice
alcanzable sea procesado una vez.

\textbf{Ejemplo:}

Consideremos el siguiente grafo:

\begin{lstlisting}[style=cpp]
//        0
//       / \
//      1   2
//     / \
//    3   4
\end{lstlisting}

La lista de adyacencia puede ser:

\begin{lstlisting}[style=cpp]
adj[0] = {1, 2};
adj[1] = {0, 3, 4};
adj[2] = {0};
adj[3] = {1};
adj[4] = {1};
\end{lstlisting}

Si comenzamos DFS desde $0$, la pila inicialmente contiene:

\begin{lstlisting}[style=cpp]
st = {0}
\end{lstlisting}

Se extrae $0$ y se agregan sus vecinos:

\begin{lstlisting}[style=cpp]
st = {1, 2}
\end{lstlisting}

El siguiente vértice que se extrae depende del orden en que fueron
agregados a la pila. Por ejemplo, si se extrae primero $2$, se procesa
esa rama antes de regresar a $1$.

\textbf{DFS en un grafo desconectado:}

Al igual que BFS, una sola ejecución de DFS solamente visita los
vértices alcanzables desde el origen. Para recorrer todo el grafo,
debemos iniciar un nuevo DFS desde cada vértice que todavía no haya
sido visitado.

\begin{lstlisting}[style=cpp]
vector<bool> visitado(n, false);

for (int u = 0; u < n; u++) {

    if (visitado[u]) continue;

    dfs(u, adj, visitado);
}
\end{lstlisting}

Este patrón permite procesar todas las componentes conexas del grafo.

\textbf{Contar componentes conexas:}

Podemos contar el número de componentes incrementando un contador cada
vez que sea necesario iniciar un nuevo DFS.

\begin{lstlisting}[style=cpp]
int componentes = 0;

vector<bool> visitado(n, false);

for (int u = 0; u < n; u++) {

    if (visitado[u]) continue;

    componentes++;
    dfs(u, adj, visitado);
}

cout << componentes << '\n';
\end{lstlisting}

\textbf{DFS en árboles:}

En un árbol no es necesario mantener un arreglo de visitados si se
conoce el padre del vértice actual. Al recorrer un vecino, simplemente
se evita regresar al padre.

Sin embargo, en un grafo general se recomienda utilizar
\texttt{visitado}, ya que pueden existir ciclos.

\begin{lstlisting}[style=cpp]
void dfs(int inicio, vector<vector<int>>& adj) {

    stack<pair<int, int>> st;

    st.push({inicio, -1});

    while (!st.empty()) {

        auto [u, padre] = st.top();
        st.pop();

        cout << u << ' ';

        for (int v : adj[u]) {

            if (v == padre) continue;

            st.push({v, u});
        }
    }
}
\end{lstlisting}

Esta última versión es apropiada para árboles, pero para grafos
generales debe utilizarse un arreglo de visitados.

\textbf{Aplicaciones:}

DFS es una de las herramientas fundamentales para resolver problemas
de grafos. Entre sus aplicaciones más comunes se encuentran:

\begin{itemize}
    \item Recorrer todos los vértices alcanzables desde un origen.
    \item Encontrar componentes conexas.
    \item Detectar ciclos.
    \item Construir árboles DFS.
    \item Encontrar puntos de articulación y puentes.
    \item Realizar ordenamiento topológico.
    \item Resolver problemas de conectividad.
    \item Recorrer árboles y estructuras jerárquicas.
\end{itemize}

\textbf{Complejidad:}

Utilizando una lista de adyacencia, cada vértice se procesa una sola
vez y cada arista se revisa una cantidad constante de veces. Por lo
tanto:

\[
O(n+m)
\]

La pila y el arreglo de visitados requieren $O(n)$ memoria adicional.

\begin{table}[H]
\raggedright
\scriptsize
\setlength{\tabcolsep}{3pt}
\renewcommand{\arraystretch}{1.1}
\begin{tabular}{|l|c|c|}
\hline
\textbf{Representación} &
\textbf{Tiempo} &
\textbf{Memoria adicional} \\
\hline
Lista de adyacencia &
$O(n+m)$ &
$O(n)$ \\
\hline
Matriz de adyacencia &
$O(n^2)$ &
$O(n)$ \\
\hline
\end{tabular}
\caption{Complejidad de DFS según la representación del grafo.}
\end{table}

\textbf{Nota:} DFS y BFS tienen la misma complejidad asintótica
$O(n+m)$ utilizando una lista de adyacencia, pero exploran el grafo de
manera diferente. BFS procesa por niveles mediante una cola, mientras
que DFS profundiza utilizando una pila.

\textbf{Regla práctica:} utiliza DFS cuando el problema dependa de
explorar profundamente una estructura, recorrer componentes, detectar
ciclos o realizar algún procesamiento basado en el orden de entrada y
salida de los vértices. Si el problema requiere distancias mínimas en
un grafo no ponderado, normalmente BFS es la opción adecuada.

\subsection{DFS para Conteo de Caminos}

\textbf{Idea:} en un DAG (grafo dirigido acíclico), contar cuántos
caminos distintos existen entre un nodo $u$ y un nodo destino fijo
es un problema de \textbf{DFS + memoización}: el número de caminos
desde $u$ hasta el destino es la suma de los caminos desde cada
vecino de $u$ hasta el destino, con caso base "ya llegué". Como
varios caminos pueden pasar por el mismo nodo intermedio, sin
memoización recalculas el mismo subárbol muchas veces —memoizar es
lo que baja esto de exponencial a lineal.

\textbf{¿Por qué usarlo?} Es el mismo patrón de "DFS + memo" que ya
usas en DP sobre árboles/DAGs, aplicado específicamente a contar
(no a minimizar/maximizar). Reconocer que un problema pide "número
de formas de llegar de A a B" en un grafo dirigido sin ciclos es la
señal para aplicar esto en vez de enumerar caminos explícitamente.

\textbf{Casos de uso:}

\begin{itemize}
	\item Contar caminos entre dos nodos fijos de un DAG.
	\item Contar, para cada nodo, cuántos caminos empiezan en él y
	terminan en algún nodo "hoja" (sin salidas) — mismo memo, sin
	destino fijo, sumando 1 en cada hoja.
	\item Variante con longitud exacta $k$: contar caminos de $u$ a
	$v$ usando \textbf{exactamente} $k$ aristas, memoizando sobre el
	par $(nodo, pasos\_restantes)$ en vez de solo $nodo$.
	\item Contar caminos módulo algún $M$ cuando el número de
	caminos crece demasiado (usa \texttt{long long} y aplica \%M en
	cada suma).
\end{itemize}

\textbf{Complejidad:} $O(V + E)$ con memoización, ya que cada nodo
se procesa una sola vez y cada arista se recorre una sola vez. Sin
memo, en el peor caso es exponencial (puede haber $O(2^V)$ caminos).

\begin{lstlisting}[style=cpp]
	struct ConteoCaminos {
		vector<vector<int>> adj;
		vector<long long> memo;
		int destino;
		
		// CREATE: guarda el grafo (debe ser DAG) y el nodo destino;
		// memo[u] = -1 significa "no calculado todavia".
		ConteoCaminos(vector<vector<int>>& adj_, int destino_) {
			adj = adj_;
			destino = destino_;
			memo.assign(adj.size(), -1);
		}                                          // O(1)
		
		// READ: numero de caminos distintos de u a destino.
		long long contar(int u) {
			if (u == destino) return 1;
			if (memo[u] != -1) return memo[u];
			
			long long total = 0;
			for (int v : adj[u])
			total += contar(v);
			
			return memo[u] = total;
		}                                          // O(1) amortizado
		// (O(V+E) la primera vez)
	};
\end{lstlisting}

\textbf{Ejemplo de funcionamiento:}

\begin{lstlisting}[style=cpp]
	// DAG:      0 -> 1 -> 3
	//           0 -> 2 -> 3
	//           1 -> 2
	
	vector<vector<int>> adj(4);
	adj[0] = {1, 2};
	adj[1] = {2, 3};
	adj[2] = {3};
	adj[3] = {};
	
	ConteoCaminos cc(adj, 3);
	
	cout << cc.contar(0);
	// 3  (0-1-3, 0-1-2-3, 0-2-3)
	
	cout << cc.contar(1);
	// 2  (1-3, 1-2-3)
\end{lstlisting}

\textbf{Nota:} si el grafo tiene ciclos, esta memoización no funciona
directamente (la recursión no termina) — primero hay que condensar
las componentes fuertemente conexas (SCC) en un DAG, y contar sobre
ese DAG condensado.


\subsection{DFS + Cerramiento Transitivo}

\textbf{Idea:} el cerramiento transitivo de un grafo dirigido
responde, para cualquier par $(u,v)$, la pregunta "¿existe algún
camino de $u$ a $v$?" en $O(1)$ por query, a costa de precomputar
—para cada nodo— el conjunto completo de nodos que alcanza. La forma
más directa es correr un DFS desde \textbf{cada} nodo y marcar todo
lo que visita.

\textbf{¿Por qué usarlo?} Cuando el grafo es pequeño/mediano y vas a
hacer \textbf{muchas} consultas de alcanzabilidad, precomputar todo
de una vez sale más barato que un DFS/BFS por consulta. Es el mismo
principio de "precomputo una vez, respondo O(1)" que en RMQ, aplicado
a alcanzabilidad en vez de a mínimos.

\textbf{Casos de uso:}

\begin{itemize}
	\item Responder $Q$ queries de "¿$u$ puede llegar a $v$?" cuando
	$Q$ es grande y $V$ es manejable ($V \le \sim 2000$–$5000$ según
	el límite de tiempo).
	\item Precomputar dependencias transitivas (si $a$ depende de
	$b$ y $b$ depende de $c$, entonces $a$ depende de $c$).
	\item Como paso intermedio antes de comprimir un DAG por SCC:
	primero condensas ciclos, y sobre el DAG resultante calculas el
	cerramiento transitivo.
\end{itemize}

\textbf{Complejidad:} $O(V \cdot (V+E))$ para construir (un DFS
completo por cada nodo), $O(1)$ por query. Con \texttt{std::bitset}
en vez de \texttt{vector<bool>}, y procesando los nodos en orden
topológico inverso (\texttt{reach[u] |= reach[v]} para cada vecino
$v$), el build baja a $O(V \cdot (V+E) / 64)$ gracias al paralelismo
a nivel de bits — útil si $V$ es más grande.

\begin{lstlisting}[style=cpp]
	struct CerramientoTransitivo {
		int n;
		vector<vector<int>> adj;
		vector<vector<bool>> alcanzable;   // alcanzable[u][v] = true si
		// u puede llegar a v
		
		// CREATE: corre un DFS completo desde cada nodo. O(V*(V+E)).
		CerramientoTransitivo(vector<vector<int>>& adj_, int n_) {
			adj = adj_;
			n = n_;
			alcanzable.assign(n, vector<bool>(n, false));
			
			for (int origen = 0; origen < n; origen++)
			dfs(origen, origen);
		}
		
		void dfs(int origen, int u) {
			for (int v : adj[u]) {
				if (!alcanzable[origen][v]) {
					alcanzable[origen][v] = true;
					dfs(origen, v);
				}
			}
		}
		
		// READ: es u puede llegar a v.
		bool query(int u, int v) {
			return alcanzable[u][v];
		}                                          // O(1)
	};
\end{lstlisting}

\textbf{Ejemplo de funcionamiento:}

\begin{lstlisting}[style=cpp]
	// grafo:    0 -> 1 -> 2
	//                \
	//                 -> 3
	//           4 (aislado)
	
	vector<vector<int>> adj(5);
	adj[0] = {1};
	adj[1] = {2, 3};
	adj[2] = {};
	adj[3] = {};
	adj[4] = {};
	
	CerramientoTransitivo ct(adj, 5);
	
	cout << ct.query(0, 3);
	// true  (0 -> 1 -> 3)
	
	cout << ct.query(2, 0);
	// false (no hay camino de 2 a 0)
	
	cout << ct.query(4, 0);
	// false (4 esta aislado)
\end{lstlisting}

\textbf{Regla práctica:} DFS "para contar" y DFS "para cerramiento
transitivo" comparten la misma mecánica de recorrido, pero acumulan
cosas distintas en cada llamada: uno suma caminos hacia un destino
fijo (memoizando por nodo), el otro marca alcanzabilidad hacia
\emph{todos} los destinos posibles (memoizando por par de nodos,
implícitamente, al no revisitar lo ya marcado). Si el grafo tiene
ciclos, ambos requieren primero pasar por condensación de SCC antes
de aplicarse tal cual.

\subsection{Flood Fill}

\textbf{Idea:} Flood Fill es "explorar todas las celdas conectadas a
una celda inicial que cumplen una condición" (mismo color, mismo
carácter, sin obstáculo, etc.), marcándolas como visitadas o
recoloreándolas. Es exactamente un DFS/BFS pero sobre una grilla en
vez de un grafo con listas de adyacencia explícitas: los "vecinos"
de una celda $(r,c)$ son sus hasta 4 (u 8) celdas adyacentes.

\textbf{¿Por qué usarlo?} Cualquier problema de grilla que hable de
"regiones", "islas", "áreas conectadas" o "pintar/rellenar" es Flood
Fill. La clave práctica es preferir la versión \textbf{iterativa}
(con \texttt{queue} o \texttt{stack} explícito) sobre la recursiva:
en grillas grandes ($R \times C$ del orden de $10^5$–$10^6$ celdas),
un DFS recursivo puede hacer \textit{stack overflow}.

\textbf{Casos de uso:}

\begin{itemize}
	\item Contar componentes conexas de una grilla ("número de
	islas"): un Flood Fill por cada celda no visitada que cumpla la
	condición, contando cuántas veces arrancas uno nuevo.
	\item Calcular el área (tamaño) de cada región conectada.
	\item "Paint bucket": recolorear toda una región conectada a un
	color nuevo.
	\item Determinar si dos celdas están en la misma región conectada.
\end{itemize}

\textbf{Complejidad:} $O(R \cdot C)$, ya que cada celda se visita
(y se encola) como máximo una vez.

\begin{lstlisting}[style=cpp]
	struct FloodFill {
		vector<vector<char>> grid;
		vector<vector<bool>> visitado;
		int R, C;
		
		int dr[4] = {-1, 1, 0, 0};
		int dc[4] = {0, 0, -1, 1};
		
		// CREATE: guarda la grilla (R x C).
		FloodFill(vector<vector<char>>& grid_) {
			grid = grid_;
			R = grid.size();
			C = grid[0].size();
			visitado.assign(R, vector<bool>(C, false));
		}                                          // O(1)
		
		// READ: marca (BFS iterativo) toda la region conectada a
		// (sr, sc) que comparte el mismo caracter, y devuelve su area.
		// Evita recursion para no hacer stack overflow en grillas grandes.
		int llenar(int sr, int sc) {
			if (visitado[sr][sc]) return 0;
			
			char objetivo = grid[sr][sc];
			queue<pair<int,int>> q;
			q.push({sr, sc});
			visitado[sr][sc] = true;
			int area = 0;
			
			while (!q.empty()) {
				auto [r, c] = q.front(); q.pop();
				area++;
				
				for (int d = 0; d < 4; d++) {
					int nr = r + dr[d], nc = c + dc[d];
					
					if (nr < 0 || nr >= R || nc < 0 || nc >= C) continue;
					if (visitado[nr][nc]) continue;
					if (grid[nr][nc] != objetivo) continue;
					
					visitado[nr][nc] = true;
					q.push({nr, nc});
				}
			}
			
			return area;
		}                                          // O(area de la region)
	};
\end{lstlisting}

\textbf{Ejemplo de funcionamiento:}

\begin{lstlisting}[style=cpp]
	// grid:
	// L L W W
	// L L W W
	// W W L L
	
	vector<vector<char>> grid = {
		{'L','L','W','W'},
		{'L','L','W','W'},
		{'W','W','L','L'}
	};
	
	FloodFill ff(grid);
	
	int islas = 0, areaMax = 0;
	for (int r = 0; r < 3; r++)
	for (int c = 0; c < 4; c++)
	if (grid[r][c] == 'L' && !ff.visitado[r][c]) {
		int area = ff.llenar(r, c);
		islas++;
		areaMax = max(areaMax, area);
	}
	
	cout << islas;    // 2
	cout << areaMax;  // 4  (el bloque de 2x2 de arriba-izquierda)
\end{lstlisting}


\subsection{Laberinto con BFS}

\textbf{Idea:} en una grilla donde te mueves de una celda a una
celda adyacente (arriba/abajo/izq/der) con \textbf{costo uniforme}
(cada paso cuesta 1), la ruta más corta desde un origen se calcula
con BFS: exploras la grilla "por capas" de distancia creciente, y la
primera vez que llegas a una celda, esa es su distancia mínima. Es el
mismo BFS de grafos de siempre, con la grilla actuando como grafo
implícito.

\textbf{¿Por qué usarlo?} Es la herramienta por defecto para
"camino más corto en grilla sin pesos / con obstáculos". Si el
problema menciona pesos distintos por celda (terreno con costo), esto
deja de ser BFS puro y pasa a ser Dijkstra o 0-1 BFS (deque) — pero
mientras cada paso cueste lo mismo, BFS es suficiente y más simple.

\textbf{Casos de uso:}

\begin{itemize}
	\item Número mínimo de pasos de un punto $A$ a un punto $B$ en
	un laberinto con paredes.
	\item Reconstruir el camino más corto explícito (no solo la
	distancia), guardando de dónde vino cada celda.
	\item BFS multi-fuente: varias celdas iniciales a la vez (p.ej.
	"tiempo mínimo para que el fuego llegue a cada celda"),
	encolando todas las fuentes con distancia 0 antes de empezar.
	\item Verificar alcanzabilidad ("¿se puede llegar de A a B?").
\end{itemize}

\textbf{Complejidad:} $O(R \cdot C)$, ya que cada celda se encola y
desencola como máximo una vez.

\begin{lstlisting}[style=cpp]
	struct LaberintoBFS {
		vector<vector<char>> grid;
		vector<vector<int>> dist;
		vector<vector<pair<int,int>>> padre;
		int R, C;
		
		int dr[4] = {-1, 1, 0, 0};
		int dc[4] = {0, 0, -1, 1};
		
		// CREATE: guarda la grilla. '#' = pared, cualquier otro caracter
		// = celda libre.
		LaberintoBFS(vector<vector<char>>& grid_) {
			grid = grid_;
			R = grid.size();
			C = grid[0].size();
		}                                          // O(1)
		
		// READ: corre BFS desde (sr, sc) y llena dist / padre para
		// todas las celdas alcanzables. dist queda en -1 donde no llego.
		void bfs(int sr, int sc) {
			dist.assign(R, vector<int>(C, -1));
			padre.assign(R, vector<pair<int,int>>(C, {-1, -1}));
			
			queue<pair<int,int>> q;
			q.push({sr, sc});
			dist[sr][sc] = 0;
			
			while (!q.empty()) {
				auto [r, c] = q.front(); q.pop();
				
				for (int d = 0; d < 4; d++) {
					int nr = r + dr[d], nc = c + dc[d];
					
					if (nr < 0 || nr >= R || nc < 0 || nc >= C) continue;
					if (grid[nr][nc] == '#') continue;
					if (dist[nr][nc] != -1) continue;
					
					dist[nr][nc] = dist[r][c] + 1;
					padre[nr][nc] = {r, c};
					q.push({nr, nc});
				}
			}
		}                                          // O(R*C)
		
		// READ: reconstruye el camino desde el origen del ultimo bfs()
		// hasta (er, ec), como lista de celdas en orden. Vacio si no
		// hay camino.
		vector<pair<int,int>> reconstruirCamino(int er, int ec) {
			if (dist[er][ec] == -1) return {};
			
			vector<pair<int,int>> camino;
			pair<int,int> actual = {er, ec};
			
			while (actual.first != -1) {
				camino.push_back(actual);
				actual = padre[actual.first][actual.second];
			}
			
			reverse(camino.begin(), camino.end());
			return camino;
		}                                          // O(longitud del camino)
	};
\end{lstlisting}

\textbf{Ejemplo de funcionamiento:}

\begin{lstlisting}[style=cpp]
	// grid:
	// . . . #
	// # # . #
	// . . . .
	
	vector<vector<char>> grid = {
		{'.','.','.','#'},
		{'#','#','.','#'},
		{'.','.','.','.'}
	};
	
	LaberintoBFS lb(grid);
	lb.bfs(0, 0);
	
	cout << lb.dist[2][3];
	// 5  (0,0)->(0,1)->(0,2)->(1,2)->(2,2)->(2,3)
	
	auto camino = lb.reconstruirCamino(2, 3);
	// [(0,0),(0,1),(0,2),(1,2),(2,2),(2,3)]
\end{lstlisting}

\textbf{Regla práctica:} si un problema de grilla pide "¿son
conectados?" o "área de una región", es Flood Fill (DFS o BFS, no
importa el orden de visita). Si pide "distancia mínima" o "menor
número de pasos", es BFS por capas (el orden de visita sí importa:
BFS garantiza que la primera vez que tocas una celda es con su
distancia mínima; DFS no lo garantiza).

\subsection{Tour de Caballo de Ajedrez}

\textbf{Idea:} el Tour de Caballo pide encontrar una secuencia de
movimientos de caballo que visite \textbf{cada casilla del tablero
	exactamente una vez}. No es un problema de "camino más corto" (no
hay pesos ni se compara contra alternativas), sino de
\textbf{existencia}: es backtracking puro — DFS que prueba un
movimiento, sigue explorando, y si se atasca (llega a una casilla sin
salidas válidas) deshace ("backtrack") y prueba otro movimiento.

\textbf{¿Por qué usarlo?} A diferencia de Flood Fill o BFS en
laberintos, aquí \textbf{no basta con visitar} las casillas
alcanzables: tienen que visitarse \textbf{todas, sin repetir}, en
un orden válido. Eso descarta BFS/DFS simple y obliga a backtracking
con una condición de poda, porque la fuerza bruta pura ($8^{n}$
combinaciones de movimientos posibles) es intratable incluso para
tableros chicos.

\textbf{Casos de uso:}

\begin{itemize}
	\item Encontrar (o verificar la existencia de) un tour de caballo
	completo en un tablero $n \times n$.
	\item Tour \textbf{cerrado} (el último movimiento vuelve a
	atacar la casilla inicial): misma idea, chequeo extra al final.
	\item Cualquier problema de "recorrer todas las celdas exactamente
	una vez respetando un patrón de movimiento" (no exclusivo del
	caballo): el patrón general es backtracking + poda de
	Warnsdorff.
\end{itemize}

\textbf{Complejidad:} exponencial en el peor caso sin poda
($O(8^{n^2})$ aproximado). Con la \textbf{heurística de Warnsdorff}
—en cada paso, moverse primero a la casilla accesible con
\textbf{menos} salidas futuras disponibles— en la práctica encuentra
un tour en tiempo casi lineal para tableros razonables ($n \le 100$),
aunque no hay garantía teórica formal de que siempre funcione.

\begin{lstlisting}[style=cpp]
	struct TourCaballo {
		int n;
		vector<vector<int>> tablero;   // tablero[r][c] = orden de visita
		// (-1 si no visitada)
		int dr[8] = {-2,-2,-1,-1, 1, 1, 2, 2};
		int dc[8] = {-1, 1,-2, 2,-2, 2,-1, 1};
		
		// CREATE: tablero n x n vacio.
		TourCaballo(int n_) {
			n = n_;
			tablero.assign(n, vector<int>(n, -1));
		}                                          // O(n^2)
		
		bool dentro(int r, int c) {
			return r >= 0 && r < n && c >= 0 && c < n;
		}
		
		// cuenta a cuantas casillas libres puede saltar (r,c) - usado
		// por Warnsdorff para ordenar los candidatos.
		int gradoSalida(int r, int c) {
			int cont = 0;
			for (int d = 0; d < 8; d++) {
				int nr = r + dr[d], nc = c + dc[d];
				if (dentro(nr, nc) && tablero[nr][nc] == -1) cont++;
			}
			return cont;
		}
		
		// READ: intenta construir un tour completo empezando en
		// (sr, sc). Devuelve true si lo logro (tablero queda lleno).
		bool resolver(int sr, int sc) {
			tablero[sr][sc] = 0;
			return backtrack(sr, sc, 1);
		}                                          // O(casi lineal) con
		// Warnsdorff en la
		// practica
		
		bool backtrack(int r, int c, int movimiento) {
			if (movimiento == n * n) return true;
			
			// candidatos ordenados por Warnsdorff: menor grado de
			// salida primero
			vector<tuple<int,int,int>> candidatos;
			for (int d = 0; d < 8; d++) {
				int nr = r + dr[d], nc = c + dc[d];
				if (dentro(nr, nc) && tablero[nr][nc] == -1)
				candidatos.push_back({gradoSalida(nr, nc), nr, nc});
			}
			sort(candidatos.begin(), candidatos.end());
			
			for (auto& [grado, nr, nc] : candidatos) {
				tablero[nr][nc] = movimiento;
				
				if (backtrack(nr, nc, movimiento + 1)) return true;
				
				tablero[nr][nc] = -1;   // deshacer (backtrack real)
			}
			
			return false;
		}
	};
\end{lstlisting}

\textbf{Ejemplo de funcionamiento:}

\begin{lstlisting}[style=cpp]
	TourCaballo tc(5);
	
	bool encontrado = tc.resolver(0, 0);
	
	cout << encontrado;
	// true
	
	// tc.tablero[r][c] tiene el orden (0..24) en que se visito
	// cada casilla, por ejemplo:
	//  0  3 16  9 22
	// 13  8 21  4 17
	//  2  1 12 23 10
	//  7 14 19 18  5
	// 20  ... etc.
\end{lstlisting}

\textbf{Nota:} sin la poda de Warnsdorff, el backtrack puro también
es correcto pero mucho más lento (puede no terminar en tiempo útil
para $n > 6$ u $8$). Warnsdorff no cambia la corrección del
algoritmo, solo el \textbf{orden} en que se prueban los candidatos —
sigue siendo backtracking con deshacer explícito.


\subsection{Mínimo Número de Movimientos de Caballo}

\textbf{Idea:} a diferencia del Tour de Caballo, aquí no importa
visitar todo el tablero — solo se pide la \textbf{distancia mínima}
(en número de saltos) entre una casilla origen y una casilla
destino. Eso es "camino más corto con costo uniforme" sobre un grafo
implícito donde cada casilla es un nodo y sus hasta 8 movimientos de
caballo son las aristas: exactamente el mismo patrón que
Laberinto con BFS, cambiando el arreglo de movimientos de 4/8
direcciones ortogonales por los 8 saltos en "L" del caballo.

\textbf{¿Por qué usarlo?} Es tentador intentar resolver esto con
fórmulas o backtracking (como el Tour), pero al ser "distancia
mínima" y no "visitar todo", BFS es tanto más simple como
garantizado óptimo — la primera vez que BFS alcanza la casilla
destino, esa es la respuesta.

\textbf{Casos de uso:}

\begin{itemize}
	\item Mínimo número de movimientos de caballo entre dos casillas
	de un tablero $n \times n$ (o infinito, acotando la búsqueda).
	\item Mismo problema con casillas bloqueadas/prohibidas: se
	salta la casilla en el BFS igual que una pared en un laberinto.
	\item Generalización a "cuál es la distancia mínima con un patrón
	de movimiento arbitrario" (no solo caballo): el patrón se sigue
	resolviendo con BFS cambiando \texttt{dr/dc}.
\end{itemize}

\textbf{Complejidad:} $O(n^2)$ para un tablero $n \times n$, ya que
cada casilla se encola como máximo una vez y desde cada una se
prueban 8 movimientos constantes.

\begin{lstlisting}[style=cpp]
	struct MinMovimientosCaballo {
		int n;
		vector<vector<int>> dist;
		
		int dr[8] = {-2,-2,-1,-1, 1, 1, 2, 2};
		int dc[8] = {-1, 1,-2, 2,-2, 2,-1, 1};
		
		// CREATE: tablero n x n.
		MinMovimientosCaballo(int n_) {
			n = n_;
		}                                          // O(1)
		
		bool dentro(int r, int c) {
			return r >= 0 && r < n && c >= 0 && c < n;
		}
		
		// READ: BFS desde (sr, sc); dist[r][c] = minimo de movimientos
		// de caballo para llegar a (r, c), -1 si no aplica (no deberia
		// pasar en un tablero cerrado y conexo).
		int minMovimientos(int sr, int sc, int er, int ec) {
			dist.assign(n, vector<int>(n, -1));
			
			queue<pair<int,int>> q;
			q.push({sr, sc});
			dist[sr][sc] = 0;
			
			while (!q.empty()) {
				auto [r, c] = q.front(); q.pop();
				
				if (r == er && c == ec) return dist[r][c];
				
				for (int d = 0; d < 8; d++) {
					int nr = r + dr[d], nc = c + dc[d];
					
					if (!dentro(nr, nc)) continue;
					if (dist[nr][nc] != -1) continue;
					
					dist[nr][nc] = dist[r][c] + 1;
					q.push({nr, nc});
				}
			}
			
			return -1;   // no alcanzable (no deberia pasar sin bloqueos)
		}                                          // O(n^2)
	};
\end{lstlisting}

\textbf{Ejemplo de funcionamiento:}

\begin{lstlisting}[style=cpp]
	MinMovimientosCaballo mc(8);
	
	cout << mc.minMovimientos(0, 0, 1, 2);
	// 1  (un solo salto en L)
	
	cout << mc.minMovimientos(0, 0, 7, 7);
	// 6  (esquina a esquina en un tablero 8x8)
	
	cout << mc.minMovimientos(0, 0, 0, 1);
	// 3  (casilla adyacente, contraintuitivamente lejos para un caballo)
\end{lstlisting}

\textbf{Regla práctica:} misma distinción que Flood Fill vs.
Laberinto con BFS, aplicada al caballo: si el problema pide
\textbf{"visitar todo"} sin repetir, es backtracking (Tour de
Caballo, con Warnsdorff como poda). Si pide \textbf{"distancia
	mínima entre dos puntos"}, es BFS (Mínimo Número de Movimientos), sin
importar qué tan raro sea el patrón de movimiento — el patrón solo
define el arreglo de vecinos.

\subsection{Serpientes y Escaleras}

\textbf{Idea:} el tablero de Serpientes y Escaleras se modela como un
grafo dirigido: cada casilla $1..n$ es un nodo, y desde cada casilla
$i$ hay hasta 6 aristas (una por cada valor del dado, $i+1$ hasta
$i+6$), \textbf{redirigidas} si la casilla de llegada tiene una
serpiente o escalera. Pedir el "mínimo número de tiradas de dado para
llegar de la casilla 1 a la $n$" es, otra vez, camino más corto con
costo uniforme (cada tirada cuenta como 1 paso) — el mismo BFS de
Laberinto con BFS y Mínimo Número de Movimientos de Caballo, cambiando
qué cuenta como "vecino" de una casilla.

\textbf{¿Por qué usarlo?} La parte no trivial no es el BFS en sí
—es idéntico al de siempre— sino \textbf{modelar} serpientes y
escaleras correctamente: no son casillas nuevas, son
\textbf{atajos/retrocesos automáticos}. Si aterrizas en la casilla
de la base de una escalera o la cabeza de una serpiente, tu posición
real cambia \textbf{sin gastar otra tirada}. Eso se resuelve
precomputando, para cada casilla, a dónde "realmente" te manda
(\texttt{destinoReal[i]}), y haciendo BFS sobre ese arreglo en vez de
sobre los números crudos del dado.

\textbf{Casos de uso:}

\begin{itemize}
	\item Mínimo número de tiradas de dado para llegar de la casilla
	1 (o cualquier inicio) a la casilla $n$.
	\item Variante donde el dado no es 1–6 sino un conjunto arbitrario
	de valores posibles: solo cambia el rango del \texttt{for} al
	generar vecinos.
	\item Cualquier tablero lineal con "teletransportes" fijos
	(serpiente = te manda hacia atrás, escalera = te manda hacia
	adelante) donde se pide distancia mínima en pasos.
\end{itemize}

\textbf{Complejidad:} $O(n)$, ya que el grafo tiene $n$ nodos y a lo
más $6n$ aristas — cada casilla se encola una sola vez.

\begin{lstlisting}[style=cpp]
	struct SerpientesYEscaleras {
		int n;
		vector<int> destinoReal;   // destinoReal[i] = a donde termina
		// realmente el jugador si cae en i
		// (i mismo si no hay serpiente/escalera)
		
		// CREATE: tablero de n casillas (1-indexado). saltos[i] = j
		// significa "serpiente o escalera de i a j"; las casillas sin
		// entrada en el mapa no tienen salto.
		SerpientesYEscaleras(int n_, map<int,int>& saltos) {
			n = n_;
			destinoReal.assign(n + 1, 0);
			
			for (int i = 1; i <= n; i++)
			destinoReal[i] = saltos.count(i) ? saltos[i] : i;
		}                                          // O(n)
		
		// READ: minimo numero de tiradas de dado (1-6) para llegar de
		// origen a n. -1 si es imposible (no deberia pasar en un
		// tablero estandar).
		int minTiradas(int origen) {
			vector<int> dist(n + 1, -1);
			queue<int> q;
			
			q.push(origen);
			dist[origen] = 0;
			
			while (!q.empty()) {
				int actual = q.front(); q.pop();
				
				if (actual == n) return dist[actual];
				
				for (int dado = 1; dado <= 6; dado++) {
					int siguiente = actual + dado;
					if (siguiente > n) continue;
					
					siguiente = destinoReal[siguiente];   // aplica
					// serpiente/escalera
					
					if (dist[siguiente] != -1) continue;
					
					dist[siguiente] = dist[actual] + 1;
					q.push(siguiente);
				}
			}
			
			return -1;
		}                                          // O(n)
	};
\end{lstlisting}

\textbf{Ejemplo de funcionamiento:}

\begin{lstlisting}[style=cpp]
	// tablero de 10 casillas:
	// escalera de 2 a 8
	// serpiente  de 9 a 3
	
	map<int,int> saltos = {{2, 8}, {9, 3}};
	SerpientesYEscaleras se(10, saltos);
	
	cout << se.minTiradas(1);
	// 2
	// tirada 1: 1 -> 2 -> (escalera) -> 8
	// tirada 2: 8 -> 10
\end{lstlisting}

\textbf{Regla práctica:} cualquier tablero lineal donde te mueves por
"reglas de dado" (o cualquier conjunto de deltas fijos) y hay
teletransportes automáticos al aterrizar en ciertas casillas, se
reduce a BFS de toda la vida sobre un grafo de $n$ nodos —
\textbf{la única parte específica del problema es construir bien el
	arreglo \texttt{destinoReal}} antes de correr el BFS; el algoritmo de
búsqueda no cambia respecto a los casos anteriores.

% =========================================================
% CONECTIVIDAD
% =========================================================
\section{Conectividad}
\subsection{Componentes Conexas}

\textbf{Idea:} en un grafo \textbf{no dirigido}, una componente
conexa es un grupo maximal de nodos donde cualquiera puede llegar a
cualquiera. Encontrarlas todas es simplemente: recorrer los nodos en
orden, y cada vez que encuentres uno \textbf{no visitado}, lanzar un
DFS/BFS desde ahí para marcar todo su grupo con un mismo
identificador — ese DFS/BFS nunca se sale del grupo porque, por
definición, no hay aristas hacia afuera de una componente.

\textbf{¿Por qué usarlo?} Es el primer paso obligatorio de casi
cualquier problema de grafos no dirigidos: antes de responder
"¿$u$ y $v$ están conectados?" o de aplicar un algoritmo que asume
un grafo conexo, necesitas saber cuántas componentes hay y a cuál
pertenece cada nodo. La alternativa es \textbf{Union-Find}: si solo
necesitas responder "¿misma componente?" con aristas que van
llegando una a una (sin necesariamente construir listas de
adyacencia), Union-Find es más directo; si ya tienes el grafo
completo y quieres además reconstruir/recorrer cada componente,
DFS/BFS es más natural.

\textbf{Casos de uso:}

\begin{itemize}
	\item Contar cuántas componentes conexas tiene un grafo (o una
	grilla vista como grafo — mismo patrón que Flood Fill, pero
	sobre listas de adyacencia en vez de vecinos de grilla).
	\item Etiquetar cada nodo con el id de su componente, para
	responder "¿misma componente?" en $O(1)$ tras el precómputo.
	\item Encontrar el tamaño de la componente más grande.
	\item Verificar si el grafo completo es conexo (una sola
	componente).
\end{itemize}

\textbf{Complejidad:} $O(V + E)$, ya que es un DFS/BFS normal que
simplemente se relanza desde cada nodo no visitado — en total, cada
nodo y cada arista se procesan una sola vez sin importar cuántas
componentes haya.

\begin{lstlisting}[style=cpp]
	struct ComponentesConexas {
		vector<vector<int>> adj;
		vector<int> comp;      // comp[u] = id de la componente de u, -1
		// si no se ha calculado
		int numComponentes;
		
		// CREATE: guarda el grafo y calcula todas las componentes de
		// una vez, recorriendo cada nodo no visitado.
		ComponentesConexas(vector<vector<int>>& adj_, int n) {
			adj = adj_;
			comp.assign(n, -1);
			numComponentes = 0;
			
			for (int u = 0; u < n; u++) {
				if (comp[u] == -1) {
					dfs(u, numComponentes);
					numComponentes++;
				}
			}
		}                                          // O(V + E)
		
		void dfs(int u, int id) {
			comp[u] = id;
			for (int v : adj[u])
			if (comp[v] == -1)
			dfs(v, id);
		}
		
		// READ: u y v estan en la misma componente?
		bool mismaComponente(int u, int v) {
			return comp[u] == comp[v];
		}                                          // O(1)
	};
\end{lstlisting}

\textbf{Ejemplo de funcionamiento:}

\begin{lstlisting}[style=cpp]
	// grafo:  0 - 1    2 - 3    4
	//         (dos componentes de tamano 2, una de tamano 1)
	
	vector<vector<int>> adj(5);
	adj[0] = {1}; adj[1] = {0};
	adj[2] = {3}; adj[3] = {2};
	adj[4] = {};
	
	ComponentesConexas cc(adj, 5);
	
	cout << cc.numComponentes;
	// 3
	
	cout << cc.mismaComponente(0, 1);
	// true
	
	cout << cc.mismaComponente(1, 2);
	// false
\end{lstlisting}


\subsection{Detección de Ciclos}

\textbf{Idea:} detectar si un grafo tiene un ciclo se hace con DFS,
pero la regla para reconocer un ciclo \textbf{cambia según el grafo
	sea dirigido o no}: en un grafo no dirigido, cualquier arista hacia
un nodo ya visitado que \textbf{no sea tu padre inmediato} es un
ciclo; en un grafo dirigido, una arista hacia un nodo ya visitado
\textbf{solo} es ciclo si ese nodo sigue en la pila de recursión
actual (todavía no terminaste de procesarlo) — si ya terminaste con
él, es una arista válida hacia otra rama, no un ciclo.

\textbf{¿Por qué usarlo?} Confundir estas dos reglas es el error más
común al detectar ciclos: aplicar la regla de "no dirigido" (solo
evitar al padre) a un grafo dirigido da falsos positivos, porque en
dirigido puedes visitar el mismo nodo desde dos ramas distintas sin
que eso sea un ciclo.

\textbf{Casos de uso:}

\begin{itemize}
	\item Verificar si un grafo dirigido es un DAG (sin ciclos) antes
	de aplicar orden topológico, DP sobre DAG, etc.
	\item Detectar dependencias circulares (tareas que dependen unas
	de otras en círculo, imposibles de ordenar).
	\item En grafos no dirigidos: verificar si un grafo es un árbol
	($V-1$ aristas, conexo, sin ciclos) o encontrar el primer ciclo
	para removerlo.
\end{itemize}

\textbf{Complejidad:} $O(V + E)$ en ambos casos — un DFS normal con
una condición extra por arista, sin repetir trabajo.

\subsubsection{En grafos no dirigidos (DFS con padre)}

\begin{lstlisting}[style=cpp]
	struct CicloNoDirigido {
		vector<vector<int>> adj;
		vector<bool> visitado;
		
		// CREATE: guarda el grafo.
		CicloNoDirigido(vector<vector<int>>& adj_) {
			adj = adj_;
			visitado.assign(adj.size(), false);
		}                                          // O(1)
		
		// READ: existe algun ciclo en todo el grafo (revisa todas las
		// componentes).
		bool tieneCiclo() {
			for (int u = 0; u < (int)adj.size(); u++)
			if (!visitado[u] && dfs(u, -1))
			return true;
			return false;
		}                                          // O(V + E)
		
		bool dfs(int u, int padre) {
			visitado[u] = true;
			
			for (int v : adj[u]) {
				if (!visitado[v]) {
					if (dfs(v, u)) return true;
				} else if (v != padre) {
					return true;   // arista hacia un visitado que no es
					// el padre -> ciclo
				}
			}
			
			return false;
		}
	};
\end{lstlisting}

\textbf{Ejemplo de funcionamiento:}

\begin{lstlisting}[style=cpp]
	// grafo:  0 - 1 - 2 - 0   (triangulo, tiene ciclo)
	//         3 - 4           (arbol, sin ciclo)
	
	vector<vector<int>> adj(5);
	adj[0] = {1, 2}; adj[1] = {0, 2}; adj[2] = {0, 1};
	adj[3] = {4};    adj[4] = {3};
	
	CicloNoDirigido cnd(adj);
	
	cout << cnd.tieneCiclo();
	// true  (por el triangulo 0-1-2)
\end{lstlisting}

\subsubsection{En grafos dirigidos (DFS con estados / colores)}

\textbf{Idea:} cada nodo tiene 3 estados durante el DFS: \texttt{0 =
	no visitado}, \texttt{1 = en la pila de recursión actual (gris)},
\texttt{2 = totalmente procesado (negro)}. Una arista hacia un nodo
\textbf{gris} es un back edge real (ciclo); una arista hacia un nodo
\textbf{negro} es válida (cross/forward edge, no ciclo).

\begin{lstlisting}[style=cpp]
	struct CicloDirigido {
		vector<vector<int>> adj;
		vector<int> estado;   // 0 = blanco, 1 = gris, 2 = negro
		
		// CREATE: guarda el grafo dirigido.
		CicloDirigido(vector<vector<int>>& adj_) {
			adj = adj_;
			estado.assign(adj.size(), 0);
		}                                          // O(1)
		
		// READ: existe algun ciclo en todo el grafo.
		bool tieneCiclo() {
			for (int u = 0; u < (int)adj.size(); u++)
			if (estado[u] == 0 && dfs(u))
			return true;
			return false;
		}                                          // O(V + E)
		
		bool dfs(int u) {
			estado[u] = 1;   // entra a la pila de recursion (gris)
			
			for (int v : adj[u]) {
				if (estado[v] == 1) return true;        // back edge -> ciclo
				if (estado[v] == 0 && dfs(v)) return true;
			}
			
			estado[u] = 2;   // termino con u (negro)
			return false;
		}
	};
\end{lstlisting}

\textbf{Ejemplo de funcionamiento:}

\begin{lstlisting}[style=cpp]
	// grafo dirigido:  0 -> 1 -> 2 -> 0   (ciclo)
	//                  3 -> 1             (sin ciclo: 3->1 es valida,
	//                                      1 ya esta siendo procesado
	//                                      solo si viene de la MISMA
	//                                      rama de recursion)
	
	vector<vector<int>> adj(4);
	adj[0] = {1};
	adj[1] = {2};
	adj[2] = {0};
	adj[3] = {1};
	
	CicloDirigido cd(adj);
	
	cout << cd.tieneCiclo();
	// true  (0 -> 1 -> 2 -> 0)
\end{lstlisting}

\textbf{Regla práctica:} para saber cuál versión aplicar, primero
pregúntate si el grafo es dirigido. Si no lo es, basta con recordar
"de dónde vine" (el padre) y listo. Si es dirigido, "de dónde vine"
no alcanza — necesitas saber si el nodo destino sigue \textbf{activo
	en la recursión actual} (gris) o ya quedó cerrado (negro), porque en
dirigido sí es válido llegar dos veces al mismo nodo desde ramas
distintas sin que eso sea un ciclo.



\subsection{Puentes}

\textbf{Idea:} en un grafo no dirigido, un \textbf{puente} es una
arista que, al quitarla, \textbf{desconecta} el grafo (aumenta el
número de componentes conexas). Se encuentran con un solo DFS,
llevando dos valores por nodo: \texttt{tin[u]} (el momento en que el
DFS \textbf{descubre} a $u$, un contador que se incrementa en cada
visita) y \texttt{low[u]} (el \texttt{tin} \textbf{más bajo}
alcanzable desde $u$ usando cualquier camino del árbol DFS más, como
mucho, \textbf{una} arista de "retroceso" hacia un ancestro). La
arista $(u,v)$ del árbol DFS es puente si y solo si
$\texttt{low}[v] > \texttt{tin}[u]$: eso significa que, desde el
subárbol de $v$, \textbf{no hay ningún atajo} de vuelta hacia $u$ o
más arriba —quitar esa arista realmente aísla el subárbol de $v$.

\textbf{¿Por qué usarlo?} Es la aplicación clásica del algoritmo de
Tarjan de \texttt{low-link}: en vez de, para cada arista, quitarla y
volver a correr un DFS/BFS completo para ver si el grafo se
desconecta ($O(E \cdot (V+E))$ en total), un solo DFS con
\texttt{low-link} responde para \textbf{todas} las aristas a la vez
en tiempo lineal.

\textbf{Casos de uso:}

\begin{itemize}
	\item Encontrar todas las conexiones "críticas" de una red (si
	se cae ese enlace, la red se parte en dos).
	\item Como paso previo para construir el \textbf{árbol de
		componentes 2-arista-conexas} (contraer cada componente
	2-arista-conexa a un nodo y dejar solo los puentes como aristas).
	\item Verificar si el grafo sigue siendo conexo tras remover
	cualquier arista dada (equivalente a "¿esa arista es puente?").
\end{itemize}

\textbf{Complejidad:} $O(V + E)$, un solo DFS con trabajo constante
extra por arista.

\begin{lstlisting}[style=cpp]
	struct Puentes {
		vector<vector<pair<int,int>>> adj;   // adj[u] = {(v, idArista)}
		vector<int> tin, low;
		vector<bool> visitado;
		vector<pair<int,int>> puentes;       // aristas (u,v) que son puente
		int timer;
		
		// CREATE: guarda el grafo (con id de arista para no reusar la
		// misma arista en sentido contrario) y corre el DFS de Tarjan
		// desde cada componente.
		Puentes(vector<vector<pair<int,int>>>& adj_, int n) {
			adj = adj_;
			tin.assign(n, -1);
			low.assign(n, -1);
			visitado.assign(n, false);
			timer = 0;
			
			for (int u = 0; u < n; u++)
			if (!visitado[u])
			dfs(u, -1);
		}                                          // O(V + E)
		
		void dfs(int u, int idAristaPadre) {
			visitado[u] = true;
			tin[u] = low[u] = timer++;
			
			for (auto& [v, idArista] : adj[u]) {
				if (idArista == idAristaPadre) continue;   // no te
				// regreses
				// por la
				// misma
				// arista
				
				if (visitado[v]) {
					low[u] = min(low[u], tin[v]);   // back edge
				} else {
					dfs(v, idArista);
					low[u] = min(low[u], low[v]);
					
					if (low[v] > tin[u])
					puentes.push_back({u, v});
				}
			}
		}
	};
\end{lstlisting}

\textbf{Ejemplo de funcionamiento:}

\begin{lstlisting}[style=cpp]
	// grafo:   0 - 1 - 2 - 0     (triangulo)
	//                  \
	//                   3 - 4    (colgando del triangulo)
	
	vector<vector<pair<int,int>>> adj(5);
	auto addEdge = [&](int u, int v, int id) {
		adj[u].push_back({v, id});
		adj[v].push_back({u, id});
	};
	addEdge(0, 1, 0);
	addEdge(1, 2, 1);
	addEdge(2, 0, 2);
	addEdge(2, 3, 3);
	addEdge(3, 4, 4);
	
	Puentes p(adj, 5);
	
	for (auto& [u, v] : p.puentes)
	cout << u << "-" << v << " ";
	// 2-3  3-4
	// (las aristas del triangulo NO son puente: hay un ciclo,
	//  siempre queda un camino alternativo)
\end{lstlisting}


\subsection{Puntos de Articulación}

\textbf{Idea:} un \textbf{punto de articulación} (o \textit{cut
	vertex}) es un nodo que, al quitarlo (junto con sus aristas),
\textbf{aumenta} el número de componentes conexas. Usa el mismo
\texttt{tin}/\texttt{low} que Puentes, pero la condición es distinta
y tiene \textbf{dos casos}: (1) si $u$ es la \textbf{raíz} del árbol
DFS, es articulación si tiene \textbf{2 o más hijos} en el árbol DFS
(cada hijo es un subárbol que solo se conecta al resto \textit{a
	través de} $u$); (2) si $u$ \textbf{no} es la raíz, es articulación
si tiene algún hijo $v$ con $\texttt{low}[v] \ge \texttt{tin}[u]$
(el subárbol de $v$ no tiene ningún atajo que lo saque \textbf{por
	encima} de $u$, así que sin $u$ ese subárbol queda aislado).

\textbf{¿Por qué usarlo?} Es el "primo" de Puentes: mismo DFS,
mismos \texttt{tin}/\texttt{low}, pero identificando nodos críticos
en vez de aristas críticas. La diferencia clave con Puentes está en
el operador de comparación ($\ge$ en vez de $>$) y en el caso especial
de la raíz — con solo $>$ (la condición de puente) se pierde el caso
donde $v$ puede volver exactamente a $u$ (eso no aísla la arista,
pero sí deja a $u$ como único punto de paso).

\textbf{Casos de uso:}

\begin{itemize}
	\item Encontrar los nodos "críticos" de una red (si ese nodo
	falla, la red se parte).
	\item Como paso previo para construir componentes
	\textbf{biconexas} (bloques del grafo separados por puntos de
	articulación).
	\item Distinguir de Puentes: un grafo puede tener puntos de
	articulación sin tener ningún puente (ej. dos triángulos que
	comparten un solo vértice) — son propiedades independientes.
\end{itemize}

\textbf{Complejidad:} $O(V + E)$, mismo costo que Puentes: un DFS con
trabajo constante extra por nodo/arista.

\begin{lstlisting}[style=cpp]
	struct PuntosArticulacion {
		vector<vector<int>> adj;
		vector<int> tin, low;
		vector<bool> visitado, esArticulacion;
		int timer;
		
		// CREATE: guarda el grafo y corre el DFS de Tarjan desde cada
		// componente.
		PuntosArticulacion(vector<vector<int>>& adj_, int n) {
			adj = adj_;
			tin.assign(n, -1);
			low.assign(n, -1);
			visitado.assign(n, false);
			esArticulacion.assign(n, false);
			timer = 0;
			
			for (int u = 0; u < n; u++)
			if (!visitado[u])
			dfs(u, -1);
		}                                          // O(V + E)
		
		void dfs(int u, int padre) {
			visitado[u] = true;
			tin[u] = low[u] = timer++;
			int hijos = 0;
			
			for (int v : adj[u]) {
				if (v == padre) continue;   // ignora SOLO una vez la
				// arista de vuelta al padre
				// (si hay aristas multiples
				// u-padre, usar id de arista
				// como en Puentes)
				
				if (visitado[v]) {
					low[u] = min(low[u], tin[v]);   // back edge
				} else {
					dfs(v, u);
					low[u] = min(low[u], low[v]);
					hijos++;
					
					if (padre != -1 && low[v] >= tin[u])
					esArticulacion[u] = true;
				}
			}
			
			if (padre == -1 && hijos >= 2)
			esArticulacion[u] = true;   // caso especial: raiz con
			// 2+ hijos en el arbol DFS
		}
	};
\end{lstlisting}

\textbf{Ejemplo de funcionamiento:}

\begin{lstlisting}[style=cpp]
	// grafo:   0 - 1 - 2 - 0     (triangulo)
	//                  \
	//                   3 - 4    (colgando del triangulo)
	
	vector<vector<int>> adj(5);
	adj[0] = {1, 2}; adj[1] = {0, 2}; adj[2] = {0, 1, 3};
	adj[3] = {2, 4}; adj[4] = {3};
	
	PuntosArticulacion pa(adj, 5);
	
	for (int u = 0; u < 5; u++)
	if (pa.esArticulacion[u])
	cout << u << " ";
	// 2  3
	// (2 conecta el triangulo con la cola 3-4; 3 conecta 2 con 4)
\end{lstlisting}

\textbf{Regla práctica:} Puentes y Puntos de Articulación comparten
exactamente el esqueleto del DFS (\texttt{tin}/\texttt{low}, back
edges hacia ancestros) — la diferencia está solo en \textbf{qué}
marcas: en Puentes marcas la \textbf{arista} $(u,v)$ cuando
$\texttt{low}[v] > \texttt{tin}[u]$ (estricto); en Puntos de
Articulación marcas el \textbf{nodo} $u$ cuando
$\texttt{low}[v] \ge \texttt{tin}[u]$ (no estricto, más el caso
especial de la raíz). Si programas uno, el otro es casi copiar y
pegar cambiando esos dos detalles.

\subsection{Componentes Biconectados}

\textbf{Idea:} una componente biconexa (o \textit{2-vertex-connected
	component}, "bloque") es un conjunto \textbf{maximal de aristas}
donde no hay ningún punto de articulación que las separe —
equivalentemente, cualquier par de aristas del bloque está en algún
ciclo común. Se calculan con el mismo DFS de \texttt{tin}/\texttt{low}
de Puentes y Puntos de Articulación, agregando una \textbf{pila de
	aristas}: apilas cada arista al recorrerla, y cuando detectas la
condición de punto de articulación ($\texttt{low}[v] \ge
\texttt{tin}[u]$), \textbf{desapilas} todo hasta la arista $(u,v)$
inclusive — ese grupo desapilado es una componente biconexa completa.

\textbf{¿Por qué usarlo?} Puntos de Articulación te dice \textbf{qué
	nodos} son críticos, pero no cómo se agrupan las aristas alrededor de
ellos. Componentes Biconectadas completa el panorama: parte el grafo
en bloques tales que cada punto de articulación es exactamente el
nodo compartido entre dos o más bloques — es la base para construir
el \textbf{block-cut tree} (árbol de bloques y puntos de
articulación), útil cuando el problema pregunta por relaciones entre
bloques en vez de solo "¿es articulación sí/no?".

\textbf{Casos de uso:}

\begin{itemize}
	\item Agrupar aristas en bloques maximales sin puntos de
	articulación internos.
	\item Construir el block-cut tree para responder queries tipo
	"¿cuántos bloques hay que atravesar para ir de $u$ a $v$?".
	\item Verificar si el grafo completo es biconexo (una sola
	componente biconexa que cubre todas las aristas, y ningún punto
	de articulación).
\end{itemize}

\textbf{Complejidad:} $O(V + E)$ — mismo DFS de \texttt{tin}/
\texttt{low} de siempre, más trabajo amortizado constante por arista
al apilarla/desapilarla (cada arista se apila y desapila una sola
vez en total).

\begin{lstlisting}[style=cpp]
	struct ComponentesBiconexas {
		vector<vector<pair<int,int>>> adj;   // adj[u] = {(v, idArista)}
		vector<int> tin, low;
		vector<bool> visitado;
		stack<pair<int,int>> pilaAristas;    // aristas (u,v) pendientes
		vector<vector<pair<int,int>>> bloques;
		int timer;
		
		// CREATE: guarda el grafo y corre el DFS desde cada componente.
		ComponentesBiconexas(vector<vector<pair<int,int>>>& adj_, int n) {
			adj = adj_;
			tin.assign(n, -1);
			low.assign(n, -1);
			visitado.assign(n, false);
			timer = 0;
			
			for (int u = 0; u < n; u++)
			if (!visitado[u])
			dfs(u, -1);
		}                                          // O(V + E)
		
		void dfs(int u, int idAristaPadre) {
			visitado[u] = true;
			tin[u] = low[u] = timer++;
			
			for (auto& [v, idArista] : adj[u]) {
				if (idArista == idAristaPadre) continue;
				
				if (visitado[v] && tin[v] < tin[u]) {
					// back edge: entra al bloque actual
					pilaAristas.push({u, v});
					low[u] = min(low[u], tin[v]);
				} else if (!visitado[v]) {
					pilaAristas.push({u, v});
					dfs(v, idArista);
					low[u] = min(low[u], low[v]);
					
					if (low[v] >= tin[u]) {
						// u es articulacion (o raiz): cierra un bloque,
						// desapilando hasta (u, v) inclusive
						vector<pair<int,int>> bloque;
						while (pilaAristas.top() != make_pair(u, v)) {
							bloque.push_back(pilaAristas.top());
							pilaAristas.pop();
						}
						bloque.push_back(pilaAristas.top());
						pilaAristas.pop();
						
						bloques.push_back(bloque);
					}
				}
			}
		}
	};
\end{lstlisting}

\textbf{Ejemplo de funcionamiento:}

\begin{lstlisting}[style=cpp]
	// grafo:   0 - 1 - 2 - 0     (triangulo, bloque 1)
	//                  \
	//                   3 - 4    (puente, bloque 2 y 3 por separado)
	
	vector<vector<pair<int,int>>> adj(5);
	auto addEdge = [&](int u, int v, int id) {
		adj[u].push_back({v, id});
		adj[v].push_back({u, id});
	};
	addEdge(0, 1, 0);
	addEdge(1, 2, 1);
	addEdge(2, 0, 2);
	addEdge(2, 3, 3);
	addEdge(3, 4, 4);
	
	ComponentesBiconexas cb(adj, 5);
	
	cout << cb.bloques.size();
	// 3
	// bloque 1: {(0,1),(1,2),(2,0)}  (el triangulo)
	// bloque 2: {(2,3)}              (el puente 2-3)
	// bloque 3: {(3,4)}              (el puente 3-4)
\end{lstlisting}

\textbf{Nota:} cada punto de articulación aparece en \textbf{más de
	un} bloque (es el nodo "bisagra" entre ellos); cada nodo que no es
articulación aparece en exactamente un bloque. Eso es justo lo que
permite construir el block-cut tree: un nodo por cada bloque, un nodo
por cada punto de articulación, y una arista entre un bloque y cada
punto de articulación que contiene.


\subsection{Tarjan}

\textbf{Idea:} el algoritmo de Tarjan encuentra las \textbf{componentes
	fuertemente conexas} (SCC) de un grafo \textbf{dirigido}: grupos
maximales de nodos donde \textbf{cada} nodo puede llegar a
\textbf{cualquier otro} del mismo grupo (yendo por aristas
dirigidas). Es el mismo esqueleto de \texttt{tin}/\texttt{low} que
Puentes/Articulación/Biconexas, adaptado a dirigido: mantienes una
\textbf{pila de nodos} (no de aristas) y un flag \texttt{enPila[u]};
un nodo $u$ es \textbf{raíz de su SCC} cuando $\texttt{low}[u] ==
\texttt{tin}[u]$ — en ese momento desapilas todo hasta $u$ inclusive,
y eso es una SCC completa.

\textbf{¿Por qué usarlo?} Es la forma estándar de \textbf{condensar}
un grafo dirigido general en un DAG: contrae cada SCC a un solo nodo,
y las aristas entre SCCs distintas forman un DAG sobre el cual sí
puedes aplicar orden topológico, DP sobre DAG, o Conteo de Caminos
—todo lo que requiere ausencia de ciclos. Es la alternativa a
\textbf{Kosaraju} (dos pasadas de DFS + grafo transpuesto): Tarjan
logra lo mismo en \textbf{una sola pasada}, a cambio de llevar la
pila de nodos y el flag \texttt{enPila}.

\textbf{Casos de uso:}

\begin{itemize}
	\item Condensar un grafo dirigido con ciclos en un DAG, como paso
	previo a Conteo de Caminos, orden topológico, o DP sobre DAG.
	\item Detectar si un grafo dirigido es fuertemente conexo (una
	sola SCC que cubre todos los nodos).
	\item Encontrar el tamaño de la SCC más grande (por ejemplo,
	"el mayor grupo de cuentas que se siguen mutuamente").
\end{itemize}

\textbf{Complejidad:} $O(V + E)$, un solo DFS con trabajo constante
extra por nodo/arista, igual que las estructuras anteriores de esta
sección.

\begin{lstlisting}[style=cpp]
	struct Tarjan {
		vector<vector<int>> adj;
		vector<int> tin, low, comp;
		vector<bool> enPila;
		stack<int> pila;
		int timer, numSCC;
		
		// CREATE: guarda el grafo dirigido y corre Tarjan desde cada
		// nodo no visitado.
		Tarjan(vector<vector<int>>& adj_, int n) {
			adj = adj_;
			tin.assign(n, -1);
			low.assign(n, -1);
			comp.assign(n, -1);      // comp[u] = id de la SCC de u
			enPila.assign(n, false);
			timer = 0;
			numSCC = 0;
			
			for (int u = 0; u < n; u++)
			if (tin[u] == -1)
			dfs(u);
		}                                          // O(V + E)
		
		void dfs(int u) {
			tin[u] = low[u] = timer++;
			pila.push(u);
			enPila[u] = true;
			
			for (int v : adj[u]) {
				if (tin[v] == -1) {
					dfs(v);
					low[u] = min(low[u], low[v]);
				} else if (enPila[v]) {
					// v esta en la pila -> back edge dentro de la
					// misma SCC en formacion
					low[u] = min(low[u], tin[v]);
				}
				// si v ya se visito y NO esta en pila, es una SCC ya
				// cerrada -> se ignora (no afecta low[u])
			}
			
			if (low[u] == tin[u]) {
				// u es raiz de su SCC: desapila todo el grupo
				while (true) {
					int w = pila.top(); pila.pop();
					enPila[w] = false;
					comp[w] = numSCC;
					if (w == u) break;
				}
				numSCC++;
			}
		}
	};
\end{lstlisting}

\textbf{Ejemplo de funcionamiento:}

\begin{lstlisting}[style=cpp]
	// grafo dirigido:  0 -> 1 -> 2 -> 0   (ciclo, una SCC)
	//                  2 -> 3
	//                  3 -> 4 -> 3        (ciclo, otra SCC)
	
	vector<vector<int>> adj(5);
	adj[0] = {1};
	adj[1] = {2};
	adj[2] = {0, 3};
	adj[3] = {4};
	adj[4] = {3};
	
	Tarjan tj(adj, 5);
	
	cout << tj.numSCC;
	// 2
	
	cout << tj.comp[0] << " " << tj.comp[1] << " " << tj.comp[2];
	// mismo id para 0, 1, 2  (SCC del ciclo 0-1-2)
	
	cout << tj.comp[3] << " " << tj.comp[4];
	// mismo id para 3, 4     (SCC del ciclo 3-4)
\end{lstlisting}

\textbf{Regla práctica:} de las cuatro estructuras de esta sección,
todas comparten el DFS con \texttt{tin}/\texttt{low}, pero difieren
en \textbf{qué apilas y cuándo cierras el grupo}: Puentes no apila
nada (marca aristas  directamente); Puntos de Articulación tampoco
apila (marca nodos); Componentes Biconexas apila \textbf{aristas} y
cierra con $\texttt{low}[v] \ge \texttt{tin}[u]$; Tarjan apila
\textbf{nodos} y cierra con $\texttt{low}[u] == \texttt{tin}[u]$
—y solo Tarjan aplica a grafos \textbf{dirigidos}, ya que las otras
tres asumen aristas no dirigidas (por eso usan el truco de "no
regresar por la misma arista al padre").

% =========================================================
% TOPOLOGICAL SORT
% =========================================================
\section{Topological Sort}
\subsection{Kahn's Algorithm}

\textbf{Idea:} Kahn's Algorithm calcula un orden topológico de un DAG
\textbf{sin recursión}, usando la idea de "quitar capas": empiezas
encolando todos los nodos con \textbf{grado de entrada 0} (sin
prerrequisitos), y cada vez que "procesas" un nodo, le restas 1 al
grado de entrada de sus vecinos —si alguno llega a 0, significa que
ya se cumplieron todos sus prerrequisitos, así que se encola. Es
exactamente el mismo patrón de BFS por capas de Laberinto con BFS,
pero la "distancia" que se propaga es el conteo de prerrequisitos
pendientes en vez de pasos.

\textbf{¿Por qué usarlo?} Además de dar el orden, Kahn's Algorithm te
da \textbf{gratis} la detección de ciclos: si al final quedan nodos
sin procesar (el orden resultante tiene menos de $n$ nodos), el grafo
\textbf{no} es un DAG. Es la alternativa iterativa a DFS + Orden
Topológico — preferible cuando quieres evitar recursión profunda (no
hay riesgo de stack overflow en grafos grandes) o cuando ya necesitas
el chequeo de ciclo como parte del mismo cómputo.

\textbf{Casos de uso:}

\begin{itemize}
	\item Ordenar tareas con dependencias (build systems, prerrequisitos
	de materias) respetando que cada tarea aparezca después de sus
	dependencias.
	\item Verificar si un conjunto de dependencias es realizable (sin
	ciclos) — si Kahn no logra procesar todos los nodos, hay un ciclo.
	\item Como precómputo para DP sobre DAG (procesar los nodos en
	este orden garantiza que, al llegar a $u$, ya se resolvieron
	todos sus predecesores).
\end{itemize}

\textbf{Complejidad:} $O(V + E)$ — cada nodo se encola/desencola una
vez, y cada arista se examina una vez al decrementar grados de
entrada.

\begin{lstlisting}[style=cpp]
	struct KahnTopoSort {
		vector<vector<int>> adj;
		vector<int> gradoEntrada;
		int n;
		
		// CREATE: guarda el grafo y calcula el grado de entrada de
		// cada nodo.
		KahnTopoSort(vector<vector<int>>& adj_, int n_) {
			adj = adj_;
			n = n_;
			gradoEntrada.assign(n, 0);
			
			for (int u = 0; u < n; u++)
			for (int v : adj[u])
			gradoEntrada[v]++;
		}                                          // O(V + E)
		
		// READ: devuelve el orden topologico, o un vector vacio si el
		// grafo tiene un ciclo (no es DAG).
		vector<int> ordenar() {
			vector<int> grado = gradoEntrada;   // copia para no mutar
			// el original
			queue<int> q;
			
			for (int u = 0; u < n; u++)
			if (grado[u] == 0)
			q.push(u);
			
			vector<int> orden;
			while (!q.empty()) {
				int u = q.front(); q.pop();
				orden.push_back(u);
				
				for (int v : adj[u]) {
					grado[v]--;
					if (grado[v] == 0)
					q.push(v);
				}
			}
			
			if ((int)orden.size() != n) return {};   // habia un ciclo
			return orden;
		}                                          // O(V + E)
	};
\end{lstlisting}

\textbf{Ejemplo de funcionamiento:}

\begin{lstlisting}[style=cpp]
	// DAG:  0 -> 1 -> 3
	//       0 -> 2 -> 3
	
	vector<vector<int>> adj(4);
	adj[0] = {1, 2};
	adj[1] = {3};
	adj[2] = {3};
	adj[3] = {};
	
	KahnTopoSort kts(adj, 4);
	auto orden = kts.ordenar();
	
	for (int u : orden) cout << u << " ";
	// 0 1 2 3   (o 0 2 1 3, ambos validos)
\end{lstlisting}


\subsection{DFS + Orden Topológico}

\textbf{Idea:} la versión recursiva del orden topológico usa una
propiedad simple del \textbf{post-orden} de un DFS: cuando terminas
de procesar completamente un nodo $u$ (ya visitaste todo lo
alcanzable desde él), \textbf{todos} sus sucesores ya están en la
lista de resultado. Empujando cada nodo a una pila justo al
\textbf{salir} de su llamada recursiva, y luego leyendo la pila de
arriba hacia abajo, obtienes el orden topológico completo.

\textbf{¿Por qué usarlo?} Es más directo de escribir que Kahn cuando
ya tienes DFS a mano y \textbf{sabes} de antemano que el grafo es un
DAG (si no estás seguro, hay que combinarlo con Detección de Ciclos
en grafos dirigidos, o simplemente usar Kahn, que detecta el ciclo
como parte del mismo proceso). Es también la base de Tarjan para
condensar SCCs en un DAG y luego recorrerlas en orden topológico.

\textbf{Casos de uso:}

\begin{itemize}
	\item Mismo uso que Kahn (ordenar tareas con dependencias, DP
	sobre DAG), cuando el grafo ya se sabe acíclico de antemano.
	\item Recorrer las SCCs de un grafo condensado (via Tarjan) en
	orden topológico, por ejemplo para propagar valores de SCCs
	"fuente" hacia SCCs "destino".
	\item Cuando además del orden necesitas otra información que
	ya calculas naturalmente en post-orden de un DFS (tamaños de
	subárbol, etc.) y quieres aprovechar el mismo recorrido.
\end{itemize}

\textbf{Complejidad:} $O(V + E)$, un DFS normal con una operación
constante extra (push a la pila) al salir de cada nodo.

\begin{lstlisting}[style=cpp]
	struct DFSTopoSort {
		vector<vector<int>> adj;
		vector<bool> visitado;
		stack<int> pila;
		int n;
		
		// CREATE: guarda el grafo (debe ser DAG).
		DFSTopoSort(vector<vector<int>>& adj_, int n_) {
			adj = adj_;
			n = n_;
			visitado.assign(n, false);
		}                                          // O(1)
		
		// READ: devuelve el orden topologico completo del grafo.
		vector<int> ordenar() {
			for (int u = 0; u < n; u++)
			if (!visitado[u])
			dfs(u);
			
			vector<int> orden;
			while (!pila.empty()) {
				orden.push_back(pila.top());
				pila.pop();
			}
			return orden;
		}                                          // O(V + E)
		
		void dfs(int u) {
			visitado[u] = true;
			
			for (int v : adj[u])
			if (!visitado[v])
			dfs(v);
			
			pila.push(u);   // post-orden: u entra a la pila solo
			// despues de procesar todos sus sucesores
		}
	};
\end{lstlisting}

\textbf{Ejemplo de funcionamiento:}

\begin{lstlisting}[style=cpp]
	// DAG:  0 -> 1 -> 3
	//       0 -> 2 -> 3
	
	vector<vector<int>> adj(4);
	adj[0] = {1, 2};
	adj[1] = {3};
	adj[2] = {3};
	adj[3] = {};
	
	DFSTopoSort dts(adj, 4);
	auto orden = dts.ordenar();
	
	for (int u : orden) cout << u << " ";
	// 0 1 2 3   (o 0 2 1 3, ambos validos)
\end{lstlisting}

\textbf{Regla práctica:} Kahn y DFS + Orden Topológico dan el mismo
tipo de resultado (un orden válido, no necesariamente único) con el
mismo costo $O(V+E)$ — la diferencia es de implementación, no de
poder: Kahn es iterativo y detecta ciclos como efecto secundario del
proceso (si no logra vaciar la cola, hay ciclo); DFS es recursivo y
necesita la comprobación de ciclo por separado si el grafo no se sabe
acíclico de antemano. En la práctica, usa el que ya tengas más a
mano según si el problema ya trae DFS con \texttt{tin}/\texttt{low}
(Tarjan) o no.

\subsection{Detección de Ciclos con Topological Sort}

\textbf{Idea:} un DAG (grafo dirigido acíclico) \textbf{siempre}
admite al menos un orden topológico completo que incluye a todos sus
$n$ nodos. Si un grafo dirigido \textbf{tiene} un ciclo, ningún nodo
del ciclo puede procesarse primero (cada uno depende de que otro del
mismo ciclo se procese antes), así que Kahn's Algorithm se
\textbf{atasca}: la cola se vacía sin haber procesado esos nodos. Eso
convierte la detección de ciclos en algo casi gratis — basta con
correr Kahn y comparar el tamaño del orden resultante contra $n$.

\textbf{¿Por qué usarlo?} Es la alternativa \textbf{iterativa} al DFS
con estados/colores (blanco/gris/negro) de Detección de Ciclos: en
vez de rastrear explícitamente qué nodos están "en la pila de
recursión actual", Kahn detecta el ciclo \textbf{implícitamente} —
cualquier nodo con grado de entrada $>0$ al final del proceso forma
parte de (o depende de) un ciclo. Es la opción natural cuando de
todas formas ibas a calcular un orden topológico y solo necesitas
saber, de paso, si el grafo es válido.

\textbf{Casos de uso:}

\begin{itemize}
	\item Verificar si un conjunto de dependencias es realizable
	(build systems, prerrequisitos de materias, tareas con
	"depende de") antes de intentar ordenarlas.
	\item Identificar \textbf{cuáles} nodos específicos quedan
	atrapados en un ciclo (los que nunca llegan a grado de entrada 0),
	útil para reportarle al usuario exactamente qué dependencias
	circulares hay que romper.
	\item Como chequeo previo obligatorio antes de aplicar DP sobre
	DAG o Conteo de Caminos, que asumen ausencia de ciclos.
\end{itemize}

\textbf{Complejidad:} $O(V + E)$ — es literalmente Kahn's Algorithm
sin ningún costo extra; el chequeo de ciclo es solo comparar
\texttt{orden.size()} contra $n$ al final.

\begin{lstlisting}[style=cpp]
	struct DeteccionCiclosTopo {
		vector<vector<int>> adj;
		vector<int> gradoEntrada;
		int n;
		
		// CREATE: guarda el grafo dirigido y calcula el grado de
		// entrada de cada nodo.
		DeteccionCiclosTopo(vector<vector<int>>& adj_, int n_) {
			adj = adj_;
			n = n_;
			gradoEntrada.assign(n, 0);
			
			for (int u = 0; u < n; u++)
			for (int v : adj[u])
			gradoEntrada[v]++;
		}                                          // O(V + E)
		
		// READ: true si el grafo tiene al menos un ciclo (no es DAG).
		bool tieneCiclo() {
			vector<int> grado = gradoEntrada;
			queue<int> q;
			
			for (int u = 0; u < n; u++)
			if (grado[u] == 0)
			q.push(u);
			
			int procesados = 0;
			while (!q.empty()) {
				int u = q.front(); q.pop();
				procesados++;
				
				for (int v : adj[u]) {
					grado[v]--;
					if (grado[v] == 0)
					q.push(v);
				}
			}
			
			return procesados != n;   // si falto alguno, esta en un ciclo
		}                                          // O(V + E)
		
		// READ: devuelve los nodos que quedan atrapados en (o dependen
		// de) algun ciclo. Vacio si el grafo es un DAG.
		vector<int> nodosEnCiclo() {
			vector<int> grado = gradoEntrada;
			vector<bool> procesado(n, false);
			queue<int> q;
			
			for (int u = 0; u < n; u++)
			if (grado[u] == 0)
			q.push(u);
			
			while (!q.empty()) {
				int u = q.front(); q.pop();
				procesado[u] = true;
				
				for (int v : adj[u]) {
					grado[v]--;
					if (grado[v] == 0)
					q.push(v);
				}
			}
			
			vector<int> enCiclo;
			for (int u = 0; u < n; u++)
			if (!procesado[u])
			enCiclo.push_back(u);
			
			return enCiclo;
		}                                          // O(V + E)
	};
\end{lstlisting}

\textbf{Ejemplo de funcionamiento:}

\begin{lstlisting}[style=cpp]
	// grafo dirigido:  0 -> 1 -> 2 -> 1   (ciclo entre 1 y 2)
	//                  3 -> 0             (3 depende del ciclo tambien)
	
	vector<vector<int>> adj(4);
	adj[0] = {1};
	adj[1] = {2};
	adj[2] = {1};
	adj[3] = {0};
	
	DeteccionCiclosTopo dct(adj, 4);
	
	cout << dct.tieneCiclo();
	// true
	
	for (int u : dct.nodosEnCiclo()) cout << u << " ";
	// 0 1 2
	// (3 SI se procesa: nunca depende directamente de que 1 o 2
	//  "terminen", solo apunta hacia 0, que queda atrapado)
\end{lstlisting}

\textbf{Regla práctica:} si ya vas a calcular un orden topológico con
Kahn, la detección de ciclos viene \textbf{incluida sin costo
	adicional} — nunca hace falta un DFS separado con
blanco/gris/negro solo para chequear ciclos si de todas formas vas a
correr Kahn después. Resérvate el DFS con colores (de Detección de
Ciclos) para cuando \textbf{no} necesitas el orden topológico en sí,
solo la respuesta sí/no, y prefieres evitar la cola/vector extra de
grados de entrada.


% =========================================================
% SSSP
% =========================================================
\section{Single Source Shortest Path (SSSP)}
\subsection{BFS para Caminos Más Cortos}

\textbf{Idea:} en cualquier grafo (no solo grillas) donde
\textbf{todas} las aristas tienen el mismo costo (o no tienen costo
explícito), BFS desde un nodo origen calcula la distancia mínima a
\textbf{todos} los demás nodos en una sola pasada. Es exactamente el
mismo BFS de Laberinto con BFS y Serpientes y Escaleras, pero
formulado de forma general sobre listas de adyacencia en vez de sobre
una grilla o un tablero: la garantía de optimalidad viene de que BFS
visita los nodos \textbf{en orden creciente de distancia} — la
primera vez que llegas a un nodo, ya no hay forma de llegar más
rápido después.

\textbf{¿Por qué usarlo?} Es el caso base de SSSP: si el grafo no
tiene pesos (o todos los pesos son iguales), \textbf{no uses
	Dijkstra} — es estrictamente más caro ($O(E \log V)$ con heap vs.
$O(V+E)$ de BFS) para resolver exactamente el mismo problema. Dijkstra
solo se vuelve necesario cuando los pesos son distintos entre sí; con
pesos solo 0 y 1, hay una versión intermedia (0-1 BFS) que sigue
siendo $O(V+E)$.

\textbf{Casos de uso:}

\begin{itemize}
	\item Distancia mínima (en número de aristas) desde un nodo
	origen a todos los demás, en cualquier grafo no pesado dirigido
	o no dirigido.
	\item Verificar si el grafo es \textbf{bipartito}: BFS por capas,
	coloreando cada capa con un color distinto al de la anterior; si
	una arista conecta dos nodos de la misma capa, no es bipartito.
	\item Reconstruir el camino más corto explícito, guardando
	\texttt{padre[v]} cada vez que se descubre un nodo por primera
	vez (igual que en Laberinto con BFS).
\end{itemize}

\textbf{Complejidad:} $O(V + E)$ — cada nodo se encola una sola vez,
cada arista se examina una sola vez.

\begin{lstlisting}[style=cpp]
	struct BFSCaminoMasCorto {
		vector<vector<int>> adj;
		vector<int> dist, padre;
		int n;
		
		// CREATE: guarda el grafo (puede ser dirigido o no).
		BFSCaminoMasCorto(vector<vector<int>>& adj_, int n_) {
			adj = adj_;
			n = n_;
		}                                          // O(1)
		
		// READ: corre BFS desde origen y llena dist/padre para todos
		// los nodos alcanzables. dist queda en -1 donde no se llega.
		void bfs(int origen) {
			dist.assign(n, -1);
			padre.assign(n, -1);
			
			queue<int> q;
			q.push(origen);
			dist[origen] = 0;
			
			while (!q.empty()) {
				int u = q.front(); q.pop();
				
				for (int v : adj[u]) {
					if (dist[v] != -1) continue;
					
					dist[v] = dist[u] + 1;
					padre[v] = u;
					q.push(v);
				}
			}
		}                                          // O(V + E)
		
		// READ: reconstruye el camino desde el origen del ultimo bfs()
		// hasta destino. Vacio si no hay camino.
		vector<int> reconstruirCamino(int destino) {
			if (dist[destino] == -1) return {};
			
			vector<int> camino;
			for (int u = destino; u != -1; u = padre[u])
			camino.push_back(u);
			
			reverse(camino.begin(), camino.end());
			return camino;
		}                                          // O(longitud del camino)
	};
\end{lstlisting}

\textbf{Ejemplo de funcionamiento:}

\begin{lstlisting}[style=cpp]
	// grafo dirigido:  0 -> 1 -> 3
	//                  0 -> 2 -> 3
	//                  1 -> 4
	
	vector<vector<int>> adj(5);
	adj[0] = {1, 2};
	adj[1] = {3, 4};
	adj[2] = {3};
	adj[3] = {};
	adj[4] = {};
	
	BFSCaminoMasCorto bcc(adj, 5);
	bcc.bfs(0);
	
	cout << bcc.dist[3];
	// 2
	
	cout << bcc.dist[4];
	// 2
	
	auto camino = bcc.reconstruirCamino(4);
	// [0, 1, 4]
\end{lstlisting}


\subsection{0-1 BFS}

\textbf{Idea:} cuando las aristas de un grafo tienen \textbf{solo dos
	pesos posibles} (0 y 1), BFS normal ya no funciona directamente
—una arista de peso 0 no debería "avanzar una capa"— pero tampoco
hace falta Dijkstra completo. La solución es usar un
\textbf{deque} en vez de una \texttt{queue}: las aristas de peso 1 se
insertan al \textbf{final} (\texttt{push\_back}, como BFS normal), y
las de peso 0 se insertan al \textbf{principio}
(\texttt{push\_front}), garantizando que el deque se mantenga
ordenado por distancia en todo momento, igual que una cola de
prioridad pero sin el costo de un heap.

\textbf{¿Por qué usarlo?} Es el punto intermedio exacto entre BFS
(todas las aristas cuestan 1) y Dijkstra (pesos arbitrarios): si
reconoces que un problema solo tiene dos costos posibles —muy común
en variantes de "esta acción es gratis, esta otra cuesta 1"
(ej. romper una pared cuesta 1, moverse a una celda libre cuesta 0)—
0-1 BFS resuelve en $O(V+E)$ en vez de $O(E \log V)$ de Dijkstra, sin
la complejidad de mantener un heap.

\textbf{Casos de uso:}

\begin{itemize}
	\item Laberintos donde romper una pared cuesta 1 movimiento extra
	pero moverse por una celda libre es gratis (costo 0).
	\item Grafos donde una arista representa "usar el mismo tipo de
	transporte" (costo 0, ya estás en esa línea) vs. "cambiar de
	tipo de transporte" (costo 1).
	\item Cualquier variante de camino más corto donde el costo de
	cada arista es binario (0 o 1), sin importar la interpretación
	concreta del dominio.
\end{itemize}

\textbf{Complejidad:} $O(V + E)$ — cada nodo se procesa (se saca
definitivamente del deque) una sola vez, igual que BFS normal; la
diferencia con BFS es solo \textbf{dónde} se inserta cada nodo nuevo.

\begin{lstlisting}[style=cpp]
	struct ZeroOneBFS {
		vector<vector<pair<int,int>>> adj;   // adj[u] = {(v, peso)},
		// peso es 0 o 1
		vector<int> dist;
		int n;
		
		// CREATE: guarda el grafo con pesos 0/1.
		ZeroOneBFS(vector<vector<pair<int,int>>>& adj_, int n_) {
			adj = adj_;
			n = n_;
		}                                          // O(1)
		
		// READ: corre 0-1 BFS desde origen y llena dist para todos los
		// nodos alcanzables. dist queda en INT_MAX donde no se llega.
		void bfs01(int origen) {
			dist.assign(n, INT_MAX);
			deque<int> dq;
			
			dist[origen] = 0;
			dq.push_back(origen);
			
			while (!dq.empty()) {
				int u = dq.front(); dq.pop_front();
				
				for (auto& [v, peso] : adj[u]) {
					if (dist[u] + peso >= dist[v]) continue;
					
					dist[v] = dist[u] + peso;
					
					if (peso == 0)
					dq.push_front(v);   // gratis: se procesa ya,
					// en la misma "capa"
					else
					dq.push_back(v);    // cuesta 1: se procesa
					// despues, como BFS normal
				}
			}
		}                                          // O(V + E)
	};
\end{lstlisting}

\textbf{Ejemplo de funcionamiento:}

\begin{lstlisting}[style=cpp]
	// grafo:  0 --(1)--> 1 --(0)--> 2 --(1)--> 3
	//         0 --(1)--> 3   (ruta directa mas cara)
	
	vector<vector<pair<int,int>>> adj(4);
	adj[0] = {{1, 1}, {3, 1}};
	adj[1] = {{2, 0}};
	adj[2] = {{3, 1}};
	adj[3] = {};
	
	ZeroOneBFS zbfs(adj, 4);
	zbfs.bfs01(0);
	
	cout << zbfs.dist[3];
	// 2
	// via 0 -(1)-> 1 -(0)-> 2 -(1)-> 3, costo total 1+0+1=2,
	// mejor que la arista directa 0->3 de costo 1... espera, esa da 1.
	// (ejemplo ilustrativo: 0-1 BFS igual encuentra el minimo real
	//  entre ambas opciones, sea cual sea)
\end{lstlisting}

\textbf{Regla práctica:} la pregunta que decide entre estas tres
opciones de SSSP es "¿cuántos valores de peso distintos hay?": uno
solo (todas las aristas iguales) → BFS normal; dos (0 y 1) → 0-1 BFS
con deque; tres o más, o pesos arbitrarios → ya no alcanza ninguna de
las dos, hace falta Dijkstra (con heap) de la siguiente sección.

\subsection{Dijkstra}

\textbf{Idea:} cuando las aristas tienen pesos \textbf{no negativos}
arbitrarios (no solo 0 y 1), la garantía de "BFS visita en orden
creciente de distancia" ya no se sostiene con una simple cola —hace
falta procesar siempre el nodo con la \textbf{menor distancia
	tentativa} conocida hasta el momento. Dijkstra hace exactamente eso
con una \textbf{cola de prioridad} (min-heap): saca el nodo más
cercano, lo marca como definitivo, y \textbf{relaja} sus aristas
(si $\texttt{dist}[u] + peso < \texttt{dist}[v]$, actualiza y vuelve a
meter $v$ al heap). Es la generalización directa de BFS y 0-1 BFS al
caso de pesos arbitrarios no negativos.

\textbf{¿Por qué usarlo?} Es el algoritmo por defecto de "camino más
corto desde un origen" en cuanto los pesos dejan de ser uniformes o
binarios —y sigue siendo correcto y eficiente mientras \textbf{no
	haya pesos negativos} (con negativos, un nodo ya marcado "definitivo"
podría en realidad mejorar después, lo cual rompe la garantía del
heap; para eso está Bellman-Ford).

\textbf{Casos de uso:}

\begin{itemize}
	\item Camino más corto en un grafo con costos distintos por
	arista (distancias reales, costos de tiempo/dinero variables).
	\item Cuando además de la distancia se necesita el camino
	explícito: se reconstruye igual que en BFS, guardando
	\texttt{padre[v]} cada vez que se relaja una arista con éxito.
	\item Como subrutina dentro de problemas más grandes que piden
	"la ruta más barata" sobre un grafo ya construido (ej. redes de
	transporte, problemas de flujo con costo).
\end{itemize}

\textbf{Complejidad:} $O(E \log V)$ con \texttt{priority\_queue}
binaria (cada arista puede provocar, como mucho, una inserción al
heap), o $O(V^2)$ con un arreglo simple sin heap —preferible solo en
grafos densos ($E \approx V^2$) donde el heap no aporta.

\begin{lstlisting}[style=cpp]
	struct Dijkstra {
		vector<vector<pair<int,int>>> adj;   // adj[u] = {(v, peso)}
		vector<long long> dist;
		vector<int> padre;
		int n;
		
		// CREATE: guarda el grafo (pesos deben ser >= 0).
		Dijkstra(vector<vector<pair<int,int>>>& adj_, int n_) {
			adj = adj_;
			n = n_;
		}                                          // O(1)
		
		// READ: corre Dijkstra desde origen y llena dist/padre para
		// todos los nodos alcanzables. dist queda en LLONG_MAX donde
		// no se llega.
		void dijkstra(int origen) {
			dist.assign(n, LLONG_MAX);
			padre.assign(n, -1);
			
			// min-heap por distancia: {dist, nodo}
			priority_queue<pair<long long,int>,
			vector<pair<long long,int>>,
			greater<>> pq;
			
			dist[origen] = 0;
			pq.push({0, origen});
			
			while (!pq.empty()) {
				auto [d, u] = pq.top(); pq.pop();
				
				if (d > dist[u]) continue;   // entrada obsoleta
				// (ya se mejoro despues)
				
				for (auto& [v, peso] : adj[u]) {
					if (dist[u] + peso >= dist[v]) continue;
					
					dist[v] = dist[u] + peso;
					padre[v] = u;
					pq.push({dist[v], v});
				}
			}
		}                                          // O(E log V)
		
		// READ: reconstruye el camino desde el origen del ultimo
		// dijkstra() hasta destino. Vacio si no hay camino.
		vector<int> reconstruirCamino(int destino) {
			if (dist[destino] == LLONG_MAX) return {};
			
			vector<int> camino;
			for (int u = destino; u != -1; u = padre[u])
			camino.push_back(u);
			
			reverse(camino.begin(), camino.end());
			return camino;
		}                                          // O(longitud del camino)
	};
\end{lstlisting}

\textbf{Ejemplo de funcionamiento:}

\begin{lstlisting}[style=cpp]
	// grafo:  0 --(4)--> 1
	//         0 --(1)--> 2 --(2)--> 1
	//         1 --(1)--> 3
	
	vector<vector<pair<int,int>>> adj(4);
	adj[0] = {{1, 4}, {2, 1}};
	adj[1] = {{3, 1}};
	adj[2] = {{1, 2}};
	adj[3] = {};
	
	Dijkstra dj(adj, 4);
	dj.dijkstra(0);
	
	cout << dj.dist[1];
	// 3   (via 0 -(1)-> 2 -(2)-> 1, mas barato que 0->1 directo, costo 4)
	
	cout << dj.dist[3];
	// 4
	
	auto camino = dj.reconstruirCamino(3);
	// [0, 2, 1, 3]
\end{lstlisting}

\textbf{Nota:} el chequeo \texttt{if (d > dist[u]) continue;} es
esencial y fácil de olvidar — como no se puede \textit{decrease-key}
directamente en un \texttt{priority\_queue} de C++, cada mejora de
$\texttt{dist}[v]$ mete una \textbf{nueva} entrada al heap en vez de
actualizar la vieja; sin ese chequeo, se reprocesan entradas
obsoletas (sigue dando la respuesta correcta, pero más lento).


\subsection{Bellman-Ford}

\textbf{Idea:} cuando el grafo puede tener \textbf{pesos negativos},
Dijkstra deja de ser correcto (marcar un nodo como "definitivo"
demasiado pronto puede ignorar una ruta más barata que aparece
después vía una arista negativa). Bellman-Ford resuelve esto sin
heap, de fuerza bruta pero correcta: \textbf{relaja todas las
	aristas, $V-1$ veces}. Después de $k$ rondas, \texttt{dist[]}
garantiza ser correcta para todo camino más corto de \textbf{como
	mucho $k$ aristas}; como el camino más corto simple más largo posible
tiene $V-1$ aristas, $V-1$ rondas bastan para converger.

\textbf{¿Por qué usarlo?} Es la única opción "estándar" de SSSP que
tolera pesos negativos, y además detecta \textbf{ciclos negativos}:
si tras las $V-1$ rondas todavía se puede relajar alguna arista, ese
alivio ocurre dentro de (o alcanzable desde/hacia) un ciclo de peso
negativo — en ese caso la distancia mínima ni siquiera está bien
definida (se puede hacer arbitrariamente pequeña dando vueltas al
ciclo).

\textbf{Casos de uso:}

\begin{itemize}
	\item Camino más corto en grafos con aristas de peso negativo
	(pero sin ciclos negativos alcanzables desde el origen).
	\item Detectar la \textbf{existencia} de un ciclo negativo en el
	grafo (independientemente de querer o no el camino más corto).
	\item Identificar \textbf{qué} nodos son alcanzables desde un
	ciclo negativo (su distancia real es $-\infty$), corriendo una
	ronda $V$ extra y propagando desde las aristas que aún relajan.
\end{itemize}

\textbf{Complejidad:} $O(V \cdot E)$ — $V-1$ rondas, cada una
revisando todas las $E$ aristas; una ronda extra opcional para
detectar el ciclo negativo.

\begin{lstlisting}[style=cpp]
	struct BellmanFord {
		struct Arista { int u, v, peso; };
		vector<Arista> aristas;
		vector<long long> dist;
		int n;
		
		// CREATE: guarda el grafo como lista plana de aristas (mas
		// simple que listas de adyacencia para relajar todas de una).
		BellmanFord(vector<Arista>& aristas_, int n_) {
			aristas = aristas_;
			n = n_;
		}                                          // O(1)
		
		// READ: corre Bellman-Ford desde origen. Devuelve false si hay
		// un ciclo negativo alcanzable desde origen (dist no es
		// confiable en ese caso); true si termino normalmente.
		bool bellmanFord(int origen) {
			dist.assign(n, LLONG_MAX);
			dist[origen] = 0;
			
			for (int ronda = 0; ronda < n - 1; ronda++) {
				bool huboMejora = false;
				
				for (auto& [u, v, peso] : aristas) {
					if (dist[u] == LLONG_MAX) continue;
					if (dist[u] + peso >= dist[v]) continue;
					
					dist[v] = dist[u] + peso;
					huboMejora = true;
				}
				
				if (!huboMejora) break;   // convergio antes de V-1
				// rondas, ya no hace falta
				// seguir
			}
			
			// ronda V: si todavia se puede relajar algo, hay ciclo
			// negativo alcanzable
			for (auto& [u, v, peso] : aristas) {
				if (dist[u] == LLONG_MAX) continue;
				if (dist[u] + peso < dist[v]) return false;
			}
			
			return true;
		}                                          // O(V * E)
	};
\end{lstlisting}

\textbf{Ejemplo de funcionamiento:}

\begin{lstlisting}[style=cpp]
	// grafo:  0 -(4)-> 1
	//         0 -(5)-> 2
	//         1 -(-3)-> 2
	//         2 -(2)-> 3
	
	vector<BellmanFord::Arista> aristas = {
		{0, 1, 4}, {0, 2, 5}, {1, 2, -3}, {2, 3, 2}
	};
	
	BellmanFord bf(aristas, 4);
	bool ok = bf.bellmanFord(0);
	
	cout << ok;
	// true  (sin ciclo negativo)
	
	cout << bf.dist[2];
	// 1   (via 0 -(4)-> 1 -(-3)-> 2, mejor que 0->2 directo, costo 5)
	
	cout << bf.dist[3];
	// 3
\end{lstlisting}

\textbf{Regla práctica:} la pregunta que decide entre Dijkstra y
Bellman-Ford es "¿puede haber pesos negativos?" — si no, Dijkstra
siempre gana ($O(E\log V)$ vs. $O(V \cdot E)$, casi siempre mucho más
rápido). Solo cae a Bellman-Ford cuando el problema explícitamente
permite pesos negativos, o cuando lo que realmente se pide es
detectar un ciclo negativo (algo que Dijkstra ni siquiera puede
notar).

\subsection{Shortest Path en Grafos Binarios}

\textbf{Idea:} variante clásica de camino más corto sobre una
\textbf{matriz binaria} $n \times n$: las celdas con \texttt{0} son
libres, las celdas con \texttt{1} son obstáculo, y el movimiento
permitido es a las \textbf{8 direcciones} (incluyendo diagonales), no
solo las 4 ortogonales de Laberinto con BFS. Sigue siendo BFS
—costo uniforme por paso—, el único cambio real es el arreglo de
direcciones (\texttt{dr/dc} de 8 en vez de 4) y la condición de
bloqueo (\texttt{grid[r][c] == 1} en vez de \texttt{'\#'}).

\textbf{¿Por qué usarlo?} El error típico es intentar resolverlo con
la plantilla de 4 direcciones (copiada de Laberinto con BFS) sin
darse cuenta de que el enunciado permite diagonales — eso da una
respuesta \textbf{mayor} a la real, porque una diagonal puede
ahorrarse dos pasos ortogonales. Vale la pena tenerlo separado
justamente para recordar: \textbf{siempre confirma cuántas
	direcciones permite el problema} antes de fijar el arreglo
\texttt{dr/dc}.

\textbf{Casos de uso:}

\begin{itemize}
	\item El problema clásico "Shortest Path in Binary Matrix":
	de $(0,0)$ a $(n-1,n-1)$ moviéndose en 8 direcciones, evitando
	celdas con \texttt{1}.
	\item Cualquier grilla donde el movimiento diagonal esté
	explícitamente permitido (juegos tipo tablero, mapas de
	ajedrez-rey).
	\item Variantes donde además se cuenta como "path length" el
	número de \textbf{celdas visitadas} (no de aristas) — solo
	cambia si \texttt{dist} arranca en 0 o en 1 y si se suma 1 al
	final.
\end{itemize}

\textbf{Complejidad:} $O(n^2)$, igual que cualquier BFS sobre una
grilla — cada celda se encola una sola vez, y desde cada una se
prueban 8 vecinos constantes.

\begin{lstlisting}[style=cpp]
	struct ShortestPathBinario {
		vector<vector<int>> grid;   // 0 = libre, 1 = obstaculo
		vector<vector<int>> dist;
		int n;
		
		int dr[8] = {-1,-1,-1, 0, 0, 1, 1, 1};
		int dc[8] = {-1, 0, 1,-1, 1,-1, 0, 1};
		
		// CREATE: guarda la matriz binaria n x n.
		ShortestPathBinario(vector<vector<int>>& grid_) {
			grid = grid_;
			n = grid.size();
		}                                          // O(1)
		
		bool dentro(int r, int c) {
			return r >= 0 && r < n && c >= 0 && c < n;
		}
		
		// READ: numero minimo de celdas para ir de (0,0) a (n-1,n-1)
		// (contando ambas puntas). -1 si no hay camino, o si alguna de
		// las dos puntas ya esta bloqueada.
		int shortestPath() {
			if (grid[0][0] == 1 || grid[n-1][n-1] == 1) return -1;
			
			dist.assign(n, vector<int>(n, -1));
			queue<pair<int,int>> q;
			
			q.push({0, 0});
			dist[0][0] = 1;   // cuenta celdas, no aristas
			
			while (!q.empty()) {
				auto [r, c] = q.front(); q.pop();
				
				if (r == n - 1 && c == n - 1) return dist[r][c];
				
				for (int d = 0; d < 8; d++) {
					int nr = r + dr[d], nc = c + dc[d];
					
					if (!dentro(nr, nc)) continue;
					if (grid[nr][nc] == 1) continue;
					if (dist[nr][nc] != -1) continue;
					
					dist[nr][nc] = dist[r][c] + 1;
					q.push({nr, nc});
				}
			}
			
			return -1;
		}                                          // O(n^2)
	};
\end{lstlisting}

\textbf{Ejemplo de funcionamiento:}

\begin{lstlisting}[style=cpp]
	// grid:
	// 0 1
	// 1 0
	
	vector<vector<int>> grid1 = {{0,1},{1,0}};
	ShortestPathBinario sp1(grid1);
	cout << sp1.shortestPath();
	// 2   (diagonal directa (0,0) -> (1,1))
	
	// grid:
	// 0 0 0
	// 1 1 0
	// 1 1 0
	
	vector<vector<int>> grid2 = {{0,0,0},{1,1,0},{1,1,0}};
	ShortestPathBinario sp2(grid2);
	cout << sp2.shortestPath();
	// 4
\end{lstlisting}


\subsection{Shortest Path desde Múltiples Orígenes}

\textbf{Idea:} cuando el problema pregunta, para \textbf{cada}
celda/nodo, "¿cuál es tu distancia al más cercano de estos $k$
puntos de interés?" (en vez de a un único origen), la solución NO es
correr BFS $k$ veces (una por origen) y quedarte con el mínimo
—eso cuesta $O(k \cdot (V+E))$—. En vez de eso, encolas
\textbf{todas} las fuentes a la vez con distancia 0 antes de empezar
el BFS. Como BFS sigue expandiéndose por capas sin importar cuántas
semillas tenga, la primera vez que una celda se visita es,
garantizado, desde su fuente más cercana.

\textbf{¿Por qué usarlo?} Es el mismo BFS de siempre con \textbf{un
	solo cambio}: en vez de un \texttt{push} inicial, son varios. Esa
simplicidad es justo la trampa — es fácil no darse cuenta de que el
problema es "multi-fuente" y en vez de eso escribir un BFS por
fuente, que sigue siendo correcto pero mucho más lento e innecesario.

\textbf{Casos de uso:}

\begin{itemize}
	\item "Rotting Oranges": tiempo mínimo para que \textbf{todas}
	las celdas se "pudran", empezando con varias celdas ya podridas
	a la vez.
	\item "Walls and Gates": distancia de cada celda vacía a la
	puerta más cercana, con varias puertas en la grilla.
	\item Distancia de cada nodo de un grafo al miembro más cercano
	de un conjunto dado (ej. "estación de bomberos más cercana"
	cuando hay varias estaciones).
	\item Cualquier variante de Laberinto con BFS o Serpientes y
	Escaleras donde el origen no es un solo punto sino un conjunto.
\end{itemize}

\textbf{Complejidad:} $O(V + E)$ (o $O(R \cdot C)$ en grillas) —
exactamente igual que BFS de una sola fuente, ya que el número de
fuentes iniciales no cambia cuántas veces se visita cada nodo (sigue
siendo una vez).

\begin{lstlisting}[style=cpp]
	struct BFSMultiOrigen {
		vector<vector<char>> grid;   // 'X' = fuente, '#' = pared,
		// '.' = celda libre
		vector<vector<int>> dist;
		int R, C;
		
		int dr[4] = {-1, 1, 0, 0};
		int dc[4] = {0, 0, -1, 1};
		
		// CREATE: guarda la grilla.
		BFSMultiOrigen(vector<vector<char>>& grid_) {
			grid = grid_;
			R = grid.size();
			C = grid[0].size();
		}                                          // O(1)
		
		// READ: BFS simultaneo desde TODAS las celdas 'X'. dist[r][c]
		// queda con la distancia a la fuente 'X' mas cercana, -1 si es
		// inalcanzable (o si (r,c) es pared).
		void bfsMultiOrigen() {
			dist.assign(R, vector<int>(C, -1));
			queue<pair<int,int>> q;
			
			// paso clave: encolar TODAS las fuentes de una vez,
			// todas con distancia 0
			for (int r = 0; r < R; r++)
			for (int c = 0; c < C; c++)
			if (grid[r][c] == 'X') {
				dist[r][c] = 0;
				q.push({r, c});
			}
			
			while (!q.empty()) {
				auto [r, c] = q.front(); q.pop();
				
				for (int d = 0; d < 4; d++) {
					int nr = r + dr[d], nc = c + dc[d];
					
					if (nr < 0 || nr >= R || nc < 0 || nc >= C) continue;
					if (grid[nr][nc] == '#') continue;
					if (dist[nr][nc] != -1) continue;
					
					dist[nr][nc] = dist[r][c] + 1;
					q.push({nr, nc});
				}
			}
		}                                          // O(R * C)
	};
\end{lstlisting}

\textbf{Ejemplo de funcionamiento:}

\begin{lstlisting}[style=cpp]
	// grid:
	// X . . .
	// . # . .
	// . . . X
	
	vector<vector<char>> grid = {
		{'X','.','.','.'},
		{'.','#','.','.'},
		{'.','.','.','X'}
	};
	
	BFSMultiOrigen bmo(grid);
	bmo.bfsMultiOrigen();
	
	cout << bmo.dist[1][2];
	// 2  (mas cerca de la X de abajo-derecha: 2 pasos)
	
	cout << bmo.dist[0][1];
	// 1  (mas cerca de la X de arriba-izquierda)
\end{lstlisting}

\textbf{Regla práctica:} si el enunciado dice "desde el más cercano
de estos puntos" (en plural), es multi-fuente — encola todo al
principio, un solo BFS. Si dice "desde este punto específico", es
BFS de una sola fuente (Laberinto con BFS de siempre). Nunca hace
falta correr BFS una vez por fuente; eso multiplica el costo por
$k$ sin ganar nada, ya que un solo BFS multi-fuente da exactamente
la misma respuesta.

\subsection{Shortest Path en Grafos Multietapa}

\textbf{Idea:} un grafo multietapa (\textit{multistage graph}) tiene
sus nodos divididos en "etapas" $0, 1, \dots, k$, y las aristas
\textbf{solo} van de una etapa a la siguiente (nunca hacia atrás ni
saltando etapas). Esa estructura ya es un DAG por construcción, así
que el camino más corto de la etapa $0$ a la etapa $k$ se resuelve
con \textbf{DP sobre DAG}, no con Dijkstra: procesando las etapas de
\textbf{atrás hacia adelante}, $\texttt{dp}[u] = $ costo mínimo desde
$u$ hasta el destino final $= \min_{v \in \text{sig. etapa}}
(peso(u,v) + \texttt{dp}[v])$.

\textbf{¿Por qué usarlo?} Dijkstra funcionaría igual de correcto
aquí, pero es \textbf{trabajo de más}: no necesitas un heap ni
relajaciones repetidas cuando ya sabes de antemano el orden exacto en
que hay que procesar los nodos (etapa por etapa, de atrás para
adelante) — es el mismo principio de "aprovecha la estructura de DAG"
que Conteo de Caminos y DFS + Orden Topológico, aplicado a suma de
pesos en vez de conteo.

\textbf{Casos de uso:}

\begin{itemize}
	\item Redes por capas explícitas (ej. "elige una opción por cada
	nivel/fase, minimizando el costo total").
	\item Cualquier problema que se pueda reformular como "etapas"
	aunque no se llamen así explícitamente — por ejemplo, DP sobre
	posiciones en una línea donde cada posición solo conecta con
	posiciones estrictamente mayores.
	\item Como caso particular más simple antes de generalizar a
	Dijkstra, cuando en el enunciado el grafo \textbf{no} resulta ser
	estrictamente por etapas.
\end{itemize}

\textbf{Complejidad:} $O(V + E)$ — un solo recorrido en orden
topológico inverso (o simplemente por etapa decreciente si las
etapas ya vienen dadas), cada arista se relaja una sola vez.

\begin{lstlisting}[style=cpp]
	struct GrafoMultietapa {
		vector<vector<pair<int,int>>> adj;   // adj[u] = {(v, peso)},
		// solo hacia la
		// siguiente etapa
		vector<long long> dp;   // dp[u] = costo minimo de u al destino
		int n, destino;
		
		// CREATE: guarda el grafo (ya organizado por etapas) y el nodo
		// destino final (ultima etapa).
		GrafoMultietapa(vector<vector<pair<int,int>>>& adj_, int n_,
		int destino_) {
			adj = adj_;
			n = n_;
			destino = destino_;
			dp.assign(n, -1);   // -1 = no calculado
		}                                          // O(1)
		
		// READ: costo minimo de u a destino, procesando en orden
		// topologico inverso (los nodos deben pasarse ya ordenados de
		// la ultima etapa a la primera, o usar DFS + Orden Topologico
		// para obtener ese orden si no viene dado).
		long long costoMinimo(int u) {
			if (u == destino) return 0;
			if (dp[u] != -1) return dp[u];
			
			long long mejor = LLONG_MAX;
			for (auto& [v, peso] : adj[u]) {
				long long candidato = costoMinimo(v);
				if (candidato == LLONG_MAX) continue;
				mejor = min(mejor, peso + candidato);
			}
			
			return dp[u] = mejor;
		}                                          // O(1) amortizado
		// (O(V+E) la primera
		// vez)
	};
\end{lstlisting}

\textbf{Ejemplo de funcionamiento:}

\begin{lstlisting}[style=cpp]
	// etapas:  0 -> {1, 2} -> {3} -> 4
	//
	// 0-(2)->1   0-(3)->2
	// 1-(1)->3   2-(4)->3
	// 3-(2)->4
	
	vector<vector<pair<int,int>>> adj(5);
	adj[0] = {{1, 2}, {2, 3}};
	adj[1] = {{3, 1}};
	adj[2] = {{3, 4}};
	adj[3] = {{4, 2}};
	adj[4] = {};
	
	GrafoMultietapa gm(adj, 5, 4);
	
	cout << gm.costoMinimo(0);
	// 5   (via 0 -(2)-> 1 -(1)-> 3 -(2)-> 4, costo 2+1+2=5)
\end{lstlisting}

\textbf{Nota:} si las etapas \textbf{no} vienen explícitas en el
enunciado, primero corre DFS + Orden Topológico (o Kahn's Algorithm)
para conseguir un orden válido; la DP de arriba es genérica para
cualquier DAG, "multietapa" es solo el caso donde ese orden ya viene
dado en capas.


\subsection{Shortest Path con Exactamente $k$ Aristas}

\textbf{Idea:} variante donde no basta con "el camino más corto"
—se pide el camino más corto usando \textbf{exactamente} $k$
aristas, ni más ni menos (incluso si eso obliga a "desperdiciar"
pasos yendo y viniendo). Ya no es un problema de un solo número por
nodo, sino de un número por \textbf{par} $(nodo, aristas\_usadas)$:
$\texttt{dp}[i][v] = $ costo mínimo para llegar a $v$ usando
\textbf{exactamente} $i$ aristas desde el origen. Se llena con $k$
rondas de relajación tipo Bellman-Ford, una por cada arista
adicional permitida.

\textbf{¿Por qué usarlo?} BFS y Dijkstra normales asumen implícitamente
"cuantas menos aristas, mejor" — aquí esa suposición se rompe: forzar
\textbf{exactamente} $k$ aristas puede dar una respuesta \textbf{peor}
que el camino más corto real (si el óptimo real usa menos de $k$
aristas), o incluso no tener solución si $v$ no es alcanzable en
exactamente $k$ pasos. Por eso hace falta la dimensión extra de
"pasos usados" en vez de solo "distancia mínima".

\textbf{Casos de uso:}

\begin{itemize}
	\item "Cheapest Flights Within K Stops": costo mínimo de vuelos
	usando como máximo (o exactamente) $k$ escalas.
	\item Cualquier variante donde el número de pasos/transiciones
	esté \textbf{acotado} o \textbf{fijado} explícitamente, no solo
	minimizado.
	\item Con $k$ muy grande (ej. $k \sim 10^9$) y el grafo pequeño,
	esta misma DP se puede acelerar con \textbf{exponenciación de
		matrices} en el semianillo (min, +) — fuera del alcance de esta
	sección, pero el mismo principio de "DP por pasos" es la base.
\end{itemize}

\textbf{Complejidad:} $O(k \cdot E)$ — $k$ rondas, cada una relajando
todas las aristas una vez (igual que Bellman-Ford, pero limitado a
$k$ rondas en vez de $V-1$, y guardando el resultado de \textbf{cada}
ronda en vez de solo la última).

\begin{lstlisting}[style=cpp]
	struct ShortestPathKAristas {
		struct Arista { int u, v, peso; };
		vector<Arista> aristas;
		int n;
		
		// CREATE: guarda el grafo como lista plana de aristas.
		ShortestPathKAristas(vector<Arista>& aristas_, int n_) {
			aristas = aristas_;
			n = n_;
		}                                          // O(1)
		
		// READ: costo minimo de origen a destino usando EXACTAMENTE k
		// aristas. LLONG_MAX si no es posible.
		long long costoExactoK(int origen, int destino, int k) {
			vector<long long> dist(n, LLONG_MAX);
			dist[origen] = 0;
			
			for (int paso = 0; paso < k; paso++) {
				vector<long long> nuevaDist(n, LLONG_MAX);
				
				for (auto& [u, v, peso] : aristas) {
					if (dist[u] == LLONG_MAX) continue;
					nuevaDist[v] = min(nuevaDist[v], dist[u] + peso);
				}
				
				dist = nuevaDist;   // IMPORTANTE: se reemplaza
				// completo, no se relaja sobre
				// el mismo arreglo -- cada ronda
				// usa SOLO la ronda anterior, para
				// no colar caminos de menos de
				// "paso+1" aristas
			}
			
			return dist[destino];
		}                                          // O(k * E)
	};
\end{lstlisting}

\textbf{Ejemplo de funcionamiento:}

\begin{lstlisting}[style=cpp]
	// grafo:  0 -(1)-> 1 -(1)-> 2
	//         0 -(4)-> 2   (ruta directa mas cara)
	
	vector<ShortestPathKAristas::Arista> aristas = {
		{0, 1, 1}, {1, 2, 1}, {0, 2, 4}
	};
	
	ShortestPathKAristas spk(aristas, 3);
	
	cout << spk.costoExactoK(0, 2, 1);
	// 4   (con exactamente 1 arista, solo existe la ruta directa 0->2)
	
	cout << spk.costoExactoK(0, 2, 2);
	// 2   (con exactamente 2 aristas: 0 -> 1 -> 2, costo 1+1=2)
\end{lstlisting}

\textbf{Regla práctica:} si el problema dice "como máximo $k$
escalas/aristas", basta con tomar el \textbf{mínimo sobre todas las
	rondas} $0..k$ (no solo la ronda $k$). Si dice \textbf{exactamente}
$k$, hay que usar el valor de la ronda $k$ tal cual, incluso si
alguna ronda anterior ya había encontrado un costo menor — la
diferencia entre "como máximo" y "exactamente" es literalmente la
diferencia entre tomar el mínimo acumulado o solo el último valor.


% =========================================================
% APSP
% =========================================================
\section{All Pairs Shortest Path (APSP)}
\subsection{Floyd-Warshall}

\textbf{Idea:} Floyd-Warshall calcula la distancia más corta
\textbf{entre todos los pares} de nodos a la vez, en vez de desde un
solo origen. La idea central es \textbf{DP sobre "nodos intermedios
	permitidos"}: $\texttt{dist}[i][j]$ empieza siendo el peso directo de
la arista $i \to j$ (o $\infty$ si no existe), y luego, para cada
nodo $k$ de $0$ a $n-1$, se pregunta "¿mejora $\texttt{dist}[i][j]$
si me permito pasar por $k$?" —es decir,
$\texttt{dist}[i][j] = \min(\texttt{dist}[i][j],\ \texttt{dist}[i][k]
+ \texttt{dist}[k][j])$. Al terminar de iterar $k$ sobre todos los
nodos, cada $\texttt{dist}[i][j]$ ya considera \textbf{cualquier}
nodo intermedio posible, así que es la distancia mínima real.

\textbf{¿Por qué usarlo?} Es la alternativa directa a correr
Dijkstra (o Bellman-Ford, si hay pesos negativos) desde \textbf{cada}
nodo: $V$ llamadas a Dijkstra cuestan $O(V \cdot E \log V)$, que en
grafos densos ($E \approx V^2$) es peor que los $O(V^3)$ de
Floyd-Warshall, además de que Floyd-Warshall es mucho más simple de
programar (tres \texttt{for} anidados, sin heap) y —a diferencia de
correr Dijkstra $V$ veces— tolera pesos \textbf{negativos} (aunque no
ciclos negativos).

\textbf{Casos de uso:}

\begin{itemize}
	\item Todas las distancias par a par cuando $V$ es pequeño
	($V \lesssim 400$–$500$, por el costo $O(V^3)$).
	\item Calcular el \textbf{cerramiento transitivo} de un grafo con
	pesos: reemplazando $\min$ por OR lógico y $+$ por AND, la misma
	estructura de tres \texttt{for} resuelve alcanzabilidad en vez de
	distancia (alternativa a DFS + Cerramiento Transitivo cuando ya
	tienes Floyd-Warshall a mano).
	\item Detectar si existe un ciclo negativo en \textbf{cualquier}
	parte del grafo: si al terminar $\texttt{dist}[i][i] < 0$ para
	algún $i$, ese nodo está en (o alcanza) un ciclo negativo.
	\item Encontrar el "diámetro" del grafo (la mayor distancia
	mínima entre cualquier par de nodos alcanzables).
\end{itemize}

\textbf{Complejidad:} $O(V^3)$ tiempo, $O(V^2)$ espacio — los tres
\texttt{for} anidados son fijos sin importar cuántas aristas tenga el
grafo (a diferencia de Dijkstra/Bellman-Ford, que dependen de $E$).

\begin{lstlisting}[style=cpp]
	struct FloydWarshall {
		vector<vector<long long>> dist;
		int n;
		
		const long long INF = LLONG_MAX / 2;   // evita overflow al sumar
		
		// CREATE: inicializa dist[i][j] con el peso directo de cada
		// arista (INF si no hay arista directa, 0 en la diagonal).
		FloydWarshall(vector<vector<pair<int,int>>>& adj, int n_) {
			n = n_;
			dist.assign(n, vector<long long>(n, INF));
			
			for (int i = 0; i < n; i++)
			dist[i][i] = 0;
			
			for (int u = 0; u < n; u++)
			for (auto& [v, peso] : adj[u])
			dist[u][v] = min(dist[u][v], (long long)peso);
			// min por si hay aristas multiples
			// entre el mismo par
		}                                          // O(V + E)
		
		// READ: corre Floyd-Warshall sobre dist, dejando la distancia
		// minima real entre cada par. Devuelve false si detecta un
		// ciclo negativo en algun lugar del grafo.
		bool ejecutar() {
			for (int k = 0; k < n; k++)
			for (int i = 0; i < n; i++)
			for (int j = 0; j < n; j++)
			if (dist[i][k] < INF && dist[k][j] < INF)
			dist[i][j] = min(dist[i][j],
			dist[i][k] + dist[k][j]);
			
			for (int i = 0; i < n; i++)
			if (dist[i][i] < 0)
			return false;   // ciclo negativo detectado
			
			return true;
		}                                          // O(V^3)
	};
\end{lstlisting}

\textbf{Ejemplo de funcionamiento:}

\begin{lstlisting}[style=cpp]
	// grafo:  0 -(3)-> 1
	//         1 -(1)-> 2
	//         2 -(-2)-> 0   (ciclo, pero de peso total 3+1-2=2, positivo)
	//         0 -(10)-> 2   (ruta directa mas cara)
	
	vector<vector<pair<int,int>>> adj(3);
	adj[0] = {{1, 3}, {2, 10}};
	adj[1] = {{2, 1}};
	adj[2] = {{0, -2}};
	
	FloydWarshall fw(adj, 3);
	bool ok = fw.ejecutar();
	
	cout << ok;
	// true  (no hay ciclo negativo)
	
	cout << fw.dist[0][2];
	// 4   (via 0 -(3)-> 1 -(1)-> 2, mejor que la ruta directa de 10)
	
	cout << fw.dist[2][1];
	// -1  (via 2 -(-2)-> 0 -(3)-> 1, costo -2+3=1... revisar signo:
	//      en este grafo especifico da 1, no -1 -- el punto es que
	//      dist[i][j] SIEMPRE queda como el minimo real considerando
	//      cualquier nodo intermedio, incluso via el ciclo)
\end{lstlisting}

\textbf{Nota sobre el orden de los \texttt{for}:} el orden
\textbf{$k$ primero, luego $i$, luego $j$} no es arbitrario —es lo
que garantiza la correctitud de la DP. $k$ representa "qué nodos
intermedios ya se me permite usar", así que debe ser el bucle más
externo: al procesar $k$, todas las combinaciones $\texttt{dist}[i][j]$
ya reflejan poder pasar por cualquier nodo $< k$, y solo entonces se
evalúa si pasar también por $k$ mejora algo. Cambiar el orden de los
\texttt{for} (aunque compile y "parezca" correcto) rompe esa garantía
y puede dar distancias mayores a las reales.

\textbf{Regla práctica:} Floyd-Warshall gana cuando necesitas
\textbf{todas} las distancias par a par y $V$ es chico; cuando solo
necesitas las distancias desde \textbf{un} origen (o pocos), correr
Dijkstra (o Bellman-Ford con pesos negativos) desde cada uno de esos
orígenes específicos casi siempre es más rápido que Floyd-Warshall
completo, salvo que ya de por sí necesites la matriz completa para
otra cosa.


% =========================================================
% MST
% =========================================================
\section{Minimum Spanning Tree (MST)}
\subsection{Disjoint Set Union (DSU)}

\textbf{Idea:} DSU (también llamado Union-Find) mantiene una
partición de $n$ elementos en conjuntos disjuntos, soportando dos
operaciones: \texttt{find(x)} (¿a qué conjunto pertenece $x$? —
devuelve un "representante" del conjunto) y \texttt{union(x,y)}
(fusiona los conjuntos de $x$ y $y$). Cada elemento apunta a un
\texttt{padre}; el representante de un conjunto es la raíz de ese
árbol (el nodo que es padre de sí mismo). Sin optimizaciones, esos
árboles pueden degenerar en una cadena larga; dos trucos lo evitan:
\textbf{unión por rango/tamaño} (al fusionar, el árbol más chico
cuelga del más grande, no al revés) y \textbf{compresión de camino}
(en cada \texttt{find}, aplana el camino recorrido para que futuras
consultas sean casi $O(1)$).

\textbf{¿Por qué usarlo?} Es la estructura correcta cuando las
aristas/relaciones \textbf{van llegando una a una} y solo necesitas
responder "¿misma componente?" o "¿ya conectados?" sin reconstruir el
grafo completo cada vez — a diferencia de Componentes Conexas (DFS),
que necesita la lista de adyacencia completa de antemano. Con las dos
optimizaciones juntas, cada operación es prácticamente $O(1)$
amortizado (formalmente $O(\alpha(n))$, la función inversa de Ackermann,
que para cualquier $n$ realista es $\le 4$).

\textbf{Casos de uso:}

\begin{itemize}
	\item Kruskal (siguiente sección): decidir si agregar una arista
	crearía un ciclo, sin reconstruir el grafo.
	\item Detección de ciclos incremental: procesar aristas una por
	una y detectar el momento exacto en que se forma el primer ciclo.
	\item Conectividad dinámica \textit{offline}: responder consultas
	de "¿$u$ y $v$ conectados?" intercaladas con la llegada de nuevas
	aristas (unión), sin volver a correr DFS cada vez.
	\item Agrupar elementos bajo una relación de equivalencia dada
	incrementalmente (ej. "estas cuentas son la misma persona").
\end{itemize}

\textbf{Complejidad:} $O(\alpha(n))$ amortizado por operación
(\texttt{find} y \texttt{union}) con ambas optimizaciones —
prácticamente constante en la práctica.

\begin{lstlisting}[style=cpp]
	struct DSU {
		vector<int> padre, rango;
		
		// CREATE: n elementos, cada uno en su propio conjunto.
		DSU(int n) {
			padre.resize(n);
			rango.assign(n, 0);
			iota(padre.begin(), padre.end(), 0);   // padre[i] = i
		}                                          // O(n)
		
		// READ: representante del conjunto de x, con compresion de
		// camino (aplana el camino recorrido).
		int find(int x) {
			if (padre[x] != x)
			padre[x] = find(padre[x]);   // compresion de camino
			return padre[x];
		}                                          // O(alpha(n)) amortizado
		
		// WRITE: fusiona los conjuntos de x e y. Devuelve false si ya
		// estaban en el mismo conjunto (util para detectar ciclos).
		bool unir(int x, int y) {
			int rx = find(x), ry = find(y);
			if (rx == ry) return false;   // ya conectados -> formaria ciclo
			
			// union por rango: el arbol mas chico cuelga del mas grande
			if (rango[rx] < rango[ry]) swap(rx, ry);
			padre[ry] = rx;
			if (rango[rx] == rango[ry]) rango[rx]++;
			
			return true;
		}                                          // O(alpha(n)) amortizado
		
		// READ: x e y estan en el mismo conjunto?
		bool mismoConjunto(int x, int y) {
			return find(x) == find(y);
		}                                          // O(alpha(n)) amortizado
	};
\end{lstlisting}

\textbf{Ejemplo de funcionamiento:}

\begin{lstlisting}[style=cpp]
	DSU dsu(6);   // {0} {1} {2} {3} {4} {5}
	
	dsu.unir(0, 1);   // {0,1} {2} {3} {4} {5}
	dsu.unir(2, 3);   // {0,1} {2,3} {4} {5}
	dsu.unir(1, 2);   // {0,1,2,3} {4} {5}
	
	cout << dsu.mismoConjunto(0, 3);
	// true  (0 y 3 quedaron en el mismo conjunto tras las uniones)
	
	cout << dsu.mismoConjunto(0, 4);
	// false
	
	cout << dsu.unir(0, 3);
	// false  (0 y 3 YA estaban conectados -> esta union formaria un
	//         ciclo si (0,3) fuera una arista de grafo)
\end{lstlisting}

\textbf{Nota:} el valor de retorno \texttt{bool} de \texttt{unir} no
es cosmético — es exactamente la señal que usa Kruskal para decidir
si una arista pertenece al MST: si \texttt{unir(u,v)} devuelve
\texttt{false}, esa arista cerraría un ciclo y se descarta.


\subsection{Kruskal}

\textbf{Idea:} Kruskal construye un MST (árbol de expansión mínima)
con una estrategia \textbf{greedy}: ordena \textbf{todas} las aristas
por peso ascendente, y las va agregando una por una \textbf{solo si}
no forman un ciclo con las ya agregadas (es decir, si sus dos
extremos están en componentes distintas). Ese chequeo de "¿mismo
conjunto?" es exactamente DSU — Kruskal es, en esencia, "ordenar
aristas + DSU" pegados.

\textbf{¿Por qué usarlo?} Es el algoritmo de MST más simple de
razonar y programar cuando tienes la lista de aristas completa de
antemano (no una matriz de adyacencia densa): el orden importa más
que la estructura del grafo, así que basta con un \texttt{sort} y un
DSU. La alternativa es Prim (crece un solo árbol desde un nodo,
agregando la arista más barata que lo conecta con algo nuevo) — Prim
suele preferirse en grafos \textbf{densos} representados como matriz
de adyacencia; Kruskal suele preferirse en grafos \textbf{dispersos}
representados como lista de aristas.

\textbf{Casos de uso:}

\begin{itemize}
	\item Conectar $n$ puntos (ciudades, computadoras, etc.) con el
	costo total mínimo, dado el costo de cada conexión posible.
	\item Encontrar el \textbf{cuello de botella} mínimo de un grafo:
	la arista más pesada del MST es el mínimo posible "peso máximo de
	arista" entre cualquier árbol que conecte todos los nodos (útil
	en problemas tipo "minimiza la arista más cara del camino").
	\item Verificar si un grafo es conexo: si al final el MST tiene
	menos de $n-1$ aristas, el grafo original no era conexo.
	\item Como filtro previo: procesar aristas en orden de peso y
	usar DSU para responder "¿en qué momento (con qué peso máximo) se
	conectan $u$ y $v$?" — variante de consultas offline.
\end{itemize}

\textbf{Complejidad:} $O(E \log E)$, dominado por el
\texttt{sort} inicial de las aristas; el procesamiento con DSU es
$O(E \cdot \alpha(V))$, prácticamente lineal.

\begin{lstlisting}[style=cpp]
	struct Kruskal {
		struct Arista { int u, v, peso; };
		vector<Arista> aristas;
		int n;
		
		// CREATE: guarda el grafo como lista plana de aristas.
		Kruskal(vector<Arista>& aristas_, int n_) {
			aristas = aristas_;
			n = n_;
		}                                          // O(1)
		
		// READ: construye el MST. Devuelve {costoTotal, aristasDelMST}.
		// Si el grafo no es conexo, aristasDelMST tiene menos de n-1
		// aristas (MST parcial, uno por componente -- un "MST forest").
		pair<long long, vector<Arista>> construirMST() {
			sort(aristas.begin(), aristas.end(),
			[](Arista& a, Arista& b) { return a.peso < b.peso; });
			
			DSU dsu(n);
			long long costoTotal = 0;
			vector<Arista> mst;
			
			for (auto& e : aristas) {
				if (dsu.unir(e.u, e.v)) {   // solo si NO forma ciclo
					costoTotal += e.peso;
					mst.push_back(e);
					
					if ((int)mst.size() == n - 1) break;   // MST completo,
					// no hace
					// falta seguir
				}
			}
			
			return {costoTotal, mst};
		}                                          // O(E log E)
	};
\end{lstlisting}

\textbf{Ejemplo de funcionamiento:}

\begin{lstlisting}[style=cpp]
	// grafo:  0-(1)-1   0-(4)-2   1-(2)-2   1-(5)-3   2-(3)-3
	
	vector<Kruskal::Arista> aristas = {
		{0, 1, 1}, {0, 2, 4}, {1, 2, 2}, {1, 3, 5}, {2, 3, 3}
	};
	
	Kruskal k(aristas, 4);
	auto [costo, mst] = k.construirMST();
	
	cout << costo;
	// 6   (aristas elegidas: 0-1 (1), 1-2 (2), 2-3 (3) = 1+2+3=6)
	
	cout << mst.size();
	// 3   (n-1 = 3 aristas, el MST conecta los 4 nodos)
\end{lstlisting}

\textbf{Regla práctica:} Kruskal es "ordena + DSU"; si ya tienes DSU
programado de la sección anterior, Kruskal es casi trivial de
agregar encima. Elige Kruskal sobre Prim cuando el grafo viene como
\textbf{lista de aristas dispersa}; considera Prim cuando el grafo es
\textbf{denso} y ya lo tienes como matriz de adyacencia, ya que ahí
Prim con un arreglo simple (sin heap) corre en $O(V^2)$ sin necesidad
de ordenar $E \approx V^2$ aristas de antemano.

\subsection{Prim}

\textbf{Idea:} a diferencia de Kruskal (que ordena \textbf{todas}
las aristas globalmente), Prim construye el MST \textbf{creciendo un
	solo árbol} desde un nodo inicial arbitrario: en cada paso, agrega la
arista de \textbf{menor peso} que conecta el árbol ya construido con
algún nodo \textbf{fuera} de él, hasta cubrir los $n$ nodos. Es
literalmente el mismo esqueleto de Dijkstra (heap con el "mejor
candidato" disponible), cambiando qué se minimiza: Dijkstra minimiza
\textbf{distancia acumulada desde el origen}; Prim minimiza
\textbf{el peso de la próxima arista individual} que agrega al árbol.

\textbf{¿Por qué usarlo?} Si ya sabes programar Dijkstra, Prim es
casi copiar y pegar cambiando una línea (qué valor entra al heap). La
diferencia práctica con Kruskal está en la representación del grafo:
Prim con \texttt{priority\_queue} es $O(E \log V)$ igual que
Dijkstra, pero brilla especialmente en su variante \textbf{sin heap}
($O(V^2)$) cuando el grafo es \textbf{denso} y ya lo tienes como
matriz de adyacencia — ahí evitas construir y ordenar una lista de
$E \approx V^2$ aristas como pide Kruskal.

\textbf{Casos de uso:}

\begin{itemize}
	\item Mismo problema que Kruskal (MST de costo mínimo), cuando el
	grafo viene como matriz de adyacencia densa en vez de lista de
	aristas.
	\item Cuando se necesita el MST \textbf{"creciendo" desde un nodo
		específico} de forma incremental (ej. simulaciones donde el árbol
	se construye en tiempo real).
	\item Igual que Kruskal, para encontrar el cuello de botella
	mínimo (la arista más pesada del MST resultante).
\end{itemize}

\textbf{Complejidad:} $O(E \log V)$ con \texttt{priority\_queue}
(igual que Dijkstra, mismo motivo: cada arista puede provocar como
mucho una inserción al heap), o $O(V^2)$ con un arreglo simple sin
heap — preferible en grafos densos donde $E \approx V^2$ y el heap ya
no aporta.

\begin{lstlisting}[style=cpp]
	struct Prim {
		vector<vector<pair<int,int>>> adj;   // adj[u] = {(v, peso)}
		int n;
		
		// CREATE: guarda el grafo (no dirigido, con pesos).
		Prim(vector<vector<pair<int,int>>>& adj_, int n_) {
			adj = adj_;
			n = n_;
		}                                          // O(1)
		
		// READ: construye el MST empezando desde 'inicio'. Devuelve
		// {costoTotal, aristasDelMST} como pares (u, v). Si el grafo no
		// es conexo, el arbol resultante solo cubre la componente de
		// 'inicio' (MST parcial).
		pair<long long, vector<pair<int,int>>> construirMST(int inicio) {
			vector<bool> enArbol(n, false);
			vector<long long> mejorPeso(n, LLONG_MAX);
			vector<int> mejorOrigen(n, -1);
			
			// min-heap por peso de arista: {peso, nodo}
			priority_queue<pair<long long,int>,
			vector<pair<long long,int>>,
			greater<>> pq;
			
			mejorPeso[inicio] = 0;
			pq.push({0, inicio});
			
			long long costoTotal = 0;
			vector<pair<int,int>> mst;
			
			while (!pq.empty()) {
				auto [peso, u] = pq.top(); pq.pop();
				
				if (enArbol[u]) continue;   // entrada obsoleta
				// (ya se agrego con otro peso)
				
				enArbol[u] = true;
				costoTotal += peso;
				if (mejorOrigen[u] != -1)
				mst.push_back({mejorOrigen[u], u});
				
				for (auto& [v, pesoArista] : adj[u]) {
					if (enArbol[v]) continue;
					if (pesoArista >= mejorPeso[v]) continue;
					
					mejorPeso[v] = pesoArista;
					mejorOrigen[v] = u;
					pq.push({pesoArista, v});
				}
			}
			
			return {costoTotal, mst};
		}                                          // O(E log V)
	};
\end{lstlisting}

\textbf{Ejemplo de funcionamiento:}

\begin{lstlisting}[style=cpp]
	// grafo:  0-(1)-1   0-(4)-2   1-(2)-2   1-(5)-3   2-(3)-3
	
	vector<vector<pair<int,int>>> adj(4);
	auto addEdge = [&](int u, int v, int peso) {
		adj[u].push_back({v, peso});
		adj[v].push_back({u, peso});
	};
	addEdge(0, 1, 1);
	addEdge(0, 2, 4);
	addEdge(1, 2, 2);
	addEdge(1, 3, 5);
	addEdge(2, 3, 3);
	
	Prim p(adj, 4);
	auto [costo, mst] = p.construirMST(0);
	
	cout << costo;
	// 6   (mismo MST que Kruskal: 0-1 (1), 1-2 (2), 2-3 (3) = 6)
	
	cout << mst.size();
	// 3
\end{lstlisting}

\textbf{Nota:} igual que en Dijkstra, el chequeo
\texttt{if (enArbol[u]) continue;} es esencial por el mismo motivo:
como \texttt{priority\_queue} en C++ no soporta \textit{decrease-key},
cada mejora de \texttt{mejorPeso[v]} mete una entrada \textbf{nueva}
al heap en vez de actualizar la vieja; sin el chequeo, se
reprocesarían entradas obsoletas (sigue siendo correcto, solo más
lento).

\textbf{Regla práctica:} Kruskal y Prim siempre dan un MST de
\textbf{el mismo costo total} (puede haber empates en qué aristas
exactas eligen si hay pesos repetidos, pero el costo mínimo es único).
La elección entre ambos es de \textbf{representación de datos}, no de
corrección: lista de aristas / grafo disperso → Kruskal; matriz de
adyacencia / grafo denso → Prim (idealmente la variante $O(V^2)$ sin
heap en ese caso).


% =========================================================
% SCC
% =========================================================
\section{Strongly Connected Components (SCC)}
\subsection{Kosaraju}

\textbf{Idea:} Kosaraju encuentra las mismas componentes fuertemente
conexas que Tarjan, pero con \textbf{dos pasadas de DFS} en vez de
una sola. Primera pasada: DFS normal sobre el grafo original,
guardando cada nodo en una pila al \textbf{salir} de su llamada
recursiva (exactamente el post-orden de DFS + Orden Topológico).
Segunda pasada: DFS sobre el grafo \textbf{transpuesto} (todas las
aristas invertidas), procesando los nodos en el orden en que salen de
la pila (de arriba hacia abajo); cada árbol de DFS que se forma en
esta segunda pasada es exactamente una SCC completa.

\textbf{¿Por qué usarlo?} La intuición clave es: si $u$ y $v$ están
en la misma SCC, hay camino de $u$ a $v$ \textbf{y} de $v$ a $u$ —eso
significa que, invirtiendo todas las aristas, esa propiedad se
\textbf{preserva} (ambos caminos siguen existiendo, solo cambiados de
dirección). Procesar los nodos en el grafo transpuesto empezando por
el que terminó \textbf{más tarde} en la primera pasada garantiza que
cada DFS de la segunda pasada no se "escape" hacia otra SCC —se queda
exactamente dentro de la suya. Es más simple de razonar y de
programar que Tarjan (no hay que llevar \texttt{tin}/\texttt{low}/pila
simultáneamente), a cambio de necesitar \textbf{dos} recorridos y
construir explícitamente el grafo transpuesto.

\textbf{Casos de uso:}

\begin{itemize}
	\item Exactamente los mismos que Tarjan: condensar un grafo
	dirigido con ciclos en un DAG, detectar si el grafo es fuertemente
	conexo, encontrar el tamaño de la SCC más grande.
	\item Cuando ya tienes DFS + Orden Topológico programado y
	prefieres reutilizar ese patrón en vez de aprender
	\texttt{tin}/\texttt{low} de Tarjan.
	\item Cuando el problema ya requiere el grafo transpuesto para
	otra cosa (algunos problemas de alcanzabilidad dual lo piden de
	todas formas).
\end{itemize}

\textbf{Complejidad:} $O(V + E)$ — dos DFS completos (uno por
pasada) más $O(V+E)$ para construir el grafo transpuesto; la
constante es mayor que Tarjan (dos recorridos en vez de uno), pero el
orden asintótico es el mismo.

\begin{lstlisting}[style=cpp]
	struct Kosaraju {
		vector<vector<int>> adj, adjT;   // adjT = grafo transpuesto
		vector<bool> visitado;
		stack<int> orden;                // orden de salida (post-orden)
		vector<int> comp;
		int numSCC;
		
		// CREATE: guarda el grafo original y construye el transpuesto.
		Kosaraju(vector<vector<int>>& adj_, int n) {
			adj = adj_;
			adjT.assign(n, {});
			
			for (int u = 0; u < n; u++)
			for (int v : adj[u])
			adjT[v].push_back(u);   // invierte cada arista
			
			visitado.assign(n, false);
			comp.assign(n, -1);
			numSCC = 0;
			
			// pasada 1: DFS en el grafo original, llenando 'orden'
			for (int u = 0; u < n; u++)
			if (!visitado[u])
			dfs1(u);
			
			// pasada 2: DFS en el transpuesto, en orden de salida
			// (de arriba hacia abajo de la pila = de mayor a menor
			// tiempo de salida de la pasada 1)
			fill(visitado.begin(), visitado.end(), false);
			while (!orden.empty()) {
				int u = orden.top(); orden.pop();
				if (!visitado[u]) {
					dfs2(u, numSCC);
					numSCC++;
				}
			}
		}                                          // O(V + E)
		
		void dfs1(int u) {
			visitado[u] = true;
			for (int v : adj[u])
			if (!visitado[v])
			dfs1(v);
			orden.push(u);   // post-orden, igual que DFS + Orden Topologico
		}
		
		void dfs2(int u, int id) {
			visitado[u] = true;
			comp[u] = id;
			for (int v : adjT[u])
			if (!visitado[v])
			dfs2(v, id);
		}
	};
\end{lstlisting}

\textbf{Ejemplo de funcionamiento:}

\begin{lstlisting}[style=cpp]
	// grafo dirigido:  0 -> 1 -> 2 -> 0   (ciclo, una SCC)
	//                  2 -> 3
	//                  3 -> 4 -> 3        (ciclo, otra SCC)
	
	vector<vector<int>> adj(5);
	adj[0] = {1};
	adj[1] = {2};
	adj[2] = {0, 3};
	adj[3] = {4};
	adj[4] = {3};
	
	Kosaraju ks(adj, 5);
	
	cout << ks.numSCC;
	// 2
	
	cout << ks.comp[0] << " " << ks.comp[1] << " " << ks.comp[2];
	// mismo id para 0, 1, 2  (SCC del ciclo 0-1-2)
	
	cout << ks.comp[3] << " " << ks.comp[4];
	// mismo id para 3, 4     (SCC del ciclo 3-4)
\end{lstlisting}

\textbf{Regla práctica:} Kosaraju y Tarjan siempre encuentran
\textbf{exactamente las mismas} SCCs (el algoritmo no cambia el
resultado, solo cómo se llega a él) — la elección entre ambos es de
preferencia de implementación, no de corrección ni de complejidad
asintótica ($O(V+E)$ los dos). Usa Tarjan si prefieres un solo DFS y
ya te sientes cómodo con \texttt{tin}/\texttt{low}/pila (además te da
gratis los ids de SCC en \textbf{orden topológico inverso}, útil para
recorrer el DAG condensado directamente); usa Kosaraju si prefieres
separar mentalmente "encontrar el orden" de "agrupar por
alcanzabilidad" en dos pasos explícitos.

\subsection{Tarjan SCC}

\textbf{Idea:} Tarjan encuentra las componentes fuertemente conexas
en \textbf{una sola pasada} de DFS, usando el mismo mecanismo de
\texttt{tin}/\texttt{low} que Puentes, Puntos de Articulación y
Componentes Biconectadas (ver sección de Conectividad), pero apilando
\textbf{nodos} en vez de aristas: cada nodo se apila al descubrirse,
con un flag \texttt{enPila}; un nodo $u$ es \textbf{raíz de su SCC}
exactamente cuando $\texttt{low}[u] == \texttt{tin}[u]$ —en ese
momento se desapila todo hasta $u$ inclusive, y ese grupo es una SCC
completa. A diferencia de Kosaraju, no hace falta construir el grafo
transpuesto ni un segundo recorrido.

\textbf{¿Por qué usarlo?} Es más eficiente en constante que Kosaraju
(un solo DFS en vez de dos, sin transponer el grafo), y trae una
ventaja extra: las SCCs que va cerrando Tarjan salen naturalmente en
\textbf{orden topológico inverso} del grafo condensado (la primera
SCC en cerrarse es un "sumidero" del DAG condensado, la última en
cerrarse es una "fuente") — útil si después necesitas recorrer el
DAG condensado sin volver a ordenar nada.

\textbf{Casos de uso:} los mismos que Kosaraju — condensar un grafo
dirigido en un DAG, verificar conectividad fuerte total, encontrar la
SCC más grande — con la diferencia de que aquí el orden de cierre de
las SCCs ya viene topológicamente ordenado (inverso), lo cual
aprovecha directamente Condensation Graph más abajo.

\textbf{Complejidad:} $O(V + E)$, un solo DFS con trabajo constante
extra por nodo/arista.

\begin{lstlisting}[style=cpp]
	struct TarjanSCC {
		vector<vector<int>> adj;
		vector<int> tin, low, comp;
		vector<bool> enPila;
		stack<int> pila;
		int timer, numSCC;
		
		// CREATE: guarda el grafo dirigido y corre Tarjan desde cada
		// nodo no visitado. comp[u] queda asignado en orden topologico
		// INVERSO del grafo condensado (comp 0 = la primera en cerrarse
		// = un sumidero del DAG condensado).
		TarjanSCC(vector<vector<int>>& adj_, int n) {
			adj = adj_;
			tin.assign(n, -1);
			low.assign(n, -1);
			comp.assign(n, -1);
			enPila.assign(n, false);
			timer = 0;
			numSCC = 0;
			
			for (int u = 0; u < n; u++)
			if (tin[u] == -1)
			dfs(u);
		}                                          // O(V + E)
		
		void dfs(int u) {
			tin[u] = low[u] = timer++;
			pila.push(u);
			enPila[u] = true;
			
			for (int v : adj[u]) {
				if (tin[v] == -1) {
					dfs(v);
					low[u] = min(low[u], low[v]);
				} else if (enPila[v]) {
					low[u] = min(low[u], tin[v]);
				}
			}
			
			if (low[u] == tin[u]) {
				while (true) {
					int w = pila.top(); pila.pop();
					enPila[w] = false;
					comp[w] = numSCC;
					if (w == u) break;
				}
				numSCC++;
			}
		}
	};
\end{lstlisting}

\textbf{Ejemplo de funcionamiento:}

\begin{lstlisting}[style=cpp]
	// grafo dirigido:  0 -> 1 -> 2 -> 0   (ciclo, una SCC)
	//                  2 -> 3
	//                  3 -> 4 -> 3        (ciclo, otra SCC)
	
	vector<vector<int>> adj(5);
	adj[0] = {1};
	adj[1] = {2};
	adj[2] = {0, 3};
	adj[3] = {4};
	adj[4] = {3};
	
	TarjanSCC tj(adj, 5);
	
	cout << tj.numSCC;
	// 2
	
	cout << tj.comp[3] << " " << tj.comp[0];
	// comp[3] < comp[0]
	// (la SCC {3,4} se cierra primero -- es un sumidero del DAG
	//  condensado, ya que 2 -> 3 pero nada sale de {3,4} hacia {0,1,2})
\end{lstlisting}


\subsection{Condensation Graph}

\textbf{Idea:} el \textit{condensation graph} (grafo condensado) es
el resultado de \textbf{contraer cada SCC a un solo nodo}: si $u$ y
$v$ están en la misma SCC, se fusionan en un nodo; si hay una arista
original de una SCC $A$ a una SCC $B$ distinta, queda una arista
$A \to B$ en el grafo condensado. El resultado \textbf{siempre} es un
DAG —no puede haber ciclos entre SCCs distintas, porque si los
hubiera, esas SCCs en realidad deberían haber sido una sola SCC más
grande (contradicción con que ya son maximales).

\textbf{¿Por qué usarlo?} Es el paso que conecta Tarjan/Kosaraju con
\textbf{todo} lo que ya sabes hacer sobre DAGs: una vez condensado,
puedes aplicar DFS + Orden Topológico, Conteo de Caminos, DP sobre
DAG, o Shortest Path en Grafos Multietapa directamente sobre el
condensado —sin preocuparte de que el grafo original tuviera ciclos.
Es la forma estándar de resolver problemas que dicen "el grafo puede
tener ciclos" pero en realidad piden algo que solo tiene sentido en
un DAG (camino más largo, conteo de caminos, orden de dependencias).

\textbf{Casos de uso:}

\begin{itemize}
	\item Cualquier problema sobre un grafo dirigido con ciclos que
	pida algo intrínsecamente de DAG (camino más largo, número de
	caminos, orden topológico de "grupos").
	\item Sumar/combinar un valor asociado a cada nodo dentro de su
	SCC (ej. "peso total" de cada grupo) antes de propagarlo por el
	DAG condensado.
	\item Determinar cuántas SCCs son alcanzables o pueden alcanzar
	a una SCC dada, tratándolo como alcanzabilidad en un DAG (DFS +
	Cerramiento Transitivo sobre el condensado, mucho más barato que
	sobre el grafo original con ciclos).
\end{itemize}

\textbf{Complejidad:} $O(V + E)$ — una vez calculadas las SCCs
(con Tarjan o Kosaraju, ambos $O(V+E)$), construir el condensado es
un solo recorrido de las aristas originales, agregando una arista
condensada por cada arista original que cruza de una SCC a otra
distinta (con deduplicación opcional si no se quieren aristas
múltiples).

\begin{lstlisting}[style=cpp]
	struct GrafoCondensado {
		vector<vector<int>> adjCond;   // adjCond[SCC] = SCCs vecinas
		int numSCC;
		
		// CREATE: a partir del grafo original y las SCCs ya calculadas
		// (por ejemplo, con TarjanSCC o Kosaraju), construye el DAG
		// condensado. 'comp' asigna a cada nodo original el id de su SCC.
		GrafoCondensado(vector<vector<int>>& adj, vector<int>& comp,
		int numSCC_) {
			numSCC = numSCC_;
			adjCond.assign(numSCC, {});
			
			set<pair<int,int>> vistas;   // evita aristas condensadas
			// duplicadas
			
			for (int u = 0; u < (int)adj.size(); u++) {
				for (int v : adj[u]) {
					if (comp[u] == comp[v]) continue;   // misma SCC,
					// no es arista
					// del condensado
					
					if (vistas.count({comp[u], comp[v]})) continue;
					vistas.insert({comp[u], comp[v]});
					
					adjCond[comp[u]].push_back(comp[v]);
				}
			}
		}                                          // O((V + E) log E)
		// por el set; O(V+E)
		// si se acepta tener
		// aristas duplicadas
	};
\end{lstlisting}

\textbf{Ejemplo de funcionamiento:}

\begin{lstlisting}[style=cpp]
	// grafo dirigido:  0 -> 1 -> 2 -> 0   (SCC A = {0,1,2})
	//                  2 -> 3
	//                  3 -> 4 -> 3        (SCC B = {3,4})
	
	vector<vector<int>> adj(5);
	adj[0] = {1};
	adj[1] = {2};
	adj[2] = {0, 3};
	adj[3] = {4};
	adj[4] = {3};
	
	TarjanSCC tj(adj, 5);
	GrafoCondensado gc(adj, tj.comp, tj.numSCC);
	
	// el condensado tiene exactamente 2 nodos (una por SCC) y 1 arista:
	// SCC{3,4} -> SCC{0,1,2}   (ya que Tarjan cierra {3,4} primero,
	// su id es menor -- ver nota de orden topologico inverso arriba)
	
	cout << gc.adjCond[tj.comp[3]][0] == tj.comp[0];
	// true  (la SCC de 3 tiene una arista hacia la SCC de 0)
\end{lstlisting}

\textbf{Regla práctica:} construir el condensado es siempre el mismo
patrón de dos pasos —(1) calcular SCCs con Tarjan o Kosaraju, (2)
recorrer las aristas originales una vez, quedándote solo con las que
cruzan de SCC —sin importar qué algoritmo de SCC hayas usado. Si
usaste Tarjan, aprovecha que los ids de \texttt{comp} ya vienen en
orden topológico inverso del condensado: puedes iterar
\texttt{numSCC-1} hasta \texttt{0} directamente para procesar el DAG
condensado sin correr otro DFS + Orden Topológico encima.


% =========================================================
% BIPARTITE
% =========================================================
\section{Bipartite Graphs}
\subsection{2-Coloring}

\textbf{Idea:} un grafo es \textbf{bipartito} si sus nodos se pueden
partir en dos grupos tales que \textbf{toda} arista conecta un nodo
de un grupo con uno del otro (nunca dos nodos del mismo grupo). Se
verifica con BFS o DFS "por capas", asignando colores
\textbf{alternados} (0/1) a medida que se visita: el color de un
vecino siempre debe ser el \textbf{opuesto} al del nodo actual. Si en
algún momento un vecino ya visitado tiene el \textbf{mismo} color que
el nodo actual, el grafo \textbf{no} es bipartito —esa arista conecta
dos nodos que deberían estar en grupos distintos y no lo están.

\textbf{¿Por qué usarlo?} Es la misma idea que BFS para Caminos Más
Cortos ("BFS por capas"), reinterpretando la paridad de la capa como
color: nodos en capas pares van a un grupo, nodos en capas impares al
otro. La razón teórica de por qué esto detecta bipartición
correctamente: un grafo es bipartito \textbf{si y solo si} no tiene
ningún \textbf{ciclo de longitud impar} —y un ciclo impar es
exactamente lo que provoca el conflicto de color al cerrar el ciclo.

\textbf{Casos de uso:}

\begin{itemize}
	\item Verificar si un grafo admite partición en dos grupos sin
	conflictos (asignación de tareas mutuamente excluyentes, horarios
	sin choques, etc.).
	\item Como \textbf{prerrequisito obligatorio} de Bipartite
	Matching (siguiente sección): solo tiene sentido buscar un
	matching bipartito si el grafo realmente lo es, y 2-Coloring
	además te da gratis los dos lados $L$/$R$ de la partición.
	\item Modelar restricciones tipo "estos dos elementos no pueden
	estar juntos" como aristas, y usar 2-Coloring para ver si existe
	una asignación válida en dos categorías.
\end{itemize}

\textbf{Complejidad:} $O(V + E)$ — un BFS/DFS normal con una
comparación de color extra por arista, revisando todas las
componentes del grafo (un grafo desconexo puede tener partes
bipartitas y partes no, o varias componentes bipartitas
independientes).

\begin{lstlisting}[style=cpp]
	struct DosColoreo {
		vector<vector<int>> adj;
		vector<int> color;   // -1 = sin colorear, 0 o 1 = los dos grupos
		int n;
		
		// CREATE: guarda el grafo.
		DosColoreo(vector<vector<int>>& adj_, int n_) {
			adj = adj_;
			n = n_;
			color.assign(n, -1);
		}                                          // O(1)
		
		// READ: true si el grafo completo es bipartito. Colorea todas
		// las componentes en el proceso (revisa cada nodo no coloreado).
		bool esBipartito() {
			for (int u = 0; u < n; u++)
			if (color[u] == -1 && !bfs(u))
			return false;
			return true;
		}                                          // O(V + E)
		
		bool bfs(int origen) {
			queue<int> q;
			q.push(origen);
			color[origen] = 0;
			
			while (!q.empty()) {
				int u = q.front(); q.pop();
				
				for (int v : adj[u]) {
					if (color[v] == -1) {
						color[v] = 1 - color[u];   // color opuesto
						q.push(v);
					} else if (color[v] == color[u]) {
						return false;   // mismo color en una arista
						// -> ciclo impar -> no bipartito
					}
				}
			}
			
			return true;
		}
	};
\end{lstlisting}

\textbf{Ejemplo de funcionamiento:}

\begin{lstlisting}[style=cpp]
	// grafo bipartito:  0 - 1     0 - 3
	//                    2 - 1     2 - 3
	// (es un ciclo de longitud 4, par -> bipartito)
	
	vector<vector<int>> adj(4);
	adj[0] = {1, 3}; adj[1] = {0, 2};
	adj[2] = {1, 3}; adj[3] = {0, 2};
	
	DosColoreo dc(adj, 4);
	cout << dc.esBipartito();
	// true
	
	// grafo NO bipartito: triangulo 0-1-2-0 (ciclo de longitud 3, impar)
	vector<vector<int>> adj2(3);
	adj2[0] = {1, 2}; adj2[1] = {0, 2}; adj2[2] = {0, 1};
	
	DosColoreo dc2(adj2, 3);
	cout << dc2.esBipartito();
	// false
\end{lstlisting}


\subsection{Bipartite Matching}

\textbf{Idea:} dado un grafo bipartito con lados $L$ y $R$, un
\textbf{matching} es un subconjunto de aristas donde ningún nodo se
repite (cada nodo de $L$ emparejado con \textbf{a lo más un} nodo de
$R$, y viceversa). El \textbf{matching máximo} se encuentra con el
algoritmo de Kuhn: para cada nodo de $L$ sin emparejar, intenta un
DFS buscando un \textbf{camino aumentante} —si llega a un nodo de $R$
libre, lo empareja directamente; si llega a uno ya emparejado,
intenta "desplazar" a su pareja actual (recursivamente, buscándole
otra opción) para liberar espacio.

\textbf{¿Por qué usarlo?} Es la herramienta estándar para cualquier
problema de "asignación óptima uno-a-uno" entre dos conjuntos
—trabajadores a tareas, estudiantes a proyectos, jugadores a
posiciones— donde cada asignación válida es una arista en un grafo
bipartito. La garantía clave del algoritmo (teorema de Berge): un
matching es máximo \textbf{si y solo si} no existe ningún camino
aumentante —por eso repetir "buscar camino aumentante desde cada nodo
libre de $L$" hasta que ya ninguno encuentre uno, termina
garantizadamente en el matching máximo.

\textbf{Casos de uso:}

\begin{itemize}
	\item Asignación máxima uno-a-uno entre dos conjuntos con
	compatibilidades dadas (trabajadores-tareas, alumnos-cursos).
	\item Como subrutina de problemas más generales: cobertura mínima
	de vértices en grafos bipartitos (teorema de König: tamaño del
	matching máximo = tamaño de la cobertura mínima de vértices).
	\item Problemas de "cuántos pares disjuntos puedo formar" cuando
	la compatibilidad se puede modelar como bipartita (ej. tareas con
	horarios que no chocan, visto como grafo).
\end{itemize}

\textbf{Complejidad:} $O(V \cdot E)$ con Kuhn simple (un DFS por
cada nodo de $L$, en el peor caso $O(E)$ cada uno) — suficiente para
$V \lesssim 1000$–$2000$. Para grafos más grandes existe
Hopcroft-Karp, $O(E\sqrt{V})$, fuera del alcance de esta sección.

\begin{lstlisting}[style=cpp]
	struct BipartiteMatching {
		vector<vector<int>> adj;   // adj[u] = vecinos en R, para cada
		// u en L
		vector<int> matchR;        // matchR[v] = nodo de L emparejado
		// con v en R, -1 si libre
		vector<bool> visitado;     // marca de la busqueda actual (se
		// reinicia en cada intento)
		int nL;
		
		// CREATE: guarda el grafo bipartito (nL nodos en L, adj[u]
		// lista los vecinos de u en R). nR = tamano del lado R.
		BipartiteMatching(vector<vector<int>>& adj_, int nL_, int nR) {
			adj = adj_;
			nL = nL_;
			matchR.assign(nR, -1);
		}                                          // O(1)
		
		// busca un camino aumentante desde u (nodo de L), intentando
		// desplazar parejas ya asignadas si hace falta.
		bool intentar(int u) {
			for (int v : adj[u]) {
				if (visitado[v]) continue;
				visitado[v] = true;
				
				// v esta libre, o su pareja actual puede reubicarse
				if (matchR[v] == -1 || intentar(matchR[v])) {
					matchR[v] = u;
					return true;
				}
			}
			return false;
		}
		
		// READ: calcula el matching maximo. Devuelve su tamano;
		// matchR queda con la asignacion final.
		int matchingMaximo() {
			int total = 0;
			
			for (int u = 0; u < nL; u++) {
				visitado.assign(matchR.size(), false);
				if (intentar(u)) total++;
			}
			
			return total;
		}                                          // O(V * E)
	};
\end{lstlisting}

\textbf{Ejemplo de funcionamiento:}

\begin{lstlisting}[style=cpp]
	// L = {0, 1, 2}, R = {0, 1, 2}
	// 0 - {0, 1}
	// 1 - {0}
	// 2 - {1, 2}
	
	vector<vector<int>> adj(3);
	adj[0] = {0, 1};
	adj[1] = {0};
	adj[2] = {1, 2};
	
	BipartiteMatching bm(adj, 3, 3);
	int maximo = bm.matchingMaximo();
	
	cout << maximo;
	// 3
	// posible asignacion: 1-0 (desplaza a 0), 0-1, 2-2
	// (0 originalmente probo emparejarse con 0 en R, pero al llegar 1
	//  que SOLO puede ir a 0, 0 se desplaza hacia 1 en R, liberando
	//  espacio para que 1 tome 0 en R)
\end{lstlisting}

\textbf{Nota:} el vector \texttt{visitado} se reinicia en \textbf{cada}
llamada a \texttt{matchingMaximo()} por cada nodo de $L$ —no es el
mismo \texttt{visitado} de un DFS normal que se llena una sola vez
para todo el grafo. Aquí marca "ya considerado \textbf{en este
	intento}", para no reprocesar el mismo nodo de $R$ dos veces dentro
de una sola búsqueda de camino aumentante, aunque sí se vuelve a
visitar en el siguiente intento (desde el próximo nodo de $L$).

\textbf{Regla práctica:} 2-Coloring y Bipartite Matching son
secuenciales: primero verificas (y separas en $L$/$R$) con
2-Coloring, y solo si el grafo realmente es bipartito tiene sentido
correr Matching sobre él —Kuhn asume esa estructura y no detecta por
sí solo si el grafo no lo es.


% =========================================================
% TREES
% =========================================================
\section{Trees}
\subsection{DFS en Árboles}

\textbf{Idea:} un árbol es simplemente un grafo conexo sin ciclos
—eso significa que un DFS sobre un árbol \textbf{nunca} necesita
marcar "visitado" con un arreglo separado: basta con no regresar por
la misma arista al \textbf{padre} inmediato, ya que no hay ningún
otro camino de vuelta (no hay ciclos que puedan reconectar el DFS con
un nodo ya procesado). Esa simplificación es lo que distingue DFS en
árboles del DFS de grafos generales, y es la base de casi todo el DP
sobre árboles: recorres en post-orden, y cada nodo combina los
resultados ya calculados de sus hijos.

\textbf{¿Por qué usarlo?} Es el patrón detrás de prácticamente
\textbf{todo} problema de árboles: tamaño de subárbol, profundidad,
diámetro, suma de valores por subárbol, DP sobre árboles en general.
Reconocer "esto es información que depende de mis hijos" es la señal
para post-orden DFS —calculas la respuesta de cada hijo primero, y el
nodo actual combina esos resultados al \textbf{salir} de la
recursión.

\textbf{Casos de uso:}

\begin{itemize}
	\item Tamaño de cada subárbol: $\texttt{tam}[u] = 1 + \sum
	\texttt{tam}[hijo]$.
	\item Profundidad/altura de cada nodo, o altura total del árbol.
	\item Suma o valor agregado de cada subárbol (útil para
	responder queries tipo "¿cuánto suma el subárbol de $u$?").
	\item Diámetro del árbol (la distancia más larga entre dos
	nodos cualquiera): con dos DFS —desde cualquier nodo encuentra el
	más lejano $a$, desde $a$ encuentra el más lejano $b$; la
	distancia $a$-$b$ es el diámetro (propiedad de árboles).
\end{itemize}

\textbf{Complejidad:} $O(n)$ — cada nodo se visita exactamente una
vez, cada arista se recorre exactamente una vez (a diferencia de un
grafo general, aquí no hace falta ni siquiera el arreglo
\texttt{visitado}, basta con excluir al padre).

\begin{lstlisting}[style=cpp]
	struct DFSArbol {
		vector<vector<int>> adj;
		vector<int> tam, profundidad;
		int n;
		
		// CREATE: guarda el arbol (n nodos, n-1 aristas).
		DFSArbol(vector<vector<int>>& adj_, int n_) {
			adj = adj_;
			n = n_;
			tam.assign(n, 0);
			profundidad.assign(n, 0);
		}                                          // O(1)
		
		// READ: corre DFS desde raiz, llenando tam[] y profundidad[]
		// de una sola pasada.
		void dfs(int u, int padre, int d) {
			profundidad[u] = d;
			tam[u] = 1;
			
			for (int v : adj[u]) {
				if (v == padre) continue;   // unica exclusion necesaria,
				// no hace falta 'visitado'
				
				dfs(v, u, d + 1);
				tam[u] += tam[v];   // combina resultados de los hijos
				// DESPUES de que ya terminaron
				// (post-orden)
			}
		}                                          // O(n)
	};
\end{lstlisting}

\textbf{Ejemplo de funcionamiento:}

\begin{lstlisting}[style=cpp]
	// arbol:      0
	//            / \
	//           1   2
	//          / \
	//         3   4
	
	vector<vector<int>> adj(5);
	adj[0] = {1, 2};
	adj[1] = {0, 3, 4};
	adj[2] = {0};
	adj[3] = {1};
	adj[4] = {1};
	
	DFSArbol da(adj, 5);
	da.dfs(0, -1, 0);
	
	cout << da.tam[0];
	// 5  (todo el arbol)
	
	cout << da.tam[1];
	// 3  (1, 3, 4)
	
	cout << da.profundidad[4];
	// 2
\end{lstlisting}


\subsection{BFS en Árboles}

\textbf{Idea:} el mismo principio de "sin ciclos, no hace falta
\texttt{visitado} extra, solo excluir al padre" aplica igual a BFS
sobre árboles —con la ventaja de que BFS recorre \textbf{por
	niveles}, así que es la herramienta natural cuando lo que se pide
depende explícitamente de la \textbf{profundidad/nivel} de cada nodo
en vez de depender de sus subárboles.

\textbf{¿Por qué usarlo?} DFS y BFS en árboles calculan las mismas
distancias/profundidades ($O(n)$ ambos), pero BFS procesa los nodos
en \textbf{orden creciente de profundidad}, lo cual es justo lo que
necesitas cuando el problema pide algo "nivel por nivel" (ej.
"imprime el árbol por niveles", "suma de valores en cada
profundidad") — con DFS tendrías que agrupar por profundidad
manualmente después; con BFS, ya vienen agrupados por construcción.

\textbf{Casos de uso:}

\begin{itemize}
	\item Recorrido por niveles ("level order traversal"): imprimir o
	procesar el árbol nivel por nivel.
	\item Distancia mínima desde la raíz a cada nodo (aunque en un
	árbol, DFS también la da correctamente — la ventaja de BFS es
	solo el orden de procesamiento por capas).
	\item Encontrar el nodo más profundo, o todos los nodos a una
	profundidad exacta $k$.
	\item Como una de las dos pasadas del cálculo de diámetro de
	árbol (BFS desde cualquier nodo para hallar el extremo más
	lejano, igual que con DFS — aquí ambos sirven).
\end{itemize}

\textbf{Complejidad:} $O(n)$ — cada nodo se encola y desencola una
sola vez.

\begin{lstlisting}[style=cpp]
	struct BFSArbol {
		vector<vector<int>> adj;
		vector<int> profundidad, padre;
		vector<vector<int>> porNivel;   // porNivel[d] = nodos a
		// profundidad d
		int n;
		
		// CREATE: guarda el arbol.
		BFSArbol(vector<vector<int>>& adj_, int n_) {
			adj = adj_;
			n = n_;
		}                                          // O(1)
		
		// READ: corre BFS desde raiz, llenando profundidad[], padre[]
		// y agrupando los nodos por nivel.
		void bfs(int raiz) {
			profundidad.assign(n, -1);
			padre.assign(n, -1);
			porNivel.clear();
			
			queue<int> q;
			q.push(raiz);
			profundidad[raiz] = 0;
			
			while (!q.empty()) {
				int u = q.front(); q.pop();
				
				if ((int)porNivel.size() <= profundidad[u])
				porNivel.push_back({});
				porNivel[profundidad[u]].push_back(u);
				
				for (int v : adj[u]) {
					if (v == padre[u]) continue;   // igual que en DFS,
					// basta excluir al
					// padre
					
					padre[v] = u;
					profundidad[v] = profundidad[u] + 1;
					q.push(v);
				}
			}
		}                                          // O(n)
	};
\end{lstlisting}

\textbf{Ejemplo de funcionamiento:}

\begin{lstlisting}[style=cpp]
	// arbol:      0
	//            / \
	//           1   2
	//          / \
	//         3   4
	
	vector<vector<int>> adj(5);
	adj[0] = {1, 2};
	adj[1] = {0, 3, 4};
	adj[2] = {0};
	adj[3] = {1};
	adj[4] = {1};
	
	BFSArbol ba(adj, 5);
	ba.bfs(0);
	
	cout << ba.porNivel.size();
	// 3   (niveles 0, 1, 2)
	
	for (int u : ba.porNivel[1]) cout << u << " ";
	// 1 2   (nivel 1: hijos directos de la raiz)
	
	for (int u : ba.porNivel[2]) cout << u << " ";
	// 3 4   (nivel 2: nietos de la raiz)
\end{lstlisting}

\textbf{Regla práctica:} en árboles, tanto DFS como BFS corren en
$O(n)$ y ninguno necesita \texttt{visitado} —basta con excluir al
padre en cada paso, porque la ausencia de ciclos lo garantiza. Elige
DFS cuando lo que calculas depende de \textbf{combinar resultados de
	subárboles} (post-orden: tamaños, sumas, diámetro vía DP); elige BFS
cuando lo que calculas depende de \textbf{agrupar por
	nivel/profundidad} explícitamente (recorridos por niveles, "todos
los nodos a distancia $k$ de la raíz").

\subsection{Diámetro de un Árbol}

\textbf{Idea:} el diámetro de un árbol es la distancia (en número de
aristas) entre el par de nodos \textbf{más lejano} entre sí. Se
calcula con una propiedad no obvia pero muy útil de los árboles:
\textbf{desde cualquier nodo $x$, el nodo más lejano $a$ es siempre
	un extremo del diámetro}. Eso da el algoritmo de \textbf{dos
	pasadas}: (1) BFS/DFS desde un nodo cualquiera para encontrar el nodo
más lejano $a$; (2) BFS/DFS desde $a$ para encontrar el nodo más
lejano $b$ —la distancia $a$-$b$ \textbf{es} el diámetro. Nótese que
$x$ inicial puede ser literalmente cualquier nodo, no hace falta
adivinar ni probar todos.

\textbf{¿Por qué usarlo?} La alternativa ingenua —correr BFS desde
\textbf{cada} nodo y quedarte con la distancia máxima global— cuesta
$O(n^2)$. El truco de las dos pasadas baja eso a $O(n)$ aprovechando
la estructura de árbol (sin ciclos): la garantía de que "el más
lejano desde cualquier punto es un extremo del diámetro" \textbf{no}
es cierta en grafos generales, solo en árboles.

\textbf{Casos de uso:}

\begin{itemize}
	\item Distancia máxima entre cualquier par de nodos de un árbol.
	\item Encontrar el par de nodos que realiza esa distancia máxima
	(reconstruyendo el camino de $a$ a $b$ con \texttt{padre[]},
	igual que en BFS en Árboles).
	\item Como subrutina de Centro de un Árbol (siguiente sección):
	el centro es literalmente el punto medio del camino del
	diámetro.
	\item Problemas de "radio mínimo" o "tiempo máximo para que algo
	se propague desde el peor nodo posible" —el diámetro acota estos
	valores.
\end{itemize}

\textbf{Complejidad:} $O(n)$ — dos BFS/DFS completos, cada uno
$O(n)$.

\begin{lstlisting}[style=cpp]
	struct DiametroArbol {
		vector<vector<int>> adj;
		int n;
		
		// CREATE: guarda el arbol.
		DiametroArbol(vector<vector<int>>& adj_, int n_) {
			adj = adj_;
			n = n_;
		}                                          // O(1)
		
		// BFS auxiliar: devuelve {nodo mas lejano, distancia} desde src.
		pair<int,int> bfsMasLejano(int src) {
			vector<int> dist(n, -1);
			queue<int> q;
			q.push(src);
			dist[src] = 0;
			
			int mejorNodo = src, mejorDist = 0;
			
			while (!q.empty()) {
				int u = q.front(); q.pop();
				
				if (dist[u] > mejorDist) {
					mejorDist = dist[u];
					mejorNodo = u;
				}
				
				for (int v : adj[u]) {
					if (dist[v] != -1) continue;
					dist[v] = dist[u] + 1;
					q.push(v);
				}
			}
			
			return {mejorNodo, mejorDist};
		}                                          // O(n)
		
		// READ: diametro del arbol (dos pasadas de BFS).
		int diametro() {
			auto [a, _] = bfsMasLejano(0);       // paso 1: cualquier
			// nodo -> extremo a
			auto [b, dist] = bfsMasLejano(a);    // paso 2: a -> extremo b
			
			return dist;
		}                                          // O(n)
	};
\end{lstlisting}

\textbf{Ejemplo de funcionamiento:}

\begin{lstlisting}[style=cpp]
	// arbol:      0
	//            / \
	//           1   2
	//          /
	//         3
	//        /
	//       4
	
	vector<vector<int>> adj(5);
	adj[0] = {1, 2};
	adj[1] = {0, 3};
	adj[2] = {0};
	adj[3] = {1, 4};
	adj[4] = {3};
	
	DiametroArbol da(adj, 5);
	
	cout << da.diametro();
	// 4   (camino 4-3-1-0-2, 4 aristas)
\end{lstlisting}

\textbf{Nota:} existe una alternativa de \textbf{una sola pasada} con
DP (post-orden, como en DFS en Árboles): para cada nodo, calcular las
dos ramas más largas de sus hijos y sumar; el diámetro es el máximo
de esa suma sobre todos los nodos. Es igual de $O(n)$ pero evita el
segundo BFS —útil si ya estás haciendo otro DP sobre el árbol y
quieres aprovechar el mismo recorrido. El método de dos pasadas es
preferible cuando además necesitas identificar \textbf{qué} par de
nodos $(a,b)$ realiza el diámetro, no solo su longitud.


\subsection{Centro de un Árbol}

\textbf{Idea:} el centro de un árbol es el nodo (o los dos nodos)
que \textbf{minimiza la distancia máxima} a cualquier otro nodo del
árbol —el punto "más equilibrado". Un árbol tiene \textbf{uno o dos}
centros nunca más: si el diámetro tiene longitud \textbf{par}, hay un
único centro (el punto medio exacto del camino del diámetro); si es
\textbf{impar}, hay dos centros adyacentes (los dos nodos de en medio
del camino, que no tiene un punto medio exacto).

\textbf{¿Por qué usarlo?} Hay dos formas equivalentes de encontrarlo,
ambas $O(n)$: (1) \textbf{vía el diámetro} —calculas el camino del
diámetro con Diámetro de un Árbol, y el/los nodo(s) en la posición
central de ese camino son el centro; (2) \textbf{"pelado de hojas"}
(\textit{leaf trimming})—quitas todas las hojas actuales del árbol
capa por capa (como un BFS multi-fuente desde las hojas hacia
adentro), y lo que sobrevive al final (1 o 2 nodos) es el centro. La
segunda es más directa de programar cuando ya tienes BFS
multi-fuente a mano (mismo patrón que Shortest Path desde Múltiples
Orígenes, pero "encogiendo" el árbol en vez de calculando distancias).

\textbf{Casos de uso:}

\begin{itemize}
	\item Elegir la raíz "óptima" de un árbol quiere minimizar la
	altura resultante (rerootear en el centro da la altura mínima
	posible).
	\item Problemas de "ubica un servidor/recurso central" que
	minimice la distancia máxima a cualquier nodo cliente.
	\item Como building block de \textbf{centroid decomposition}
	(fuera del alcance de esta sección, pero el centro es el primer
	paso de esa técnica).
\end{itemize}

\textbf{Complejidad:} $O(n)$ con cualquiera de los dos métodos.

\begin{lstlisting}[style=cpp]
	struct CentroArbol {
		vector<vector<int>> adj;
		vector<int> grado;
		int n;
		
		// CREATE: guarda el arbol y el grado inicial de cada nodo
		// (numero de vecinos).
		CentroArbol(vector<vector<int>>& adj_, int n_) {
			adj = adj_;
			n = n_;
			grado.assign(n, 0);
			
			for (int u = 0; u < n; u++)
			grado[u] = adj[u].size();
		}                                          // O(n)
		
		// READ: encuentra el/los centro(s) "pelando" hojas capa por
		// capa hasta que quedan 1 o 2 nodos.
		vector<int> centro() {
			if (n == 1) return {0};
			
			vector<int> g = grado;   // copia para no mutar el original
			queue<int> hojas;
			
			for (int u = 0; u < n; u++)
			if (g[u] == 1)
			hojas.push(u);
			
			int restantes = n;
			vector<int> capaActual;
			
			while (restantes > 2) {
				int cantidadHojas = hojas.size();
				restantes -= cantidadHojas;
				
				for (int i = 0; i < cantidadHojas; i++) {
					int u = hojas.front(); hojas.pop();
					
					for (int v : adj[u]) {
						g[v]--;
						if (g[v] == 1)
						hojas.push(v);
					}
				}
			}
			
			// lo que queda en la cola (1 o 2 nodos) es el centro
			vector<int> resultado;
			while (!hojas.empty()) {
				resultado.push_back(hojas.front());
				hojas.pop();
			}
			return resultado;
		}                                          // O(n)
	};
\end{lstlisting}

\textbf{Ejemplo de funcionamiento:}

\begin{lstlisting}[style=cpp]
	// arbol (cadena):  0 - 1 - 2 - 3 - 4
	// diametro = 4 (par) -> un solo centro, el punto medio exacto
	
	vector<vector<int>> adj(5);
	adj[0] = {1}; adj[1] = {0, 2}; adj[2] = {1, 3};
	adj[3] = {2, 4}; adj[4] = {3};
	
	CentroArbol ca(adj, 5);
	for (int u : ca.centro()) cout << u << " ";
	// 2   (unico centro, justo en medio de la cadena de 5 nodos)
	
	// arbol (cadena):  0 - 1 - 2 - 3
	// diametro = 3 (impar) -> dos centros adyacentes
	
	vector<vector<int>> adj2(4);
	adj2[0] = {1}; adj2[1] = {0, 2};
	adj2[2] = {1, 3}; adj2[3] = {2};
	
	CentroArbol ca2(adj2, 4);
	for (int u : ca2.centro()) cout << u << " ";
	// 1 2   (los dos nodos centrales)
\end{lstlisting}

\textbf{Regla práctica:} el pelado de hojas y el método vía diámetro
\textbf{siempre} coinciden en el resultado —son dos formas de llegar
al mismo centro. Usa pelado de hojas si el problema ya viene en
términos de "ir quitando hojas" o si prefieres no calcular
explícitamente el camino del diámetro; usa el método vía diámetro si
de todas formas ya necesitas el diámetro para otra parte del
problema, ya que entonces el centro sale casi gratis indexando al
medio del camino ya reconstruido.

\subsection{Combinatoria en Árboles Ordenados}

\textbf{Idea:} dado un árbol enraizado de $n$ nodos, un
\textbf{orden topológico} es cualquier permutación de sus nodos donde
cada nodo aparece \textbf{antes} que todos sus descendientes (por
ejemplo, "en qué orden pudieron haberse insertado" si el árbol se
construyó agregando nodos de a uno, siempre como hijo de algo ya
presente). El número total de órdenes topológicos válidos tiene una
fórmula cerrada muy elegante:
$$\frac{n!}{\prod_{v} \texttt{tam}[v]}$$
donde $\texttt{tam}[v]$ es el tamaño del subárbol de $v$ (incluyéndose
a sí mismo). Se calcula con un solo DFS post-orden (igual que DFS en
Árboles) para obtener los tamaños, y luego un producto/cociente
modular sobre todos los nodos.

\textbf{¿Por qué usarlo?} La fórmula evita tener que \textbf{contar}
las permutaciones una por una (imposible, son hasta $n!$) o hacer DP
combinatoria compleja: reduce el problema entero a "un DFS de
tamaños + un producto modular", el mismo patrón de precómputo que ya
usas para factoriales e inversos modulares. Es la herramienta directa
cuando un problema pregunta "¿de cuántas formas se pudo haber
construido/insertado/etiquetado este árbol respetando el orden
padre-antes-que-hijo?".

\textbf{Casos de uso:}

\begin{itemize}
	\item Contar de cuántas formas se pudo insertar $n$ elementos en
	una estructura tipo árbol (heap, trie, etc.) para terminar con la
	forma dada.
	\item Contar etiquetados válidos $1..n$ de los nodos de un árbol
	tales que cada nodo tenga una etiqueta menor que todos sus hijos
	("árbol como heap mínimo").
	\item Como building block de problemas más grandes que combinan
	esta cuenta con otras restricciones (ej. multiplicar por
	factoriales adicionales si además el orden \textbf{entre}
	hermanos importa de alguna forma extra).
\end{itemize}

\textbf{Complejidad:} $O(n)$ para el DFS de tamaños, más
$O(n \log MOD)$ si se calculan inversos modulares individuales con
exponenciación rápida (o $O(n)$ total si se precomputan inversos de
factoriales una sola vez).

\begin{lstlisting}[style=cpp]
	const long long MOD = 1e9 + 7;
	
	struct CombinatoriaArbol {
		vector<vector<int>> adj;
		vector<long long> tam;
		int n;
		
		// CREATE: guarda el arbol enraizado.
		CombinatoriaArbol(vector<vector<int>>& adj_, int n_) {
			adj = adj_;
			n = n_;
			tam.assign(n, 0);
		}                                          // O(1)
		
		long long potMod(long long base, long long exp, long long mod) {
			long long res = 1;
			base %= mod;
			while (exp > 0) {
				if (exp & 1) res = res * base % mod;
				base = base * base % mod;
				exp >>= 1;
			}
			return res;
		}
		
		long long inversoModular(long long a) {
			return potMod(a, MOD - 2, MOD);   // Fermat, MOD debe ser primo
		}
		
		void dfsTam(int u, int padre) {
			tam[u] = 1;
			for (int v : adj[u]) {
				if (v == padre) continue;
				dfsTam(v, u);
				tam[u] += tam[v];
			}
		}                                          // O(n)
		
		// READ: numero de ordenes topologicos validos del arbol
		// enraizado en 'raiz', modulo MOD.
		long long contarOrdenes(int raiz) {
			dfsTam(raiz, -1);
			
			long long factorialN = 1;
			for (int i = 2; i <= n; i++)
			factorialN = factorialN * i % MOD;
			
			long long denominador = 1;
			for (int v = 0; v < n; v++)
			denominador = denominador * tam[v] % MOD;
			
			return factorialN * inversoModular(denominador) % MOD;
		}                                          // O(n log MOD)
	};
\end{lstlisting}

\textbf{Ejemplo de funcionamiento:}

\begin{lstlisting}[style=cpp]
	// arbol:      0
	//            / \
	//           1   2
	//          /
	//         3
	
	vector<vector<int>> adj(4);
	adj[0] = {1, 2};
	adj[1] = {0, 3};
	adj[2] = {0};
	adj[3] = {1};
	
	CombinatoriaArbol ca(adj, 4);
	
	cout << ca.contarOrdenes(0);
	// 3
	// las 3 inserciones validas (etiquetando por orden de insercion):
	// 0,1,2,3 -- 0,1,3,2 -- 0,2,1,3
	// (0 siempre primero; 2 puede ir en cualquier posicion relativa a
	//  la rama 1-3, pero 3 siempre despues de 1)
\end{lstlisting}

\textbf{Nota:} esta fórmula asume que el orden \textbf{entre}
hermanos no distingue resultados distintos (a diferencia de contar
árboles \textbf{ordenados} donde la secuencia exacta de hijos
importa). Si el problema sí distingue el orden de los hijos entre sí
como parte de la respuesta, hay que multiplicar además por
$\prod_v (\text{cantidad de hijos de } v)!$ — ese factor extra cubre
las formas de \textbf{ordenar} los hijos de cada nodo entre sí,
encima de las formas de intercalarlos que ya cuenta la fórmula base.


\subsection{Tree DP}

\textbf{Idea:} Tree DP es el patrón general de DP sobre árboles con
DFS post-orden: cada nodo $u$ tiene uno o más "estados" (ej.
\texttt{dp[u][0]} = mejor respuesta \textbf{sin} usar $u$,
\texttt{dp[u][1]} = mejor respuesta \textbf{usando} $u$), y el valor
de cada estado se calcula \textbf{combinando} los estados ya
resueltos de los hijos —exactamente el mismo esqueleto de DFS en
Árboles (post-orden, sin \texttt{visitado}, solo excluir al padre),
pero acumulando una decisión óptima en vez de solo un tamaño o una
profundidad.

\textbf{¿Por qué usarlo?} Es la generalización de todo lo visto en
la sección de Trees: tamaño de subárbol, diámetro (una pasada), y
Combinatoria en Árboles Ordenados son casos particulares de "Tree
DP" donde el estado es simple (un solo número). El patrón completo
—múltiples estados por nodo, combinados con $\max$/$\min$/suma según
el problema— es lo que resuelve la mayoría de problemas de
optimización sobre árboles (no solo conteo o medición).

\textbf{Casos de uso:}

\begin{itemize}
	\item \textbf{Conjunto independiente máximo} (máxima suma de
	pesos de nodos, sin elegir dos adyacentes): dos estados por nodo,
	clásico ejemplo de esta sección.
	\item \textbf{Vertex cover mínimo} en árboles: variante muy
	similar, cambiando la condición de combinación entre estados.
	\item \textbf{Coloreo/asignación óptima} de nodos con
	restricciones locales entre padre e hijo (cada estado extra
	representa una opción de "color" o "categoría" del nodo).
	\item Cualquier problema donde la respuesta óptima de un subárbol
	dependa de una decisión binaria o categórica tomada en su raíz.
\end{itemize}

\textbf{Complejidad:} $O(n \cdot k)$, donde $k$ es el número de
estados por nodo (para el clásico de 2 estados, $O(n)$ liso) — cada
nodo combina los estados ya resueltos de sus hijos con trabajo
constante (o $O(k)$) por hijo.

\begin{lstlisting}[style=cpp]
	struct ConjuntoIndependienteMax {
		vector<vector<int>> adj;
		vector<int> peso;
		vector<long long> dpSin, dpCon;   // dpSin[u] = mejor sin usar u
		// dpCon[u] = mejor usando u
		int n;
		
		// CREATE: guarda el arbol y el peso de cada nodo.
		ConjuntoIndependienteMax(vector<vector<int>>& adj_,
		vector<int>& peso_, int n_) {
			adj = adj_;
			peso = peso_;
			n = n_;
			dpSin.assign(n, 0);
			dpCon.assign(n, 0);
		}                                          // O(1)
		
		// READ: corre el DFS post-orden que llena dpSin[]/dpCon[] para
		// todo el arbol enraizado en 'raiz'.
		void dfs(int u, int padre) {
			dpCon[u] = peso[u];   // usar u obliga a NO usar a sus hijos
			dpSin[u] = 0;         // no usar u: cada hijo elige libremente
			
			for (int v : adj[u]) {
				if (v == padre) continue;
				
				dfs(v, u);
				
				dpCon[u] += dpSin[v];              // si uso u, el hijo
				// v NO puede usarse
				dpSin[u] += max(dpSin[v], dpCon[v]); // si no uso u, el
				// hijo v elige lo
				// que le convenga
			}
		}                                          // O(n)
		
		// READ: respuesta final del arbol enraizado en 'raiz'.
		long long resolver(int raiz) {
			dfs(raiz, -1);
			return max(dpSin[raiz], dpCon[raiz]);
		}                                          // O(1) tras dfs()
	};
\end{lstlisting}

\textbf{Ejemplo de funcionamiento:}

\begin{lstlisting}[style=cpp]
	// arbol:        0(peso=3)
	//              /          \
	//        1(peso=5)      2(peso=4)
	//            |
	//        3(peso=6)
	
	vector<vector<int>> adj(4);
	adj[0] = {1, 2};
	adj[1] = {0, 3};
	adj[2] = {0};
	adj[3] = {1};
	
	vector<int> peso = {3, 5, 4, 6};
	
	ConjuntoIndependienteMax cim(adj, peso, 4);
	
	cout << cim.resolver(0);
	// 10
	// mejor conjunto: {2, 3} (pesos 4+6=10), no adyacentes entre si
\end{lstlisting}

\textbf{Regla práctica:} cualquier problema de árboles que pida
"óptimo global sujeto a una restricción entre padre e hijos" es Tree
DP: define qué información necesitas de cada hijo (usualmente 1-2
estados), escribe cómo se combina en el nodo actual, y el DFS
post-orden ya está —es el mismo recorrido que DFS en Árboles, solo
con más de un valor acumulándose por nodo.

\subsection{Rerooting}


% =========================================================
% LCA
% =========================================================

\section{Lowest Common Ancestor (LCA)}

\subsection{Binary Lifting}

\subsection{Euler Tour + RMQ}

\subsection{Distancia entre Dos Nodos}

\subsection{$k$-ésimo Ancestro}


% =========================================================
% TREE QUERIES
% =========================================================

\section{Tree Queries}

\subsection{Euler Tour}

\subsection{Subtree Queries}

\subsection{Path Queries}

\subsection{Heavy-Light Decomposition}

\subsection{Heavy-Light Decomposition + Fenwick}

\subsection{Heavy-Light Decomposition + Segment Tree}


% =========================================================
% DISJOINT SET UNION
% =========================================================

\section{Disjoint Set Union}

\subsection{Union-Find}

\subsection{Path Compression}

\subsection{Union by Size / Rank}

\subsection{Unión de Conjuntos Disjuntos en un Árbol}

\subsection{DSU Offline}

\subsection{DSU Rollback}


% =========================================================
% GRAPH DECOMPOSITION
% =========================================================

\section{Graph Decomposition}

\subsection{Bridge Tree}

\subsection{Block-Cut Tree}

\subsection{Condensation Graph}


% =========================================================
% EULERIAN
% =========================================================

\section{Eulerian Graphs}

\subsection{Eulerian Path}

\subsection{Eulerian Circuit}

\subsection{Hierholzer}


% =========================================================
% NETWORK FLOW
% =========================================================

\section{Network Flow}

\subsection{Ford-Fulkerson}

\subsection{Edmonds-Karp}

\subsection{Dinic}

\subsection{Min-Cut}

\subsection{Maximum Bipartite Matching}


% =========================================================
% 2-SAT
% =========================================================

\section{2-SAT}

\subsection{Implication Graph}

\subsection{2-SAT con SCC}


% =========================================================
% BACKTRACKING Y PROBLEMAS CLÁSICOS
% =========================================================

\section{Backtracking y Problemas Clásicos}

\subsection{Problema de las $N$ Reinas}

\subsection{Variaciones del Problema de las $N$ Reinas}

\subsection{Tour de Caballo}


% =========================================================
% ADVANCED
% =========================================================

\section{Advanced Graph Algorithms}

\subsection{Centroid Decomposition}

\subsection{DSU on Tree}

\subsection{Hopcroft-Karp}

\subsection{Min-Cost Max-Flow}

\subsection{DSU Rollback}

%---------------------------------------------------------------------------
\part{Ordenamiento y Búsqueda}
%---------------------------------------------------------------------------

\section{Ordenamiento y Búsqueda (Sorting \& Searching)}
Varias búsquedas lineales y sub-lineales. (La búsqueda binaria tiene
su propia sección.)

\subsection{Número faltante (Missing Number)}
\textbf{Idea:} los números $0,1,\dots,n$ deberían sumar
$\frac{n(n+1)}{2}$ (fórmula de Gauss). Si falta uno, la suma real queda
corta \emph{exactamente} en ese valor, así que el faltante es
(suma esperada $-$ suma real). Basta un recorrido para acumular la suma.
Alternativa sin riesgo de overflow: hacer XOR de $0..n$ y de todos los
elementos; lo que quede es el faltante. \textbf{Complejidad:} $O(n)$.
\begin{lstlisting}[style=cpp]
long long esp = (long long)n*(n+1)/2;          // suma 0..n
long long falta = esp - accumulate(a.begin(),a.end(),0LL);
\end{lstlisting}


\subsection{Máximo / Mínimo (Array Max/Min)}
\textbf{Idea:} recorres el arreglo una sola vez llevando el menor y el
mayor vistos hasta ahora; comparás cada elemento contra ambos y
actualizas. No requiere ordenar (ordenar sería $O(n\log n)$, peor). Con
un recorrido consigues los dos a la vez. \textbf{Complejidad:} $O(n)$.
\begin{lstlisting}[style=cpp]
int mn = *min_element(a.begin(), a.end());
int mx = *max_element(a.begin(), a.end());
\end{lstlisting}


\subsection{Par de suma absoluta mínima}
\textbf{Idea:} buscamos el par cuya suma esté \emph{más cerca de cero}.
Tras ordenar, ponemos un puntero en cada extremo. La suma de los
extremos te dice hacia dónde moverte: si es negativa, para acercarla a
$0$ hay que subir el menor (\texttt{lo++}); si es positiva, hay que
bajar el mayor (\texttt{hi--}). Cada movimiento descarta el candidato
menos prometedor, y en el camino guardas el $|suma|$ más pequeño visto.
Los dos punteros recorren el arreglo una vez ($O(n)$); domina el sort.
\textbf{Complejidad:} $O(n\log n)$.
\begin{lstlisting}[style=cpp]
sort(a.begin(), a.end());
int lo=0, hi=a.size()-1; long long mej=LLONG_MAX;
while(lo<hi){ long long s=(long long)a[lo]+a[hi];
  mej=min(mej,llabs(s));
  if(s<0) lo++; else hi--; }
\end{lstlisting}


\subsection{Par con diferencia dada (Difference Pair)}
\textbf{Idea:} ¿existe un par cuya diferencia sea exactamente $d$? Tras
ordenar, dos punteros avanzan hacia adelante: si la diferencia actual
$a_j-a_i$ es menor que $d$, adelantas el puntero lejano ($j$) para
agrandarla; si es mayor, adelantas el cercano ($i$) para achicarla; si
es igual, encontraste el par. Como el arreglo está ordenado, la
diferencia crece con $j$ y decrece con $i$, así que el barrido nunca
retrocede. \textbf{Complejidad:} $O(n\log n)$ (domina el sort).
\begin{lstlisting}[style=cpp]
sort(a.begin(),a.end()); int i=0,j=1,n=a.size();
while(i<n && j<n){
  if(i!=j && a[j]-a[i]==d) return true;
  if(a[j]-a[i]<d) j++; else i++; }
\end{lstlisting}


\subsection{Búsqueda ternaria (Ternary Search)}
\textbf{Idea:} sirve para hallar el pico de una función
\textbf{unimodal} (sube y luego baja, con un solo máximo — no para buscar
un valor en un arreglo cualquiera). Partes el rango en tres con dos
puntos $m_1<m_2$ y comparas: si $f(m_1)<f(m_2)$, el máximo no puede estar
a la izquierda de $m_1$, así que descartas ese tercio; en caso contrario
descartas el tercio derecho. Cada paso elimina $\sim\frac{1}{3}$ del
rango, por eso converge en logarítmico. (Para un mínimo, invierte la
comparación.) \textbf{Complejidad:} $O(\log n)$ — NO $O(n\log n)$.
\begin{lstlisting}[style=cpp]
while(hi-lo>2){ long long m1=lo+(hi-lo)/3, m2=hi-(hi-lo)/3;
  if(f(m1)<f(m2)) lo=m1+1; else hi=m2-1; }
// revisa [lo,hi] y devuelve el x que maximiza f
\end{lstlisting}


\subsection{Búsqueda de Fibonacci (Fibonacci Search)}
\textbf{Idea:} misma estrategia que la binaria (descartar mitades de un
arreglo ordenado), pero en vez de partir por el punto medio elige el
punto de comparación usando números de Fibonacci consecutivos. Mantiene
tres Fibonacci y en cada paso mira la posición \texttt{off+f2}; según el
elemento sea menor o mayor que el objetivo, retrocede uno o dos pasos en
la secuencia. Ventaja: solo usa \emph{sumas y restas}, nunca división,
lo que la hace útil en hardware sin división rápida o para saltar por
posiciones amigables con la caché. \textbf{Complejidad:} $O(\log n)$.
\begin{lstlisting}[style=cpp]
int f2=0,f1=1,fb=f1+f2;
while(fb<n){ f2=f1; f1=fb; fb=f1+f2; }
int off=-1;
while(fb>1){ int i=min(off+f2,n-1);
  if(a[i]<x){ fb=f1; f1=f2; f2=fb-f1; off=i; }
  else if(a[i]>x){ fb=f2; f1=f1-f2; f2=fb-f1; }
  else return i; }
if(f1==1 && off+1<n && a[off+1]==x) return off+1;
\end{lstlisting}

(En Java es idéntico al C\texttt{++}, cambiando \texttt{a.length}.)

\subsection{Búsqueda exponencial (Exponential Search)}
\textbf{Idea:} dos fases sobre un arreglo \textbf{ordenado}. Primero
\emph{acotas} el objetivo doblando el índice ($1,2,4,8,\dots$) hasta que
$a[i]$ lo supera; en ese momento sabes que el objetivo está entre $i/2$
e $i$. Segundo, haces una binaria normal dentro de ese bloque. Como el
bloque encontrado tiene tamaño proporcional a la posición del objetivo,
es ideal cuando el arreglo es enorme o de tamaño desconocido y el
objetivo está cerca del inicio. \textbf{Complejidad:} $O(\log n)$.
\begin{lstlisting}[style=cpp]
if(a[0]==x) return 0;
int i=1; while(i<n && a[i]<=x) i*=2;   // dobla el rango
// luego binaria en [i/2, min(i,n-1)]
\end{lstlisting}

\subsection{Búsqueda por saltos (Jump Search)}
\textbf{Idea:} en vez de revisar elemento por elemento, avanzas en
bloques de tamaño $\sqrt{n}$ mirando solo el \emph{último} de cada
bloque; cuando ese último ya supera al objetivo, sabes que (si existe)
está en el bloque anterior, así que lo barres linealmente. El balance
$\sqrt{n}$ saltos $+$ $\sqrt{n}$ del barrido minimiza el total. Punto
medio entre la lineal ($O(n)$) y la binaria ($O(\log n)$), útil cuando
saltar hacia atrás es costoso. \textbf{Complejidad:} $O(\sqrt{n})$.
\begin{lstlisting}[style=cpp]
int paso=max(1,(int)sqrt((double)n)), prev=0;
while(prev<n && a[min(prev+paso,n)-1]<x) prev+=paso;
for(int i=prev; i<min(prev+paso,n); i++)
  if(a[i]==x) return i;
\end{lstlisting}


\subsection{Heap Sort}
\textbf{Idea:} un \emph{heap} binario se representa en el mismo arreglo:
el hijo izquierdo de $i$ es $2i+1$ y el derecho $2i+2$; en un
\emph{max-heap} todo padre es $\geq$ sus hijos, así que el mayor está en
la raíz. Heap Sort tiene dos fases. \textbf{(1) Construir el heap}:
recorres de abajo hacia arriba aplicando \texttt{bajar} (sift-down) a
cada padre. Aunque parezca $O(n\log n)$, es $O(n)$: casi todos los nodos
están cerca de las hojas y se hunden poco (la suma
$\sum n/2^h\cdot h$ converge a $2n$). \textbf{(2) Ordenar}: repetidamente
intercambias la raíz (el mayor) con el último elemento activo —quedando
ya ordenado al final— y ``hundes'' la nueva raíz para restaurar el heap;
cada una de las $n$ extracciones cuesta $O(\log n)$. In-place, sin
memoria extra, pero \textbf{no estable}. \textbf{Complejidad:}
$O(n\log n)$ en todos los casos.
\begin{lstlisting}[style=cpp]
void bajar(vector<int>&a,int n,int i){    // hunde: O(log n)
  while(true){ int m=i,l=2*i+1,r=2*i+2;
    if(l<n&&a[l]>a[m]) m=l;
    if(r<n&&a[r]>a[m]) m=r;
    if(m==i) break; swap(a[i],a[m]); i=m; } }
void heapSort(vector<int>&a){ int n=a.size();
  for(int i=n/2-1;i>=0;i--) bajar(a,n,i);   // heap: O(n)
  for(int i=n-1;i>0;i--){ swap(a[0],a[i]);  // max al final
    bajar(a,i,0); } }
\end{lstlisting}


\subsection{QuickSelect}
\textbf{Idea:} encuentra el $k$-ésimo menor \emph{sin ordenar todo}.
Como quicksort, particiona alrededor de un pivote dejando a la izquierda
lo menor y a la derecha lo mayor; pero como sabes en qué lado cae la
posición $k$, recurres (o iteras) sobre \emph{un solo} lado y descartas
el otro. Por eso el promedio baja de $O(n\log n)$ a $O(n)$: cada fase
procesa \emph{la mitad} del anterior, y $n+\tfrac{n}{2}+\tfrac{n}{4}+
\dots = 2n$. El peor caso $O(n^2)$ aparece con pivotes pésimos (arreglo
casi ordenado + pivote extremo); \textbf{elegir el pivote al azar} lo
vuelve prácticamente imposible. La versión de abajo es iterativa (sin
pila de recursión). Si necesitas garantía $O(n)$ en el peor caso, existe
\emph{mediana de medianas} para elegir un buen pivote.
\textbf{Complejidad:} promedio $O(n)$, peor $O(n^2)$.
\begin{lstlisting}[style=cpp]
int quickSelect(vector<int> a, int k){       // k in [0,n-1]
  int lo=0, hi=a.size()-1;
  while(lo<hi){
    int p=a[lo+rand()%(hi-lo+1)], i=lo, j=hi; // pivote aleatorio
    while(i<=j){ while(a[i]<p) i++; while(a[j]>p) j--;
      if(i<=j){ swap(a[i],a[j]); i++; j--; } }
    if(k<=j) hi=j;                             // k a la izq
    else if(k>=i) lo=i;                        // k a la der
    else return a[k];
  }
  return a[lo];
}
// C++ trae nth_element(a.begin(), a.begin()+k, a.end()); // O(n)
\end{lstlisting}

(En Java igual, con \texttt{Random}; código completo en el repo.)

\section{Búsqueda Binaria (Binary Search)}
\textbf{Idea central:} mantienes un rango $[lo,hi]$ donde \emph{sabes}
que está la respuesta y en cada paso lo partes por la mitad, descartando
la mitad que no puede contenerla. Como reduces a la mitad cada vez,
llegas en $\log_2 n$ pasos. Funciona si los datos están \textbf{ordenados}
o, más general, si existe un \textbf{predicado monótono} (una condición
que pasa de falsa a verdadera una sola vez): entonces buscas la frontera
del cambio. \textbf{Complejidad:} $O(\log n)$ en discreto,
$O(\text{iter})$ en reales, $O(L\log n)$ con cadenas. Tip: usa
\texttt{lo+(hi-lo)/2} para el \texttt{mid} y evitar overflow.

\subsection{En arreglo ordenado}
\textbf{Cómo:} comparas $x$ con el elemento del medio; si es menor,
descartas la mitad derecha, si es mayor la izquierda. \texttt{lower\_bound}
te da el primer índice con valor $\geq x$ y \texttt{upper\_bound} el
primer $> x$; con eso sabes si $x$ existe y cuántas copias hay.
\begin{lstlisting}[style=cpp]
int i = lower_bound(a.begin(),a.end(),x)-a.begin(); // >= x
int j = upper_bound(a.begin(),a.end(),x)-a.begin(); // >  x
bool hay = binary_search(a.begin(),a.end(),x);
\end{lstlisting}


\subsection{Menor $x$ que cumple}
\textbf{Cómo:} aquí no buscas en un arreglo sino en el \emph{rango de
respuestas posibles}. El predicado es monótono F..F\,T..T (una vez que
algo ``sirve'', todo lo mayor también). En cada paso pruebas el
\texttt{mid}: si cumple, lo guardas como candidato y sigues buscando a la
izquierda (\texttt{hi=mid-1}) por si hay uno menor que también sirva; si
no cumple, subes (\texttt{lo=mid+1}). Al final \texttt{res} es la
frontera izquierda: el \emph{menor} que cumple. Ejemplo: ``mínima
capacidad para repartir en $\leq k$ días''.
\begin{lstlisting}[style=cpp]
long long res=-1;
while(lo<=hi){ long long mid=lo+(hi-lo)/2;
  if(cumple(mid)){res=mid; hi=mid-1;} // sirve -> izq
  else lo=mid+1; }
\end{lstlisting}


\subsection{Mayor $x$ que cumple}
\textbf{Cómo:} el caso espejo. Ahora el predicado es T..T\,F..F (sirve
hasta cierto punto y luego deja de servir). Cuando el \texttt{mid}
cumple, lo guardas y buscas a la \emph{derecha} (\texttt{lo=mid+1}) por
si hay uno mayor que también sirva; si no cumple, bajas
(\texttt{hi=mid-1}). Respecto al anterior solo cambian esas dos líneas.
Ejemplo: ``máxima longitud de corte tal que salgan $\geq k$ piezas''.
\begin{lstlisting}[style=cpp]
long long res=-1;
while(lo<=hi){ long long mid=lo+(hi-lo)/2;
  if(cumple(mid)){res=mid; lo=mid+1;} // sirve -> der
  else hi=mid-1; }
\end{lstlisting}


\subsection{Sobre reales}
\textbf{Cómo:} cuando la respuesta es un número \emph{real} (raíces,
puntos de equilibrio) no hay ``mitad entera'' ni condición $lo\leq hi$
para parar. En su lugar iteras un número fijo de veces (unas 100 basta:
el rango se divide por $2^{100}$, precisión más que suficiente),
moviendo \texttt{lo} o \texttt{hi} al \texttt{mid} según el predicado.
Al terminar, \texttt{lo}$\approx$\texttt{hi} es la frontera buscada.
\begin{lstlisting}[style=cpp]
for(int it=0; it<100; it++){ double mid=(lo+hi)/2;
  if(cumple(mid)) hi=mid; else lo=mid; }
// lo ~ hi = frontera
\end{lstlisting}


\subsection{Con strings}
\textbf{Cómo:} exactamente la misma binaria, pero comparando cadenas en
orden \emph{lexicográfico} (diccionario) en lugar de números. La única
diferencia de costo es que comparar dos cadenas cuesta $O(L)$ (longitud),
no $O(1)$, así que el total sube a $O(L\log n)$. Ojo: el arreglo debe
estar ordenado con el mismo criterio (ASCII: mayúsculas antes que
minúsculas).
\begin{lstlisting}[style=cpp]
// vector<string> s ordenado
int i = lower_bound(s.begin(),s.end(),string("caro"))
        - s.begin();
bool hay = binary_search(s.begin(),s.end(),string("caro"));
\end{lstlisting}


%---------------------------------------------------------------------------
\part{Matemáticas}
%---------------------------------------------------------------------------

\section{Matemáticas (Math)}
\pendiente

\section{Teoría de Números (Number Theory)}
\pendiente



\pendiente

\part{Geometría Computacional}



\part{Programación Dinámica (Dynamic Programming)}
\pendiente

\part{Algoritmos Voraces (Greedy)}
\pendiente

\part{Ad-hoc y Simulación}
\pendiente

%---------------------------------------------------------------------------
\part{String Processing}
%---------------------------------------------------------------------------

\section{Cadenas (Strings)}
\pendiente

\part{BitWise}

\part{Game Theory}
\end{document}
